% concatenated from Version_2.tex
\documentclass[12pt,paper]{article}
\usepackage{amsmath,amssymb,amsthm,amsfonts,fullpage}
\usepackage{enumitem}  % for using the left marging in itemize
\usepackage[
backend=bibtex,
style=numeric,
sorting=nyt
]{biblatex}
\usepackage{authblk}
\renewcommand\Affilfont{\itshape\small}
\addbibresource{Bibleo.bib}
\usepackage[hidelinks,pdfauthor={Abdelghaffar Chibloun}]{hyperref}
\hypersetup{
	colorlinks   = true,
	citecolor    = blue
}
\usepackage{relsize}
\usepackage{cleveref}
\usepackage{comment}
\usepackage{color,xcolor,soul}
\newcommand{\hlc}[2][yellow]{ {\sethlcolor{#1} \hl{#2}} }
\newcommand{\ol}[2][orange]{ {\sethlcolor{#1} \hl{#2}}}
\soulregister\cite7
\soulregister\ref7
%%
% Keywords command
\providecommand{\keywords}[1]
{
  {\small	
  \textbf{{Keywords--}} #1}
}
% MSC command
\providecommand{\msc}[1]
{
  {\small	
  \textbf{{Mathematics Subject Classification (2020) --}} #1}
}

%\def\MakeUppercaseUnsupportedInPdfStrings{\scshape}   %  to deal with conflicts between some packages and hyperref 
\usepackage{comment}
\theoremstyle{plain}
\newtheorem{theorem}{Theorem}[section] % Main counter
    \newtheorem{lemma}[theorem]{Lemma} % Shares numbering with theorem
    \newtheorem{proposition}[theorem]{Proposition}
    \newtheorem{corollary}[theorem]{Corollary}
    \newtheorem{definition}[theorem]{Definition}
    \newtheorem{example}[theorem]{Example}
    \newtheorem{remark}[theorem]{Remark}
    \newtheorem{conjecture}[theorem]{Conjecture}
	\newtheorem*{notation}{Notation}
	\newtheorem*{main}{Cain Theorem}
	\newtheorem*{problem}{Problem}

\def\dsum #1#2{\displaystyle{\sum_{#1}^{#2}}}
\def\dprod #1#2{\displaystyle{\prod_{#1}^{#2}}}
\def\dlim {\displaystyle{\lim_{{n}\to{+\infty}}}}
\def\dcup #1#2{\displaystyle{\bigcup_{#1}^{#2}}}
\def\dcap #1#2{\displaystyle{\bigcap_{#1}^{#2}}}

\newcommand{\lleft}{\overleftarrow}
\newcommand{\vecc}{\overrightarrow}
\newcommand{\eps}[1]{{#1}_{\varepsilon}}
\newcommand{\ndiv}{\hspace{-4pt}\not|\hspace{2pt}}
\newcommand{\Fq}{\mathbb{F}_{q}}
\DeclareMathOperator{\e}{e}
\DeclareMathOperator{\Z}{\mathcal{Z}}
\DeclareMathOperator{\rank}{rank}
\DeclareMathOperator{\Id}{Id}
\DeclareMathOperator{\Tr}{Tr}
\DeclareMathOperator{\ord}{ord}
\DeclareMathOperator{\card}{card}
\DeclareMathOperator{\aut}{Aut}
\DeclareMathOperator{\Per}{Per}
\DeclareMathOperator{\GL}{GL_{n}}
\DeclareMathOperator{\lcm}{lcm}
\DeclareMathOperator{\lclm}{lclm}        
\DeclareMathOperator{\gcrd}{gcrd}
\DeclareMathOperator{\diag}{diag}
\DeclareMathOperator{\id}{Id}
\DeclareMathOperator{\Frob}{Frob}
\DeclareMathOperator{\Span}{Span}
\DeclareMathOperator{\Ker}{Ker}
\DeclareMathOperator{\Com}{Com}
\DeclareMathOperator{\Fix}{Fix}
\renewcommand{\labelitemi}{$\circ$}



%\author{\textsc{Chibloun Abdelghaffar}\\\textsc{\small fs umi}}
\title{\textbf{Bounds and Equivalence of Skew Polycyclic Codes over Finite Fields}}
 \author[1]{Hassan  Ou-azzou\thanks{E-mail: \rm{hassan.ouazzou@student.unisg.ch}. Supported by a Swiss Government Excellence Scholarship (ESKAS), no:  2024.0504.}}
  \author[2]{Anna-Lena Horlemann\thanks{E-mail: \rm{anna-lena.horlemann@unisg.ch}}}
   \author[3]{Nuh Aydin\thanks{E-mail: \rm{aydinn@kenyon.edu}}}
 

 \renewcommand\Affilfont{\itshape\small}
\affil[1,2]{School of Computer Science, University of St.Gallen, 9000 St.\ Gallen, Switzerland}
%\affil[2]{School of Computer Science, University of St.Gallen, St. Gallen, Switzerland}
\affil[3]{Department of Mathematics and Statistics, Kenyon College, Gambier, OH 43022, USA}
 
 
\date{ \today}
\begin{document}



	
	
	\maketitle
    
    \begin{abstract}

    We study skew polycyclic codes over a finite field $\mathbb{F}_q$, associated with a skew polynomial $f(x) \in \mathbb{F}_q[x;\sigma]$, where $\sigma$ is an automorphism of $\mathbb{F}_q$.  We establish Roos-like bounds on the minimum Hamming and rank distances for these codes, generalizing several known distance bounds for various (consta)cyclic and skew (consta)cyclic code families. To construct codes with a prescribed minimum distance, we develop an approach based on $\mu$-closed sets, which serve as a non-commutative analog of cyclotomic cosets in the commutative case. We further study the equivalence of ambient spaces for skew polycyclic codes under both the Hamming and rank metrics. General conditions are derived under which two such ambient spaces yield equivalent code families, allowing a complete classification of equivalence classes. In the special case of skew trinomial codes, we determine when a code is equivalent to one defined by a trinomial of the form $x^n-x^{\ell}-1,$ and we describe a general procedure to determine the equivalence class of any skew polycyclic code. Several examples are provided to illustrate the application of the theoretical results on code equivalence.
    




\end{abstract}
   \keywords{Skew polycyclic codes,   skew  constacyclic codes, equivalent codes, irreducible polynomials, skew polynomials.}


  		
        		        	

\section{Introduction}
Coding theory plays a fundamental role in various applications, such as error detection and correction, data transmission, data storage and reliable communication. It involves the study of efficient encoding and decoding methods for transmitting data reliably over noisy channels. {Cyclic codes} are one of the most important families of linear codes for both theoretical and practical reasons. They establish a key link between coding theory and algebra and their structure often makes them convenient for implementations. Furthermore, many of the most important classes of codes are cyclic or related to cyclic codes. Cyclic codes were introduced in the 1950's by Prange in \cite{Prange1957} as linear codes with the property that the cyclic shift of any codeword is another codeword. They were later generalized to {constacyclic codes} in \cite{Berlekamp1968} and to {polycyclic codes} (also known as {pseudo-cyclic codes}) in \cite{10}.  Like cyclic (and constacyclic) codes, each polycyclic code over a finite field $\Fq$ can be described by an ideal of the polynomial ring $\Fq[x]/\langle f(x)\rangle,$ where $f(x)$ is a nonzero polynomial in $\Fq[x]$.  Polycyclic codes contain constacyclic codes as a special case when $f(x)=x^n-\lambda,$ for some non-zero $\lambda$ in $\Fq.$

One of the generalizations of cyclic codes are skew cyclic codes (\cite{Boucher2009,Boucher2014}). They can be viewed as left $\Fq[x;\sigma]$-submodules of $\Fq[x;\sigma] / \langle x^n-1\rangle,$ where $\sigma$ is an automorphism on $\mathbb{F}_q$, and $\langle x^n-1\rangle$ is the left ideal of $\Fq[x;\sigma]$ generated by $ x^n-1.$ Similarly to classical cyclic codes, skew cyclic codes are generalized in various ways, such as skew constacyclic codes, skew polycyclic  codes, and others. We refer to the following references for more details \cite{Almendras2018, Bag2025, Boulanouar2021, Ouazzou2023}.

In \cite{Chen2014}, an equivalence relation was introduced to study the equivalence of classes of constacyclic codes of length $n$ over $\mathbb{F}_q$, depending on the ideal defining the quotient ring. This was generalized to skew constacyclic codes over $\Fq$ \cite{Ouazzou2025}, and to the case of constacyclic codes over finite chain rings \cite{Chibloun2024}. In \cite{Boulanouar2021,Lobillo2025} some characterizations on equivalence of skew constacyclic codes were given. Regarding polycyclic codes,  several properties of trinomial codes and a number of conjectures related to their equivalence and duality are presented in \cite{Aydin2022, Shi2023}. The study of polycyclic codes was continued by extending the notion of $n$-equivalence to the case of polycyclic codes in \cite{Equiv2025} 
%\textcolor{purple}{[I added our arXiv paper here]}, 
where the authors studied the properties of this equivalence relation, computed the number of $n$-equivalence classes, and provided conditions under which two ambient spaces of polycyclic (or $\ell$-trinomial) codes are equivalent. 
%We presented additional results on the class of trinomial codes associated with the polynomials $ x^p-x-a$ and $x^q-x-a.$ 

In this paper,  we study skew polycyclic codes over a finite field $\mathbb{F}_q$, associated with a skew polynomial $f(x) \in \mathbb{F}_q[x;\sigma]$, where $\sigma$ is an automorphism of $\mathbb{F}_q$. We start  by  proving  the Roos-like bound for both the Hamming and the rank metric for this class of codes. Next, we  focus on the Hamming and rank equivalence between two classes of polycyclic codes by introducing an equivalence relation and describing its equivalence classes. 
    Finally, we present examples that illustrate applications of the theory developed in this paper.


The remainder of this paper is organized as follows. Section 2 provides a review of the basic background on skew polynomial rings and skew polycyclic codes, plus the necessary background on the Hamming and the rank metric.  In Section 3, we prove the Roos-like bound for both metrics for this class of codes and study constructions of such codes with a designed distance. In Section 4 we study the equivalence of ambient spaces for skew polycyclic codes and provide conditions under which two skew polycyclic code families are equivalent -- again for both the Hamming and the rank metric. Finally, in Section 5, we provide examples for code equivalences studied in section 4.






\section{Preliminaries}
In this section we recall the necessary background about error correcting codes in the Hamming and rank metric, skew polynomials, and skew polycyclic codes.

Throughout we will denote a finite field with $q$ elements (for $q$ being a prime power) by $\Fq$. Similarly an extension field of $\Fq$ of degree $m$ is denoted by $\mathbb F_{q^m}$.  The general linear group $\GL(\mathbb{F}_{q}) $ is the group of invertible $n\times n$ matrices with entries from $\mathbb{F}_{q}$.




 
\subsection{Hamming and rank metric codes}

We start with definitions and results on error correcting codes in the Hamming and the rank distance. For more details  we refer the reader to \cite{augot2018generalized,gabidulin2025rank,Gabidulin1985,Huffman2003} 

\begin{definition}
      A $q$-ary \emph{code} of length $n$ is simply a subset of $\Fq^n$. A \emph{linear code} is a vector subspace of $\Fq^n$.
  \end{definition}

  \begin{definition}[Hamming distance]
  Let  $\Fq$  be a finite field.
  \begin{enumerate}
       \item The \emph{Hamming weight} of a vector $x=(x_0,x_1,\ldots, x_{n-1})\in \Fq^n$,  denoted by $  w_{H}(x),$ is the number of  nonzero coordinates of $x.$
       \item The \emph{Hamming distance} $d_H(x,y)$ between two vectors $x,y \in \Fq^n$ is defined as 
       $$ d_H(x,y):=w_H(x-y) .$$
\item  The   \emph{minimum Hamming distance} of a linear code $C\subseteq \Fq^n$, denoted $ d_H(C)$, is defined as the minimum Hamming weight $w_H(c)$ over all nonzero codewords $c \in C$, i.e.,
$$
d_{H}(C):=\min \left\{w_H(x): x \in C, x \neq 0\right\} .
$$
  \end{enumerate}
  \end{definition}


The Hamming distance satisfies the following bound, called the \emph{Singleton bound} (see, e.g., \cite{Huffman2003}): 
\begin{theorem}
Let $C\subseteq \Fq^n$ be a linear code of dimension $k$. Then
    $$ d_H(C)\leq n-k+1.$$
When equality is obtained, we say that $C$ is a  \emph{Maximum Distance Separable (MDS)} code.
\end{theorem}


  While codes in the Hamming metric are defined as sets of vectors, rank metric codes can be defined as subsets of the ring of $m\times n$ matrices over $\mathbb F_q$, $ \mathbb M_{m\times n}(\Fq).$ Such a code is called \emph{linear}\footnote{The notion for linearity differs in the literature, often depending on the largest field over which a code is linear.}, if it is a linear  subspace of $
\mathbb M_{m\times n}(\Fq).$ 
Via a vector space isomorphism $\mathbb M_{m\times n}(\Fq) \cong \mathbb{F}_{q^m}^n $ rank metric codes can also be represented as subsets of $\mathbb{F}_{q^m}^n$. A linear code in the above sense is then an $\Fq$-linear subspace of $\mathbb{F}_{q^m}^n$.

%for more details see \cite{bartz2022rank,augot2013rank, augot2018generalized, Gabidulin1985, gabidulin2009rank,Neri2020,Morrison2014}. 

%In this paper, we are interested in linear rank metric codes seen as  $\mathbb{F}_{q}$-vector subspaces of $\Fq^n,$ via a vector space isomorphism $\mathbb M_{\mu\times n}(\mathbb{F}_{q_0}) \cong \mathbb{F}_{q_0^\mu}^n=\Fq^n ,$ over $\mathbb{F}_{q_0}:=\Fq^{\sigma},$ is the fixed subfield of $\Fq$ by $\sigma,$ and   $ \mu$ the order of $\sigma.$ In the following  we recall the  essential about  rank metric codes that we need throughout this paper.
  

\begin{definition}[Rank distance]\label{rank_weight}
Consider an extension field $\mathbb{F}_{q^m}$.
\begin{enumerate}
\item  The \emph{rank weight} of $ x=(x_0, \ldots, x_{n-1}) \in \mathbb{F}_{q^m}^n$ over $ \mathbb{F}_{q},$ denoted $w_{r}(x)$, is the dimension of the $ \mathbb{F}_{q}$-vector space spanned by $x_0, \ldots, x_{n-1},$ i.e., $$ w_{r}(x):=\dim_{\mathbb{F}_{q}}(\langle x_0,\ldots,x_{n-1}\rangle).$$ 
    \item The relation  $d_R(x, y):=w_r(x-y)$ for $x, y \in \mathbb F_{q^m}^n$ defines a distance on $\mathbb F_{q^m}^n$ called the \emph{rank distance.}
    \item  The \emph{minimum rank distance} $d_R(C)$ of a linear code $C\subseteq \mathbb F_{q^m}^n$ is defined as 
$$ d_R(C):=\min \left\{ w_r(x) \ : \ x \in C, x \neq 0\right\} .$$
\end{enumerate}
\end{definition}


The rank distance $d_R$ of $C$ satisfies the following bound, called the \emph{Singleton-like (or rank Singleton) bound} (see, e.g., \cite{gabidulin2025rank}):
\begin{theorem}
Let $C\subseteq \mathbb{F}_{q^m}^n$ be a linear code. Then
    $$
\dim_{\mathbb{F}_{q}}(C) \leq \max \{n, m\}(\min \{n, m \}-d_R(C)+1) .
$$
 When equality in the Singleton-like bound occurs, we say  that $C$ is a  \emph{Maximum Rank Distance (MRD)} code.
\end{theorem}

 
In the following proposition from \cite[Proposition 7]{Alfarano2021}, we recall a connection between the minimum rank distance and the minimum Hamming distance of a code.
\begin{proposition}\label{Prop_Hamming_rank}
\cite[Proposition 7]{Alfarano2021}
 Let $C$ be a linear code of length $n$ over $ \mathbb{F}_{q^m}.$ Then
$$ d_R(C)=\min \left\{d_H(C \cdot M): M \in \GL(\mathbb{F}_{q})\right\} .$$
% where $\GL(\mathbb{F}_{q}) $ is the group of invertible $n\times n$ matrices with entries in $\mathbb{F}_{q}.$
\end{proposition}



%\textcolor{purple}{I think it is nice to give a reference to Gabidulin codes here, but I don't think we need the definition, do we?}
%\textcolor{blue}{For some values of the parameters  of the Roos-like bound Theorem, we find the Gabidulin codes (rank distance) and Reed Solomon codes (Hamming distance), I will add these cases later.  }
%\textcolor{purple}{Ok, but then also put the definition of RS-codes here. Moreover, I would put the definitions of the linear (or semi-linear) isometries for both distances here (instead of in Section 4).}

\vspace{0.2cm}

The notion of equivalence is usually defined via linear or semi-linear isometries of the ambient metric space. For more information on the Hamming and rank isometries we refer the interested reader to \cite{Berger2003Isometries,Huffman2003,Neri2020}. In our work, we will restrict ourselves to linear isometries and hence get the following two types of equivalence:

\begin{definition} Let $C_1$ and $C_2$ be two linear codes of length $n$ over $\mathbb{F}_{q^m}$. Then 
\begin{enumerate}
    \item   $C_1$ and $C_2$ are called \emph{Hamming equivalent} if there exists a $n \times n$ monomial matrix $P$ with entries in $ \mathbb{F}_{q^m}$ such that $C_1=C_2 P$.
    \item  $C_1, C_2 $ are \emph{rank equivalent}  if there exists an invertible  matrix $P\in \GL(\mathbb{F}_{q}),$ and a scalar $\lambda \in \mathbb{F}_{q^m}^* $ such that $C_1= \lambda C_2 P.$
    %\textcolor{purple}{The two points should be comparable, i.e., we should really say what the isometries are here. With the sentence I have written above, we should write here: $C_1$ and $C_2$ are called \emph{rank equivalent} if there exists a matrix $P\in \GL(\Fq)$ such that $C_1=C_2 P$.}
\end{enumerate}
\end{definition} 

\subsection{Skew polynomial rings}
In this subsection we give necessary background we need in skew polynomials, for more details see \cite{Ore33,Lam1988,Lam2004-I,Leroy2004,Leroy2008}. Let $\Fq$ be the finite field with $q$ elements, where $q=p^{s}$ for some prime number $p$ and $s \in \mathbb{Z}_{\geq 1}$.  Let $\sigma : \Fq \rightarrow \Fq$ be an automorphism of $\Fq$. Recall that the  \emph{Frobenius  automorphism} of  $\Fq$   is the automorphism $ \tau : \Fq \rightarrow \Fq$ defined by $$\tau(a):=a^{p}.$$ 
 Note that $\tau $ generates the cyclic group of automorphisms of $  \Fq,$ denoted by $\aut(\Fq)$, hence each automorphism $\sigma$ is of the form $ \sigma=\tau^k, \ 0\leq k\leq  s-1.$   Since $\tau $ is a generator of the cyclic group $ \aut(\Fq)$ its order is $s:=[\Fq:\mathbb{F}_p],$ and other generators of $\aut(\Fq)$ are of the form $ \sigma = \tau^{k},$ where $ \gcd(s,k)=1.$  
 
 Recall also that an element $a\in \Fq$ is \emph{fixed} by $\sigma$ if $\sigma(a)=a,$ and the set of fixed elements by $\sigma,$ denoted by $\Fq^{\sigma},$ is a subfield of $\Fq.$   Note that if  $\sigma=\tau^r,$  then its \emph{fixed subfield} is  $\Fq^{\sigma}= \mathbb{F}_{p^{\gcd(r,s)}}.$ Hence the fixed subfield of $\Fq$ by any generator of  $\aut(\Fq)$ is $\mathbb{F}_p.$ 

 Throughout  this paper we fix the following notations: we consider an automorphism  $\sigma \in \aut(\Fq),$ such that $$\sigma(a)=a^{p^r},$$ of order $ \mu:=\frac{s}{\gcd(r,s)},$ and its fixed subfield $\Fq^{\sigma}=\mathbb{F}_{q_0}$, for $\ q_0:=p^{\gcd(s,r)}.$   
 

\begin{definition}
The \emph{$\sigma$-skew polynomial ring} $\Fq [x;\sigma]$ is the set 
$$ \Fq [x;\sigma]:=\left\{a_{0}+a_{1} x+\cdots+a_{n-1} x^{n-1}: a_{i} \in \Fq, n \in \mathbb{N}\right\}$$
endowed with  the usual polynomial addition and   multiplication defined by 
$$ xa=  \sigma(a) x ,$$
extended to all polynomials by distributivity. 
%This ring is not commutative unless $\sigma$ is the identity automorphism. 
\end{definition}

Recall that an element $f\in \Fq[x;\sigma]$ is called   \emph{central} if $f(x) g(x) = g(x) f(x)$ for all $f(x) \in \Fq[x;\sigma]$, i.e., $f$ commutes with every element of $\Fq[x;\sigma]$. The set of all elements of $\Fq[x;\sigma]$ that commute with every element of $\Fq[x;\sigma]$ is called the  \emph{center} of $\Fq[x;\sigma]$ which is defined and  denoted by
\begin{center}
 $ \Z(\Fq[x;\sigma]):=\left\lbrace f(x) \in \Fq[x;\sigma] \ : \  f(x) g(x) = g(x) f(x), \ \text{ for all} \ \ g(x) \in \Fq[x;\sigma] \right\rbrace .$
\end{center}
%\textcolor{purple}{We should not use $R$, but rather $\Fq[x;\sigma]$. This appears several times throughout the paper.}


In the following we collect some properties of  the  skew polynomial ring $\Fq[x;\sigma]$, see also \cite[p. 483-486]{Ore33}.

\begin{proposition} \label{PPP.1}
\begin{enumerate}
\item The ring $\Fq[x;\sigma]$  is non-commutative unless $\sigma $ is the identity automorphism of $\Fq.$ 
\item   An element $f \in \Fq[x;\sigma]$ is central if and only if $f$ belongs to $ \Fq^{\sigma}[x^{\mu}]$ i.e., $  \Z(\Fq[x;\sigma])= \Fq^{\sigma}[x^{\mu}],$ where  $ \Fq^{\sigma}$ is the fixed subfield of $\Fq$ by $\sigma$ and $\mu$ is the order  of $\sigma.$
\item The ring $\Fq[x;\sigma]$ is a right (respectively left)   Euclidean domain. By  "$x \mid_r y$"   we denote that $x$ right divides $y$.
\item Left and right ideals of $\Fq[x;\sigma]$ are principal.
\end{enumerate}
\end{proposition}

%\textcolor{blue}{ The greatest common right divisor and least common left multiple are unique  }

Since $\Fq[x,\sigma]$ is a right (left) Euclidean domain, we have the notion of the greatest common right divisor of polynomials ($\gcrd$) in $\Fq[x,\sigma]$ and  least common left multiple ($\lclm$).\footnote{The analog holds for greatest common left divisor and least common right multiple.} Note that they are unique and can be calculated using the known left and right division algorithms, for more details see \cite{Ore33}.

%\textcolor{purple}{Please use the same numbering for all environments, such that they all increase the counter by one.}
\begin{definition}
 Let $f(x), g(x) \in \Fq[x;\sigma]$. 
 \begin{enumerate}
     \item A monic polynomial $d(x)$ is called the \emph{greatest common right divisor} ($ \gcrd $) of $f(x)$ and $g(x)$ if  $d(x)$ is a right divisor of $f(x)$ and $g(x)$, and any right divisor  $d^{'}(x)$ of $f(x)$ and $g(x)$ is also a  right divisor of $d(x)$.
\item In particular, $f(x)$ and $ g(x) $ are said to be \emph{right coprime} if $\gcrd(f(x), g(x))=1$. 
In this case, there exist polynomials  $a(x), b(x) \in  \Fq[x;\sigma]$ such that $a(x) f(x)+b(x) g(x)=1.$
 \item Similarly, the \emph{least common left multiple} $(\lclm)$ of $f(x)$ and $g(x)$   is defined as the monic polynomial  $h(x)\in  \Fq[x;\sigma]$  of least degree such that $f(x)$ and $ g(x)$ are both  right  divisors of $h(x).$
\end{enumerate}
\end{definition}




\begin{definition}[\cite{Ore33}] Let $a(x), b(x) \in \Fq[x,\sigma]$ two skew polynomials. We say  $a(x)$ is \emph{similar} to $b(x)$ if there exists $u(x) \in \Fq[x,\sigma] $  such that $\lclm(a(x), u(x)) = b(x) u(x)$  and $\gcrd(a(x), u(x)) = 1.$
\end{definition}

\begin{theorem} \cite[Theorem 1]{Ore33}  \label{Factorization}
Every non-zero polynomial $f(x) \in \Fq[x;\sigma]$ factorizes as $f(x)=f_1(x) \cdots f_n(x)$ where $f_i(x) \in \Fq[x;\sigma]$ is irreducible for all $i \in\{1, \ldots, n\}$. Furthermore, if $f(x)=g_1(x) \cdots g_s(x)$ is any other factorization of $f(x)$ as a product of irreducible polynomials $g_i(x) \in \Fq[x;\sigma]$, then $s=n$ and there exists a permutation $\pi:\{1, \ldots, n\} \rightarrow$ $\{1, \ldots, n\}$ such that $f_i$ is similar to $g_{\pi(i)}$. In particular, $f_i$ and $g_{\pi(i)}$ have the same degree for all $i \in\{1, \ldots, n\}$.
\end{theorem}
 

It is important to note that in the skew polynomial ring $\Fq[x;\sigma]$, there are different polynomial evaluations, out of which we will use the following. 

\begin{definition} \cite[p. 310]{Lam1988}\label{def.2}
For $\alpha \in \Fq,$ and $f(x)\in \Fq[x;\sigma]$,
  the \emph{(right) evaluation} (also called  the \emph{(right) remainder evaluation})  of $f(x)$ at $\alpha$ is the element $f(\alpha ):=r \in \Fq$ such that $f(x)= q(x)(x-\alpha)+r.$
\end{definition}

In \cite[Proposition 2.9]{Lam1988}, Lam and Leroy  proved that, for any $\alpha\in \Fq$ we have   
\begin{equation}\label{e1}
f(\alpha)= a_0 + a_1 N^{\sigma}_1(\alpha) +\ldots+ a_{n-1} N^{\sigma}_{n-1}(\alpha)=  \dsum{i=0}{n} a_i N^{\sigma}_i(\alpha)
\end{equation}
where $  N^{\sigma}_{0}(\alpha):=1 ,$ and  for  every $i$ in $\mathbb{N}^*$  
%\hl{I did not look at the source [7] for this, but is there an $s$ in the next equation as written below? Please check this.} 
%\textcolor{blue}{Just a typos, there is no $s$ in this equation.}
$$
N^{\sigma}_i(\alpha):=\sigma^{i-1}(\alpha) \ldots \sigma(\alpha) \alpha.
$$
By taking  $\sigma(a)=a^{p^r} $ for all $a\in \Fq,$ we obtain  for each $ i \in \mathbb{N}^*$  that
\begin{equation}\label{Eqsigma}
N_i^{\sigma}(\alpha)=\alpha^{p^{r (i-1)}} \alpha^{p^{r (i-2)}} \ldots  \alpha^{p^{r}} \alpha =  \alpha^{ \sum_{j=0}^{i-1} p^{rj}}= \alpha^{{\frac{ p^{ri}-1}{p^r-1}}}= \alpha^{[i]_r},
\end{equation}
where $[i]_{r}:=  \dfrac{ p^{ri}-1}{p^r-1}$. Then the  right evaluation of $f(x)=\dsum{i=0}{n} a_i x^i$ at $\alpha\in \Fq$ is given by 
\begin{equation*}
f(\alpha)= \dsum{i=0}{n} a_i N^{\sigma}_i(\alpha)= \dsum{i=0}{n} a_i \alpha^{[i]_r}.
\end{equation*}



%\textcolor{purple}{Do we want to erase the following now?} %\hlc[green]{we can erase it. I made it in comment. }
\begin{comment}
Skew polynomials give rise to a special type of non-associative algebras, called \emph{Petit algebras}, first introduced in 1966 by Petit, see \cite{Brown2018, Petit1967, Petit1968}. 
\begin{definition}
    For any  $f(x),g(x) \in \Fq[x;\sigma]$ denote by $g(x) \bmod_r f(x)$ the remainder of   the right division of $ g(x) $ by $f(x)$. 
    For any  $f(x) \in \Fq[x;\sigma]$  of degree $n,$ the set  
    $$ \Fq[x;\sigma]_f:= \Fq[x;\sigma]/\Fq[x;\sigma]f= \{g \in \Fq[x;\sigma] \mid \deg(g)<n \}, $$
    endowed with the 
multiplication $*$ defined by $$a(x) * b(x):=a(x) b(x) \bmod_r f(x)$$
is a non-associative algebra over $ \Fq^{\sigma} $,   called a \emph{Petit algebra}.
In what follows, keep in mind that when we multiply elements in $ \Fq[x; \sigma]_f, $ the multiplication used is $ *, $ even if we do not explicitly mention it.

\end{definition}
\end{comment}

A skew polynomial $f(x) \in \Fq[x;\sigma]$ is called an \emph{invariant polynomial} if $f(x)\cdot \Fq[x;\sigma] = \Fq[x;\sigma]\cdot f(x),$ i.e., there exist $g(x)$ and $ g'(x) $ such that $ f(x)g(x)= g'(x)f(x).$ In particular,   each polynomial in  $ \Z(\Fq[x;\sigma])=\Fq^{\sigma}[x^{\mu}]$ is   invariant, where $\mu$ is the order of $\sigma$. Note that when $f$ is an invariant polynomial in $\Fq[x;\sigma]$  then $R_f:=\Fq[x;\sigma]/\langle f(x)\rangle$ is an associative algebra over  $ \Fq^{\sigma} $ of dimension $ \deg(f) \mu.$   If $f$ is not invariant, then the center of $ R_f$ is $ \Fq.$

\begin{comment}
\textcolor{purple}{Should we write something here that the quotient rings we will consider for our codes are such Petit algebras?}

\textcolor{blue}{ Even if the polynomial  $f \in \Z(\Fq[x;\sigma])$ (central polynomial) we can't say that our code are  subalgebras of $ R_f:= \Fq[x;\sigma]/\Fq[x;\sigma] f$ over its center $\Fq^{\sigma}.$ I fact, our codes in this case are left ideals $I=\langle g(x)\rangle$ where $g$ a right divisor $f,$ Then $I$ is an $\Fq$-vector subspace of $R_f$ and so an $\Fq^{\sigma}$-vector subspace. But its not not necsssarily clossed under multiplication in $R_f,$ let $ a_1=k_1g$ and $a_2=k_2 g$ we need to prove that $ a_1 a_2 \mod f$ is an element of $I.$ Hence we can't affirm that 
$  a_1a_2=(k_1g) (k_2g) \mod f  $ is a left multiple of $g.$ We can obtain the case if $g^2\mod f \in I$, i.e, $g^2\mod f$ is left multiple of $g.$\\
If if is not central, $R_f$ has the structure of Petit algebra (not associative) over its center $ \Fq^{\sigma},$ in this case our codes are left  $[\Fq[x,\sigma]$-submodules of $R_f$ but as above we can't affirm that they are $\Fq^{\sigma}$-algebras (Petit algebras). }


\textcolor{purple}{We should definitely explain all of this. (And also for ourselves, it should be good to understand all these connections, otherwise what is the point of writing the results here? If we want a morphism, shouldn't we start with an algebra? Or do we not want an algebra morphism but just a ring morphism? Then everything changes and we do not need to care much about the Petit algebra.)\\
Note that most works I know use the semi-field notation for these types of algebras, maybe we should do this, too? Or at least mentions the works of Sheekey et al.  And again, I think there are a lot of results out there that could directly imply that our morphisms are of the form $f(\alpha x)$. \\
And then, since the rank isometries are the linear ones over the subfield, I guess the morphism isometries shoudl be the ones where $\alpha$ is in the subfield, as well. And then we can add the rank results to all the ones we have about the Hamming metric.}

%\textcolor{blue}{Our codes are $\Fq$-vector spaces of lenght $n$ over $\Fq,$ endowed with  additional polycyclic property. This additional property allows us to characterize our codes as $ \Fq[x,\sigma]$-submodules of $ \Fq[x,\sigma]/ \Fq[x,\sigma]f(x)$ for a given monic polynomial $f$ of degree $n.$ This module $ \Fq[x,\sigma]/ \Fq[x,\sigma]f(x)$  is an $\Fq$-vector space of dimension $n$ with basis $\{1,X,\ldots, X^{n-1}\}.$ 
% However, $Fq[x,\sigma]/ \Fq[x,\sigma]f(x)$
%is a Petit algebra over $ \Fq^{\sigma}$ of dimension  $n \mu.$  If the polynomial $f$ is irreducible  then the Petit   algebra  (non associative ) $\Fq[x,\sigma]/\langle f(x) \rangle $ is a division algebra, and it is known also under name of semi-field.  \\
%So for our situation we can't say that our codes are such Petit algebras, since we can't (for instance) find a one to one correspondence between the  $ \Fq[x,\sigma]$-submodules of $ \Fq[x,\sigma]/ \Fq[x,\sigma]f(x)$  and the $ \Fq^{\sigma}$-subalgebra of 
%$ \Fq[x,\sigma]/\Fq[x,\sigma].$ Also if $C=\langle g(x) \rangle$ has dimension $k$ over $\Fq, $ for instance we are not able to calculate its dimension as $ \Fq^{\sigma}$-subalgebra. }

\textcolor{purple}{Exactly, it still seems to me that we do not need the notion of Petit algebra at all then. Because we want to look at the more general module/quotient ring setup. Except -- maybe -- if we consider the rank metric. Do we then just consider the Petit algebra/semifield case, because of the specific isometries?\\
(on the other hand, when we define the set $\mathcal A_\sigma$ later, you say that the code is a left idea, but in your comment you say that we need two-sided ideals (and that's why we need the special structure and the divisibility). Aren't we then back to having Petit algebras? At least implicitly, via the idealiser or something like that?}


\textcolor{blue}{In the comment about  $ A_{\sigma},$ if the polynomial $ f(x)$ has coefficients in $A_{\sigma}$ then it is generates a two sided ideal $\langle f(x) \rangle $ in $ \Fq[x,\sigma]$ and so $ R_f:=\Fq[x,\sigma]/\langle f(x) \rangle $ is a (non-commutative ) ring, and it is  an associative algebra  over $ \Fq^{\sigma},$ but Not commutative. So,  we speak about left ideal and right ideal of this ring. \\
In our case we considered right division of polynomials so we speak about left ideals of this rings. We can also speak about right ideal of we consider the left division of polynomials. The structure of Petit algebra concerns the associativity, if the polynomial is central (two sided) then $R_f$ is associative and not associative when $f$ is not central.
}

\hlc[green]{ In the other hand I recall some results  skew polynomials used by J. Sheeky and his co-authors, \cite{JohnSheekey2016,sheekey2020new,lavrauw2013semifields}, to construct MRD codes}  
\textcolor{blue}{
They first started since 2011 by studying the   notion  of isotopism between two semifields $R_{f_1} $ and $R_{f_2},$ with $f_1$ and $f_2$ are irreducible.
An isotopism from $(\mathbb{S}, \circ)$ to $\left(\mathbb{S}^{\prime}, \circ^{\prime}\right)$ is a triple $(F, G, H)$ of bijections from $\mathbb{S}$ to $\mathbb{S}^{\prime}$, linear over the characteristic field of $\mathbb{S}$, such that
$$
F(a) \circ^{\prime} G(b)=H(a \circ b)
$$
So they proved that $\varphi_{\alpha},$ with $ \alpha\in \Fq^{\sigma},$ of the form
\begin{center}
    $\varphi(a(x)) =  a^{\rho}(\alpha x) $
\end{center}
is an is an  isotopism from  between the semifields $R_{f_1} $ and $R_{f_2}.$  For more details see \href{  https://arxiv.org/abs/1201.2528 }{Semifields from skew polynomial rings} ,  \cite[Section 5]{lavrauw2013semifields}.}
\hlc[green]{Recently they constructed many MRD codes from skew polynomials and $\sigma$-linearized polynomials.}
\textcolor{blue}{
Let  $f\in \Fq[ x;\sigma],$ 
\begin{enumerate}
    \item There exist unique two sided polynomial $f^*\in \Fq^{\sigma}[x^{\mu}],$ called the bound of $f,$ such that $f$ right divides $f^*$ in $ \Fq[ x;\sigma]$ (The existence is guaranteed  over finite fields.) 
    \item If $f$ is two sided then it coincides with its bound $f^*$
    \item For $y=x^\mu$ the bound $ f^*(y)\in \Fq^{\sigma}[y]  $ is commutative polynomial.
    \item If $f$ is irreducible in $ \Fq[x,\sigma]$ then its bound $ f^*(y)$ is irreducible as a polynomial in $\Fq^{\sigma}[y].$ 
    \item Conversely, if $ f^*(y)$ is irreducible as a polynomial in $\Fq^{\sigma}[y]$ of degree $s$ in $y,$ then
    \begin{enumerate}
        \item $ \Fq[x,\sigma]/\langle f^*(x^{\mu}) \rangle  \cong \mathbb M_{\mu}( \mathbb{F}_{q_0^s}) ,$
where $ \mathbb M_{\mu}( \mathbb{F}_{q_0^s})$ is the algebra of square matrices with entries in $ \mathbb{F}_{q_0^s}:=\Fq^{\sigma}.$
        \item Any right divisor of $ f$ in $ \Fq[x,\sigma]$ has degree $s.$
        \item The rank of any  \textbf{irreducible right divisor  $g(x)$}  of $f^*(x^\mu)$ (division in $\Fq[x,\sigma]$) in $\mathbb M_{\mu}( \mathbb{F}_{q_0^s})$ is 
        $$ \rank( g ) = \mu-1.$$
        \item  More general, \cite[Proposition 7]{sheekey2020new}, the rank of any elements  $g(x)$ ($g$ is not necessary a devisor of $f$) in $\mathbb M_{\mu}( \mathbb{F}_{q_0^s})$ is 
        $$\operatorname{rank}(g)=\mu-\frac{1}{s} \deg(\gcrd(g, f^*(x^\mu)).$$
         \end{enumerate}
          \end{enumerate}}


          
   \textcolor{blue}{       
   \begin{enumerate}      
  \item Suppose that $\sigma \in \mathbb{F}_{q^n}$ is of order $n,$ and  $f^*(x^n)$ in $\Fq[x^n]$ is irreducible ($\mathbb{F}_{q^n}^{\sigma}=\Fq).$)
  \begin{enumerate}
  \item They  considered in \cite{sheekey2020new},  $ \mathbb{F}_{q^n}/\Fq  $ and they took  with $\deg(f^*)=s$,and they proved 
$$
R_f:=\frac{\Fq[x,\sigma]}{\langle f^*(x^n)\rangle} \simeq \operatorname{End}_{ \mathbb{F}_{q^s}}\left( \mathbb{F}_{q^n}\right) \simeq M_n\left(\mathbb{F}_{q^s}\right) \simeq \textcolor{red}{ \mathbb{L}_{q^n,q}[x]/\mathbb{L}_{q^n,q}[x]( f^*(x^{q^n}) )}.
$$
 \textcolor{red}{The isomorphic  in red is proved when $ f^*(y)=y-1,$ and so we can prove it for our case for any irreducible polynomials $f^*(y)\in \Fq[y].$}
 \item In \cite[Theorem 7]{sheekey201913}, they proved  under conditions that the set
 $$ S_{n, s, k}(\eta, \rho, f^*):=\left\{a(x)+Rf^*(x^n): \deg(a(x)) \leqslant k s, \ a_{k s}=\eta \rho(a_0)\right\}$$
 is an MRD code with  minimum distance $n-k + 1,$ where $\rho$ an automorphism of $\mathbb{F}_{q^n}.$
 \item   \textcolor{red}{ It mentioned in \cite[Remark 11]{sheekey2020new} that it is possible for different choices of $f^*(y)$ we can produce different equivalent codes.}
\item So when $f^*(y)=y-1$, they got  that  $ M_n(\Fq) \simeq \frac{\mathbb F_{q^n}[x ; \sigma]}{\left(x^n-1\right)}$. So any irreducible factor $g$ of $x^n-1$ should be of degree $1,$ and generates a skew cyclic code $C$ of length $n$ over $\mathbb{F}_{q^n}$ and dimension $n-1.$ When we consider $ C \subseteq M_n\left(\mathbb{F}_{q}\right),$ the  its is of  length $n^2$ and dimension $n(n-1)$ over $\Fq$  and of min rank distance $ n-1.$ Finally we verify that $C$ is and  MRD code.\\
In connection  to what did in section 3, if we consider $\alpha$ a generator of the extension $ \mathbb{F}_{q^n}/\Fq$ and $ \beta=\alpha^{-1}\sigma(\alpha),$ we get the irreducible factorization of $x^n-1$ in $\Fq[x,\sigma]$ is given by
$$ x^n-1=(x-\beta)(x-\sigma(\beta)) \ldots (x-\sigma^{n-1}(\beta))$$
So the irreducible divisors of $x^n-1$ are $(x-\sigma^i(\beta))$ and so the rank of any one of them in $ M_n(\Fq) $  is $ n-1. $ Hence it generates a rank code of length $ n$ over $ \mathbb{F}_{q^n}$ and of dimension $ k=n-1$ and with minimal rank distance $ d_R=n-1.$ We verify that $n-k+1=d   $
and so it  generates an MRD code over $ \mathbb{F}_{q^n}.$  \\
\item With the same conditions.  In our situation (trinomial case) with the polynomial $y^p-y-1 \in \Fq[y]$ irreducible polynomial over $ \mathbb{F}_{q}:= \mathbb{F}_{q^n}^{\sigma},$ (In our first paper we have some conditions on irreducibility of trinomials)
   and then any irreducible divisor of  $ f= x^{pn}-x^{n}-1 $ in $\mathbb{F}_{q^n}[x;\sigma]$ is and   MRD code of length $ pn$ and dimension $n(p-1)$
   over $ \mathbb{F}_{q^p}$ where $p$ is the characteristic  of $\Fq. $ 
     \end{enumerate}
 \item \textcolor{red}{In  relation to the works of J. Sheeky we can interest in the case when $ f^*(y)$ is not irreducible, so we can decompose it into irreducible factors, and then decompose the code in some direct sum..., Maybe this question is far from this project we can do it in another project. Also the study with sum-rank distance should be interesting.  } Just in this context J. sheeky generlized his idea in constructing MRD code to the case of code over $\Fq(t)$ the rational fonctional field \href{https://arxiv.org/pdf/2502.13531}{Quotients of skew polynomial rings new construction}, and the problem of studying non-equivalence here is open according to what people did on equivalence on codes over  $\Fq.$     
    \end{enumerate}   }
    
\hlc[green]{Back to our situation }
\textcolor{blue}{
\begin{enumerate}    
\item  Let $f_1, f_2$ in $ \Fq[x,\sigma],$ and  let $ \vatphi$ be an $\Fq^{\sigma}$-algebra  
$$\varphi ~:~ R_{f_1}:=\Fq[x,\sigma]/ \Fq[x,\sigma] f_1 \quad  \to \quad R_{f_2}:= \Fq[x,\sigma]/ \Fq[x,\sigma]f_2,  $$ 
\begin{itemize}
    \item If  $f_1$ and $f_2$ are central, i.e with coefficient in $\mathcal{A}_{\sigma}$,  $R_{f_1}, R_{f_2}$ are rings and ($\Fq^{\sigma}$-algebras). So if $\varphi$ is an $\Fq^{\sigma}$-algebra isomorphism between $R_{f_1}$ and $R_{f_2}$ then it is an ring isomorphism. Then
    \begin{itemize}
        \item It  is a rank-isometry (not necessarily $\Fq$-linear but $\Fq^{\sigma}$-linear). Let $C$ generated by $g(x)$ then $f_1= h(x)g(x). $ So as   $\varphi $ is an $\Fq^{\sigma}$-algebra isomorphism  (the multiplication in the ring is the same as the multiplication in the $\Fq^{\sigma}$-algebra), we obtain that $ \varphi(g)$ divides $ f_2$ in $\Fq[x,\sigma].$ Then $ \varphi(C) $ is an ideal of $ R_{f_2}$ generated by $ \varphi(g_2)\in \Fq[x,\sigma].$ As $\Fq \subset R_{f_2}$ and  $ \varphi(C)$ is an ideal, then  $ \varphi(C)$  is $ \Fq$-linear. Hence $$ \dim_{\Fq}(C)= [\Fq : \Fq^{\sigma}] \dim_{\Fq^{\sigma}}(C) = [\Fq : \Fq^{\sigma}]   \dim_{\Fq^{\sigma}}(\varphi(C))=  \dim_{\Fq}(\varphi(C)). $$
        \item To become a Hamming isometry, in addition of being   $\Fq^{\sigma}$-algebra isomorphism $ \varphi$ need too be of the form
        $$ \varphi(x)= \alpha x, \qud \text{where } \alpha\in \Fq^{\sigma}).$$
      \end{itemize}
\item If  $f_1$ and $f_2$ are not central,  $R_{f_1}, R_{f_2}$ are Petit $\Fq^{\sigma}$-algebras (but not rings). So if $\varphi$ is an $\Fq^{\sigma}$-algebra isomorphism between the Petit algebras  $R_{f_1}$ and $R_{f_2}.$ Then should verify the following points:
\begin{itemize}
    \item $ \varphi(a(x)+b(x))= \varphi(a(x))+ \varphi(b(x)).$
 \item $ \varphi(a(x) *b(x))= \varphi(a(x))) *( \varphi(b(x)).$
    \item  For any $\lambda \in \Fq^{\sigma}$ we have  $ \varphi(\lambda a(x))= \lambda   \varphi(a(x)),$ since it is $\Fq^{\sigma}$-algebra morphism.
\end{itemize}
In this case,  let $C$ generated by $g(x)$ then $f_1= h(x)g(x), $ the $g(x)$ genrtes a left $\Fq[x,\sigma]$-submodules of the $ R_{f_1.}$ As   $\varphi $ is an $\Fq^{\sigma}$-algebra isomorphism  we obtain that $ \varphi(g)$ divides $ f_2$ in $\Fq[x,\sigma],$ since $ f_2=\varphi(f_1)= \varph(g)\varphi(h).$ 
Then $ \varphi(C) $ is a left $\Fq[x,\sigma]$ module  of $ R_{f_2}$ generated by $ \varphi(g_2)\in \Fq[x,\sigma].$ As $\Fq \subset \Fq[x,\sigma]$ and  $ \varphi(C)$ is a left $\Fq[x,\sigma]$-submodule of $R_{f_2}$ , then  $ \varphi(C)$  is $ \Fq$-linear.
Hence $$ \dim_{\Fq}(C)= [\Fq : \Fq^{\sigma}] \dim_{\Fq^{\sigma}}(C) = [\Fq : \Fq^{\sigma}]   \dim_{\Fq^{\sigma}}(\varphi(C))=  \dim_{\Fq}(\varphi(C)). $$
    \item To become a Hamming isometry, in addition of being   $\Fq^{\sigma}$-algebra isomorphism $ \varphi$ need too be of the form
        $$ \varphi(x)= \alpha x, \qud \text{where } \alpha\in \Fq^{\sigma}). $$ 
\end{itemize} 
\item In our approach, I think we don't need to verify the $\Fq$-morphism property if we considered $\Fq^{\sigma}$-algebras isomorphism. But in our case in the definition we defined our isometry to be $\Fq$-linear isomorphism and satisfying $\Fq$-morphism property ( $\varphi(a*b)= \varphi(a)* \varphi(b))$), but we can recover this property if we consider $\Fq^{\sigma}$-algebras and we lose the genrality of considering $ \alpha \in \Fq.$ So what do you think which approach we consider?
\end{enumerate}
}

\end{comment}

\begin{remark}
There is a well-known connection between skew polynomials and linearized polynomials (precisely $\sigma$-polynomials), and the latter have been studied extensively in the context of rank-metric codes, see, e.g.,
%In the following, we mention the relation between skew polynomial rings and the ring of linearized polynomials over finite fields. The latter play an important role in the study of rank-metric codes, 
\cite{Gabidulin1985,ko08,ore1933special,wu2013linearized}. 
For this let $\sigma$ be an automorphism of order $\mu,$ and $\mathbb{F}_{q_0}=\Fq^{\sigma}$ its fixed subfield.
The (non-commutative) ring of   $\sigma$-polynomials, $ (\mathcal{L}_{q,\sigma}[x],+,\circ),$  is  the subset  of the polynomial ring $\Fq[x]$ defined by
$$
\mathcal{L}_{q,\sigma}[x]:=\left\{\sum_{i=0}^n a_i x^{\sigma^i} := \sum_{i=0}^n a_i x^{q_0^{i}}\mid n \in \mathbb{N}, a_i \in \Fq \right\},
$$
endowed with  the usual addition $"+"$ and  composition $"\circ"$ of polynomials. Note that  $\Fq[x ; \sigma]$ and $\mathcal{L}_{q,\sigma}[x]$ are
isomorphic rings %$\mathbb{F}_{q_0}$-algebras (or simply ring isomorphisms) 
via 
$$
\Lambda: \Fq[x ; \sigma] \longrightarrow \mathcal{L}_{q,\sigma}[x], \sum_{i=0}^n a_i x^i \longmapsto \sum_{i=0}^n a_i x^{\sigma^i}=\sum_{i=0}^n a_i x^{q_0^{i}} .
$$
Analogously, we get the following isomorphism 
%between the $ \mathbb{F}_{q_0}$-algebras  (or simply ring isomorphism), see \cite[Theorem 2.2.]{wu2013linearized} 
$$ \Fq[x ; \sigma]/\langle x^{\mu}-1\rangle \cong \mathcal{L}_{q,\sigma}[x]/\langle x^{q_0^{\mu}}-x\rangle .$$ % \cong \mathbb{F}_{q_0}^{\mu \times \mu}
The above correspondence allows us to characterize skew cyclic codes of length $\mu $ over $ \Fq$ as left ideals of $\mathcal{L}_{q,\sigma}[x]/\langle x^{q_0^{\mu}}-x\rangle,$ for more details see \cite{gabidulin2009rank} and more recently \cite[Section 2.4]{martinez2017roots}. More generally, we can prove 
for any  skew polynomial $f \in \mathbb{F}_{q_0}[x,\sigma]$ of degree $ \mu$  ($f$ is defined over $\mathbb F_{q_0}$ to generate a two sided ideal in $\Fq[x;\sigma])$ the following   ring isomorphism 
$$ \Fq[x ; \sigma]/\langle f(x)\rangle \cong \mathcal{L}_{q,\sigma}[x]/\langle f(x^{q_0}) \rangle.$$ 
Note that, seen as an algebra isomorphism, this relation holds over $ \mathbb{F}_{q_0}.$
%\hlc[green]{I apologize I wasn't clear here. Even if  $f$ is central the algebra structure of $ \Fq[x ; \sigma]/\langle f(x)\rangle $ is  over $\mathbb{F}_{q_0}.$ The difference is in the type of $\mathbb{F}_{q_0}$-algebras: If $f$ central it is an associative algebra, else ( $f$ not central) it is non-associative algebra over $ \mathbb{F}_{q_0}$ (called Petit algebra). But we can prove in both cases  that if the we have an  isomorphism of  $\mathbb{F}_{q_0}-$algebras ( associative or not ), then the isomorphism is an $\Fq-$vector spaces isomorphism.}
% \textcolor{purple}{Now we do not define Petit algebras any more, hence I would also not use it here.
%For example, we could write: " Note that, seen as an algebra isomorphism, this relation holds over $ \mathbb{F}_{q_0}$, except in the case where $f$ is central, in which case it is also an $\Fq$-algebra isomorphism."}
\begin{comment}
    \textcolor{blue}{
\begin{enumerate}    
\item  Let $f_1, f_2$ in $ \Fq[x,\sigma],$ and  let $ \vatphi$ be an $\Fq^{\sigma}$-algebra  
$$\varphi ~:~ R_{f_1}:=\Fq[x,\sigma]/ \Fq[x,\sigma] f_1 \quad  \to \quad R_{f_2}:= \Fq[x,\sigma]/ \Fq[x,\sigma]f_2,  $$ 
\begin{itemize}
    \item If  $f_1$ and $f_2$ are central, i.e with coefficient in $\mathcal{A}_{\sigma}$,  $R_{f_1}, R_{f_2}$ are rings and ($\Fq^{\sigma}$-algebras). So if $\varphi$ is an $\Fq^{\sigma}$-algebra isomorphism between $R_{f_1}$ and $R_{f_2}$ then it is an ring isomorphism. Then
    \begin{itemize}
        \item It  is a rank-isometry (not necessarily $\Fq$-linear but $\Fq^{\sigma}$-linear). Let $C$ generated by $g(x)$ then $f_1= h(x)g(x). $ So as   $\varphi $ is an $\Fq^{\sigma}$-algebra isomorphism  (the multiplication in the ring is the same as the multiplication in the $\Fq^{\sigma}$-algebra), we obtain that $ \varphi(g)$ divides $ f_2$ in $\Fq[x,\sigma].$ Then $ \varphi(C) $ is an ideal of $ R_{f_2}$ generated by $ \varphi(g_2)\in \Fq[x,\sigma].$ As $\Fq \subset R_{f_2}$ and  $ \varphi(C)$ is an ideal, then  $ \varphi(C)$  is $ \Fq$-linear. Hence $$ \dim_{\Fq}(C)= [\Fq : \Fq^{\sigma}] \dim_{\Fq^{\sigma}}(C) = [\Fq : \Fq^{\sigma}]   \dim_{\Fq^{\sigma}}(\varphi(C))=  \dim_{\Fq}(\varphi(C)). $$
        \item To become a Hamming isometry, in addition of being   $\Fq^{\sigma}$-algebra isomorphism $ \varphi$ need too be of the form
        $$ \varphi(x)= \alpha x, \qud \text{where } \alpha\in \Fq^{\sigma}).$$
      \end{itemize}
\item If  $f_1$ and $f_2$ are not central,  $R_{f_1}, R_{f_2}$ are Petit $\Fq^{\sigma}$-algebras (but not rings). So if $\varphi$ is an $\Fq^{\sigma}$-algebra isomorphism between the Petit algebras  $R_{f_1}$ and $R_{f_2}.$ Then should verify the following points:
\begin{itemize}
    \item $ \varphi(a(x)+b(x))= \varphi(a(x))+ \varphi(b(x)).$
 \item $ \varphi(a(x) *b(x))= \varphi(a(x))) *( \varphi(b(x)).$
    \item  For any $\lambda \in \Fq^{\sigma}$ we have  $ \varphi(\lambda a(x))= \lambda   \varphi(a(x)),$ since it is $\Fq^{\sigma}$-algebra morphism.
\end{itemize}
In this case,  let $C$ generated by $g(x)$ then $f_1= h(x)g(x), $ the $g(x)$ genrtes a left $\Fq[x,\sigma]$-submodules of the $ R_{f_1.}$ As   $\varphi $ is an $\Fq^{\sigma}$-algebra isomorphism  we obtain that $ \varphi(g)$ divides $ f_2$ in $\Fq[x,\sigma],$ since $ f_2=\varphi(f_1)= \varph(g)\varphi(h).$ 
Then $ \varphi(C) $ is a left $\Fq[x,\sigma]$ module  of $ R_{f_2}$ generated by $ \varphi(g_2)\in \Fq[x,\sigma].$ As $\Fq \subset \Fq[x,\sigma]$ and  $ \varphi(C)$ is a left $\Fq[x,\sigma]$-submodule of $R_{f_2}$ , then  $ \varphi(C)$  is $ \Fq$-linear.
Hence $$ \dim_{\Fq}(C)= [\Fq : \Fq^{\sigma}] \dim_{\Fq^{\sigma}}(C) = [\Fq : \Fq^{\sigma}]   \dim_{\Fq^{\sigma}}(\varphi(C))=  \dim_{\Fq}(\varphi(C)). $$
    \item To become a Hamming isometry, in addition of being   $\Fq^{\sigma}$-algebra isomorphism $ \varphi$ need too be of the form
        $$ \varphi(x)= \alpha x, \qud \text{where } \alpha\in \Fq^{\sigma}). $$ 
\end{itemize} 
\item In our approach, I think we don't need to verify the $\Fq$-morphism property if we considered $\Fq^{\sigma}$-algebras isomorphism. But in our case in the definition we defined our isometry to be $\Fq$-linear isomorphism and satisfying $\Fq$-morphism property ( $\varphi(a*b)= \varphi(a)* \varphi(b))$), but we can recover this property if we consider $\Fq^{\sigma}$-algebras and we lose the genrality of considering $ \alpha \in \Fq.$ So what do you think which approach we consider?
\end{enumerate}
}
\end{comment}


\end{remark}

\begin{comment}
\textcolor{blue}{
\begin{enumerate}    
\item  Let $f_1, f_2$ in $ \Fq[x,\sigma],$ and  let $ \vatphi$ be an $\Fq^{\sigma}$-algebra  
$$\varphi ~:~ R_{f_1}:=\Fq[x,\sigma]/ \Fq[x,\sigma] f_1 \quad  \to \quad R_{f_2}:= \Fq[x,\sigma]/ \Fq[x,\sigma]f_2,  $$ 
\begin{itemize}
    \item If  $f_1$ and $f_2$ are central, i.e with coefficient in $\mathcal{A}_{\sigma}$,  $R_{f_1}, R_{f_2}$ are rings and ($\Fq^{\sigma}$-algebras). So if $\varphi$ is an $\Fq^{\sigma}$-algebra isomorphism between $R_{f_1}$ and $R_{f_2}$ then it is an ring isomorphism. Then
    \begin{itemize}
        \item It  is a rank-isometry (not necessarily $\Fq$-linear but $\Fq^{\sigma}$-linear). Let $C$ generated by $g(x)$ then $f_1= h(x)g(x). $ So as   $\varphi $ is an $\Fq^{\sigma}$-algebra isomorphism  (the multiplication in the ring is the same as the multiplication in the $\Fq^{\sigma}$-algebra), we obtain that $ \varphi(g)$ divides $ f_2$ in $\Fq[x,\sigma].$ Then $ \varphi(C) $ is an ideal of $ R_{f_2}$ generated by $ \varphi(g_2)\in \Fq[x,\sigma].$ As $\Fq \subset R_{f_2}$ and  $ \varphi(C)$ is an ideal, then  $ \varphi(C)$  is $ \Fq$-linear. Hence $$ \dim_{\Fq}(C)= [\Fq : \Fq^{\sigma}] \dim_{\Fq^{\sigma}}(C) = [\Fq : \Fq^{\sigma}]   \dim_{\Fq^{\sigma}}(\varphi(C))=  \dim_{\Fq}(\varphi(C)). $$
        \item To become a Hamming isometry, in addition of being   $\Fq^{\sigma}$-algebra isomorphism $ \varphi$ need too be of the form
        $$ \varphi(x)= \alpha x, \qud \text{where } \alpha\in \Fq^{\sigma}).$$
      \end{itemize}
\item If  $f_1$ and $f_2$ are not central,  $R_{f_1}, R_{f_2}$ are Petit $\Fq^{\sigma}$-algebras (but not rings). So if $\varphi$ is an $\Fq^{\sigma}$-algebra isomorphism between the Petit algebras  $R_{f_1}$ and $R_{f_2}.$ Then should verify the following points:
\begin{itemize}
    \item $ \varphi(a(x)+b(x))= \varphi(a(x))+ \varphi(b(x)).$
 \item $ \varphi(a(x) *b(x))= \varphi(a(x))) *( \varphi(b(x)).$
    \item  For any $\lambda \in \Fq^{\sigma}$ we have  $ \varphi(\lambda a(x))= \lambda   \varphi(a(x)),$ since it is $\Fq^{\sigma}$-algebra morphism.
\end{itemize}
In this case,  let $C$ generated by $g(x)$ then $f_1= h(x)g(x), $ the $g(x)$ genrtes a left $\Fq[x,\sigma]$-submodules of the $ R_{f_1.}$ As   $\varphi $ is an $\Fq^{\sigma}$-algebra isomorphism  we obtain that $ \varphi(g)$ divides $ f_2$ in $\Fq[x,\sigma],$ since $ f_2=\varphi(f_1)= \varph(g)\varphi(h).$ 
Then $ \varphi(C) $ is a left $\Fq[x,\sigma]$ module  of $ R_{f_2}$ generated by $ \varphi(g_2)\in \Fq[x,\sigma].$ As $\Fq \subset \Fq[x,\sigma]$ and  $ \varphi(C)$ is a left $\Fq[x,\sigma]$-submodule of $R_{f_2}$ , then  $ \varphi(C)$  is $ \Fq$-linear.
Hence $$ \dim_{\Fq}(C)= [\Fq : \Fq^{\sigma}] \dim_{\Fq^{\sigma}}(C) = [\Fq : \Fq^{\sigma}]   \dim_{\Fq^{\sigma}}(\varphi(C))=  \dim_{\Fq}(\varphi(C)). $$
    \item To become a Hamming isometry, in addition of being   $\Fq^{\sigma}$-algebra isomorphism $ \varphi$ need too be of the form
        $$ \varphi(x)= \alpha x, \qud \text{where } \alpha\in \Fq^{\sigma}). $$ 
\end{itemize} 
\item In our approach, I think we don't need to verify the $\Fq$-morphism property if we considered $\Fq^{\sigma}$-algebras isomorphism. But in our case in the definition we defined our isometry to be $\Fq$-linear isomorphism and satisfying $\Fq$-morphism property ( $\varphi(a*b)= \varphi(a)* \varphi(b))$), but we can recover this property if we consider $\Fq^{\sigma}$-algebras and we lose the genrality of considering $ \alpha \in \Fq.$ So what do you think which approach we consider?
\end{enumerate}
}



\textcolor{blue}{Similarly to  \cite{gabidulin2009rank, martinez2017roots}, by using this later isomorphism, we will be able to define different version of the Gabidulin  codes $ G_k$ ( $\mathbb{F}_{q_0}$-Frobenius automorphism) and generalized Gabidulin codes (different automorphism with same order as the $q_0$-Frobenius), these classes are defined as $ \mathbb{F}_{q}$-linear subspaces  of $\mathcal{L}_{q,\sigma}[x]/\langle x^{q_0^\mu}-x \rangle,$ which is isomorphic to $ \Fq[x,\sigma]/\langle x^\mu-1 \rangle.$ If we changes the $x^\mu-1$ by another polynomial $f\in\mathbb{F}_{q_0}[x,\sigma]$ we will obtain differents class of MRD codes and also studying their non-equivalence.   In this context of semi-fields with skew polynomials  there this recent preprint using $x^{\mu}-\lambda$ and they proved the choice of different generator of the extension fields may leads to the non-equivalence semi-fields.  \href{https://arxiv.org/pdf/2502.00770}{https://arxiv.org/pdf/2502.00770}.\\
To study the dual of these codes people use the adjoint of the linearized polynomial $ a=a_0x+a_1x^{q_0}+\ldots+ a_{k}x^{q_0^{k-1}}$ defined by
$$\widehat{a}= a_{k}x+a^{q_0}_{k-1}x^{q_0}+\ldots+ a^{q_0^{k-1}}_{0}x^{q_0^{k-1}} ,$$  and the corespendant skew polynomial is called the skew reciprocal polynomial. For $a=\dsum{i=0}{k-1} a_i x^i\in \Fq[x,\sigma]$ its $\sigma$-reciprocal polynomial is defined by
$$ a^*(x)=\dsum{i=0}{k-1} \sigma^i(a_{k-i}) x^i\in \Fq[x,\sigma].$$
}
\hlc[green]{About Isometry: }
\textcolor{blue}{
\begin{itemize}   
\item  We can also studied Equivalence between two classes of skew polycyclic codes bu considering linearized polynomials. Studying isometry between  $ LR_{f_1}:=\mathcal{L}_{q,\sigma}[x]/\langle L(f_1) \rangle$  and $LR_{f_1}:=\mathcal{L}_{q,\sigma}[x]/\langle L(f_2),$ with $L(f_i)$ the associated linearized polynomial to $f_i.$
\item Let  $$\varphi : \ LR_{f_1} \to LR_{f_1} \  ; a(x^{q_0}) \longrightarrow \varphi(a(x^{q_0})) $$
We seek for conditions on which  $ \varphi  $ be a rank-isometry. According to the existing works on rank metric codes $\varphi $ should be of the from 
$$ \varphi(x^{iq_0}) :=  P(x^{q_0}) \circ x^{iq_0} \circ Q(x^{q_0}) \quad \text{ and } \quad  \varphi( a(x^{q_0})\circ b(x^{q_0}))=  \varphi( a(x^{q_0}))\circ  \varphi( b(x^{q_0})),   $$
where $ P(x^{q_0}), Q(x^{q_0}) $ are invertible ${q_0}$-linearized  polynomials.
\item For instance if $ n\leq \mu,$  and 
$ f_1$ and $f_2$ are two sided irreducible polynomials as element of $ \Fq^{\sigma}[x^\mu].$  Then for any non zero 
Then the class of $f_1$-polyclic codes is equivalent to the class of $f_2$-polycylic codes.
\item In the point of view of skew polynomials: I  if $ n\leq \mu,$ 
$ f_1$ and $f_2$ are two sided and  irreducible polynomials, then
$$\varphi : \ \Fq[x;;\sigma]/\Fq[x;;\sigma] f_1 \to  \Fq[x;;\sigma]/\Fq[x;;\sigma] f_2 \  ; $$
is a rank-isometry if its is of the form
$$ \varphi(a(x)) = P(x) \circ a(x)\circ Q(x)$$
with $P$ and $Q$ are invertible in $  \Fq[x;;\sigma]/\Fq[x;;\sigma] f_2 ,$ which means 
$$ \rank( P)=  \rank( Q)=\mu-\gcrd( Q,f_2(x^\mu))= \mu$$
which means that $ P, Q$ should by coprime with $ f_2(x^\mu).$ 
\end{itemize}
}

\textcolor{purple}{Should we write something here that the quotient rings we will consider for our codes are such Petit algebras?}

\textcolor{blue}{ Even if the polynomial  $f \in \Z(\Fq[x;\sigma])$ (central polynomial) we can't say that our code are  subalgebras of $ R_f:= \Fq[x;\sigma]/\Fq[x;\sigma] f$ over its center $\Fq^{\sigma}.$ I fact, our codes in this case are left ideals $I=\langle g(x)\rangle$ where $g$ a right divisor $f,$ Then $I$ is an $\Fq$-vector subspace of $R_f$ and so an $\Fq^{\sigma}$-vector subspace. But its not not necsssarily clossed under multiplication in $R_f,$ let $ a_1=k_1g$ and $a_2=k_2 g$ we need to prove that $ a_1 a_2 \mod f$ is an element of $I.$ Hence we can't affirm that 
$  a_1a_2=(k_1g) (k_2g) \mod f  $ is a left multiple of $g.$ We can obtain the case if $g^2\mod f \in I$, i.e, $g^2\mod f$ is left multiple of $g.$\\
If if is not central, $R_f$ has the structure of Petit algebra (not associative) over its center $ \Fq^{\sigma},$ in this case our codes are left  $[\Fq[x,\sigma]$-submodules of $R_f$ but as above we can't affirm that they are $\Fq^{\sigma}$-algebras (Petit algebras). }


\textcolor{purple}{We should definitely explain all of this. (And also for ourselves, it should be good to understand all these connections, otherwise what is the point of writing the results here? If we want a morphism, shouldn't we start with an algebra? Or do we not want an algebra morphism but just a ring morphism? Then everything changes and we do not need to care much about the Petit algebra.)\\
Note that most works I know use the semi-field notation for these types of algebras, maybe we should do this, too? Or at least mentions the works of Sheekey et al.  And again, I think there are a lot of results out there that could directly imply that our morphisms are of the form $f(\alpha x)$. \\
And then, since the rank isometries are the linear ones over the subfield, I guess the morphism isometries shoudl be the ones where $\alpha$ is in the subfield, as well. And then we can add the rank results to all the ones we have about the Hamming metric.}

\textcolor{blue}{Our codes are $\Fq$-vector spaces of lenght $n$ over $\Fq,$ endowed with  additional polycyclic property. This additional property allows us to characterize our codes as $ \Fq[x,\sigma]$-submodules of $ \Fq[x,\sigma]/ \Fq[x,\sigma]f(x)$ for a given monic polynomial $f$ of degree $n.$ This module $ \Fq[x,\sigma]/ \Fq[x,\sigma]f(x)$  is an $\Fq$-vector space of dimension $n$ with basis $\{1,X,\ldots, X^{n-1}\}.$ 
 However, $Fq[x,\sigma]/ \Fq[x,\sigma]f(x)$
is a Petit algebra over $ \Fq^{\sigma}$ of dimension  $n \mu.$  If the polynomial $f$ is irreducible  then the Petit   algebra  (non associative ) $\Fq[x,\sigma]/\langle f(x) \rangle $ is a division algebra, and it is known also under name of semi-field.  \\
So for our situation we can't say that our codes are such Petit algebras, since we can't (for instance) find a one to one correspondence between the  $ \Fq[x,\sigma]$-submodules of $ \Fq[x,\sigma]/ \Fq[x,\sigma]f(x)$  and the $ \Fq^{\sigma}$-subalgebra of 
$ \Fq[x,\sigma]/\Fq[x,\sigma].$ Also if $C=\langle g(x) \rangle$ has dimension $k$ over $\Fq, $ for instance we are not able to calculate its dimension as $ \Fq^{\sigma}$-subalgebra. }

\textcolor{purple}{Exactly, it still seems to me that we do not need the notion of Petit algebra at all then. Because we want to look at the more general module/quotient ring setup. Except -- maybe -- if we consider the rank metric. Do we then just consider the Petit algebra/semifield case, because of the specific isometries?\\
(on the other hand, when we define the set $\mathcal A_\sigma$ later, you say that the code is a left idea, but in your comment you say that we need two-sided ideals (and that's why we need the special structure and the divisibility). Aren't we then back to having Petit algebras? At least implicitly, via the idealiser or something like that?}


\textcolor{blue}{In the comment about  $ A_{\sigma},$ if the polynomial $ f(x)$ has coefficients in $A_{\sigma}$ then it is generates a two sided ideal $\langle f(x) \rangle $ in $ \Fq[x,\sigma]$ and so $ R_f:=\Fq[x,\sigma]/\langle f(x) \rangle $ is a (non-commutative ) ring, and it is  an associative algebra  over $ \Fq^{\sigma},$ but Not commutative. So,  we speak about left ideal and right ideal of this ring. \\
In our case we considered right division of polynomials so we speak about left ideals of this rings. We can also speak about right ideal of we consider the left division of polynomials. The structure of Petit algebra concerns the associativity, if the polynomial is central (two sided) then $R_f$ is associative and not associative when $f$ is not central.
}

\hlc[green]{ In the other hand I recall some results  skew polynomials used by J. Sheeky and his co-authors, \cite{JohnSheekey2016,sheekey2020new,lavrauw2013semifields}, to construct MRD codes}  
\textcolor{blue}{
They first started since 2011 by studying the   notion  of isotopism between two semifields $R_{f_1} $ and $R_{f_2},$ with $f_1$ and $f_2$ are irreducible.
An isotopism from $(\mathbb{S}, \circ)$ to $\left(\mathbb{S}^{\prime}, \circ^{\prime}\right)$ is a triple $(F, G, H)$ of bijections from $\mathbb{S}$ to $\mathbb{S}^{\prime}$, linear over the characteristic field of $\mathbb{S}$, such that
$$
F(a) \circ^{\prime} G(b)=H(a \circ b)
$$
So they proved that $\varphi_{\alpha},$ with $ \alpha\in \Fq^{\sigma},$ of the form
\begin{center}
    $\varphi(a(x)) =  a^{\rho}(\alpha x) $
\end{center}
is an is an  isotopism from  between the semifields $R_{f_1} $ and $R_{f_2}.$  For more details see \href{  https://arxiv.org/abs/1201.2528 }{Semifields from skew polynomial rings} ,  \cite[Section 5]{lavrauw2013semifields}.}
\hlc[green]{Recently they constructed many MRD codes from skew polynomials and $\sigma$-linearized polynomials.}
\textcolor{blue}{
Let  $f\in \Fq[ x;\sigma],$ 
\begin{enumerate}
    \item There exist unique two sided polynomial $f^*\in \Fq^{\sigma}[x^{\mu}],$ called the bound of $f,$ such that $f$ right divides $f^*$ in $ \Fq[ x;\sigma]$ (The existence is guaranteed  over finite fields.) 
    \item If $f$ is two sided then it coincides with its bound $f^*$
    \item For $y=x^\mu$ the bound $ f^*(y)\in \Fq^{\sigma}[y]  $ is commutative polynomial.
    \item If $f$ is irreducible in $ \Fq[x,\sigma]$ then its bound $ f^*(y)$ is irreducible as a polynomial in $\Fq^{\sigma}[y].$ 
    \item Conversely, if $ f^*(y)$ is irreducible as a polynomial in $\Fq^{\sigma}[y]$ of degree $s$ in $y,$ then
    \begin{enumerate}
        \item $ \Fq[x,\sigma]/\langle f^*(x^{\mu}) \rangle  \cong \mathbb M_{\mu}( \mathbb{F}_{q_0^s}) ,$
where $ \mathbb M_{\mu}( \mathbb{F}_{q_0^s})$ is the algebra of square matrices with entries in $ \mathbb{F}_{q_0^s}:=\Fq^{\sigma}.$
        \item Any right divisor of $ f$ in $ \Fq[x,\sigma]$ has degree $s.$
        \item The rank of any  \textbf{irreducible right divisor  $g(x)$}  of $f^*(x^\mu)$ (division in $\Fq[x,\sigma]$) in $\mathbb M_{\mu}( \mathbb{F}_{q_0^s})$ is 
        $$ \rank( g ) = \mu-1.$$
        \item  More general, \cite[Proposition 7]{sheekey2020new}, the rank of any elements  $g(x)$ ($g$ is not necessary a devisor of $f$) in $\mathbb M_{\mu}( \mathbb{F}_{q_0^s})$ is 
        $$\operatorname{rank}(g)=\mu-\frac{1}{s} \deg(\gcrd(g, f^*(x^\mu)).$$
         \end{enumerate}
          \end{enumerate}}


          
   \textcolor{blue}{       
   \begin{enumerate}      
  \item Suppose that $\sigma \in \mathbb{F}_{q^n}$ is of order $n,$ and  $f^*(x^n)$ in $\Fq[x^n]$ is irreducible ($\mathbb{F}_{q^n}^{\sigma}=\Fq).$)
  \begin{enumerate}
  \item They  considered in \cite{sheekey2020new},  $ \mathbb{F}_{q^n}/\Fq  $ and they took  with $\deg(f^*)=s$,and they proved 
$$
R_f:=\frac{\Fq[x,\sigma]}{\langle f^*(x^n)\rangle} \simeq \operatorname{End}_{ \mathbb{F}_{q^s}}\left( \mathbb{F}_{q^n}\right) \simeq M_n\left(\mathbb{F}_{q^s}\right) \simeq \textcolor{red}{ \mathbb{L}_{q^n,q}[x]/\mathbb{L}_{q^n,q}[x]( f^*(x^{q^n}) )}.
$$
 \textcolor{red}{The isomorphic  in red is proved when $ f^*(y)=y-1,$ and so we can prove it for our case for any irreducible polynomials $f^*(y)\in \Fq[y].$}
 \item In \cite[Theorem 7]{sheekey201913}, they proved  under conditions that the set
 $$ S_{n, s, k}(\eta, \rho, f^*):=\left\{a(x)+Rf^*(x^n): \deg(a(x)) \leqslant k s, \ a_{k s}=\eta \rho(a_0)\right\}$$
 is an MRD code with  minimum distance $n-k + 1,$ where $\rho$ an automorphism of $\mathbb{F}_{q^n}.$
 \item   \textcolor{red}{ It mentioned in \cite[Remark 11]{sheekey2020new} that it is possible for different choices of $f^*(y)$ we can produce different equivalent codes.}
\item So when $f^*(y)=y-1$, they got  that  $ M_n(\Fq) \simeq \frac{\mathbb F_{q^n}[x ; \sigma]}{\left(x^n-1\right)}$. So any irreducible factor $g$ of $x^n-1$ should be of degree $1,$ and generates a skew cyclic code $C$ of length $n$ over $\mathbb{F}_{q^n}$ and dimension $n-1.$ When we consider $ C \subseteq M_n\left(\mathbb{F}_{q}\right),$ the  its is of  length $n^2$ and dimension $n(n-1)$ over $\Fq$  and of min rank distance $ n-1.$ Finally we verify that $C$ is and  MRD code.\\
In connection  to what did in section 3, if we consider $\alpha$ a generator of the extension $ \mathbb{F}_{q^n}/\Fq$ and $ \beta=\alpha^{-1}\sigma(\alpha),$ we get the irreducible factorization of $x^n-1$ in $\Fq[x,\sigma]$ is given by
$$ x^n-1=(x-\beta)(x-\sigma(\beta)) \ldots (x-\sigma^{n-1}(\beta))$$
So the irreducible divisors of $x^n-1$ are $(x-\sigma^i(\beta))$ and so the rank of any one of them in $ M_n(\Fq) $  is $ n-1. $ Hence it generates a rank code of length $ n$ over $ \mathbb{F}_{q^n}$ and of dimension $ k=n-1$ and with minimal rank distance $ d_R=n-1.$ We verify that $n-k+1=d   $
and so it  generates an MRD code over $ \mathbb{F}_{q^n}.$  \\
\item With the same conditions.  In our situation (trinomial case) with the polynomial $y^p-y-1 \in \Fq[y]$ irreducible polynomial over $ \mathbb{F}_{q}:= \mathbb{F}_{q^n}^{\sigma},$ (In our first paper we have some conditions on irreducibility of trinomials)
   and then any irreducible divisor of  $ f= x^{pn}-x^{n}-1 $ in $\mathbb{F}_{q^n}[x;\sigma]$ is and   MRD code of length $ pn$ and dimension $n(p-1)$
   over $ \mathbb{F}_{q^p}$ where $p$ is the characteristic  of $\Fq. $ 
     \end{enumerate}
 \item \textcolor{red}{In  relation to the works of J. Sheeky we can interest in the case when $ f^*(y)$ is not irreducible, so we can decompose it into irreducible factors, and then decompose the code in some direct sum..., Maybe this question is far from this project we can do it in another project. Also the study with sum-rank distance should be interesting.  } Just in this context J. sheeky generlized his idea in constructing MRD code to the case of code over $\Fq(t)$ the rational fonctional field \href{https://arxiv.org/pdf/2502.13531}{Quotients of skew polynomial rings new construction}, and the problem of studying non-equivalence here is open according to what people did on equivalence on codes over  $\Fq.$     
    \end{enumerate}   }
    


\end{comment}



 

\vspace{0.5cm}

Next we recall some important facts about algebraic sets and the class of \emph{Wedderburn (\textnormal{or} W-)polynomials}. For more details, we refer the reader to \cite{Leroy2004, Leroy2008}. 

\begin{definition}
\begin{enumerate}
    \item 
    For a skew polynomial $g \in \mathbb{F}_q[x, \sigma]$, the \emph{right vanishing set} of $g$, also called the \emph{set of right roots} of $g$, is defined by  
$
V(g) := \left\{ a \in \mathbb{F}_q : g(a) = 0 \right\}.
$
\item
For any subset $A \subseteq \mathbb{F}_q$, the set $I(A) := \{ g \in \mathbb{F}_q[x, \sigma] : g(A) = 0 \}$ is a left ideal of $\mathbb{F}_q[x, \sigma]$, and the set $A$ is called \emph{$\sigma$-algebraic} if $I(A) \neq 0$. In particular, the set $V(g)$ is $\sigma$-algebraic because $g(V(g)) = 0$. Moreover, any finite set is $\sigma$-algebraic \cite[Section 6]{Lam2000}.
\item If $A$ is $\sigma$-algebraic, then the monic generator of $I(A)$ is called the \emph{minimal polynomial} of $A$, and is denoted by $m_A$. 
\end{enumerate}
\end{definition}

The \emph{minimal polynomial} $m_A$ is the monic least left common multiple of the linear polynomials $\{x - a : a \in A\}$. Thus, it has the form  
$
m_A = (x - a_0)(x - a_1) \cdots (x - a_{n-1}),
$
where each $a_i$ is $\sigma$-conjugate to some element of $A$, and $n = \deg m_A$. Furthermore, any right root of $m_A$ is also $\sigma$-conjugate to some element of $A$.
The \emph{rank} of a set $A \subseteq \mathbb{F}_q$ is defined by $\operatorname{rank}(A) := \deg(m_A)$.

\begin{definition}
\begin{enumerate}
    \item 
A monic polynomial $g \in \mathbb{F}_q[x, \sigma]$ is called a \emph{W-polynomial} if the minimal polynomial of $V(g) = \{ a \in \mathbb{F}_q : g(a) = 0 \}$ is $g$ itself. 
\item 
An element $a \in \mathbb{F}_q$ is called \emph{$P$-dependent} on an algebraic set $A\subseteq \Fq$ if $g(a) = 0$ for every $g \in I(A)$. A $\sigma$-algebraic set $A$ is \emph{$P$-independent} if no element $b \in A$ is $P$-dependent on $A \setminus \{b\}$. 
\end{enumerate}    
Note that the minimal polynomial of a $P$-independent set is a W-polynomial.
\end{definition}

Lastly, we define \emph{exponents} of skew polynomials corresponding to \cite[Definition 3.1.]{Cherchem2016}, which will be crucial to prove the Roos-like bound on the Hamming and rank distance in the next section. (For the definition of exponents of skew polynomials over period rings, see \cite{bouzidi2021exponents}.)

\begin{definition}\cite[Definition 3.1.]{Cherchem2016}
    Let $ f(x) \in \Fq[x;\sigma], $ be a polynomial of degree $n$ with $ f(0)\neq 0 $.  Then there exists  a positive integer $ e $ such that $ f(x) $ right divides $ x^{e}-1.$  The smallest integer $e$ with this property is called the \emph{right exponent}  of $f(x),$ and we denote it by $ \ord_r(f(x))=e.$ 
    If $f(0)=0$, we write $f(x)=x^h g(x)$, where $h \in \mathbb{N}$ and $g \in \mathbb{F}_q[x ; \sigma]$ with $g(0) \neq 0$ are uniquely determined, and define $ \ord_r(f(x)):=\ord_r(g(x))$. The \emph{left exponent} is defined analogously.
\end{definition}




\subsection{Skew polycyclic codes}


We will now turn to polycyclic codes over skew polynomial rings, defining some of the necessary notions for the rest of this paper. We begin with the definition of skew polycyclic codes, based on \cite[Definition 8]{Ouazzou2023}.

  \begin{definition}
  Let $ \vec{a}= (a_0,a_1,\ldots, a_{n-1})\in \Fq^n$   and  $\sigma \in \aut(\Fq)$  be an automorphism of $\Fq.$  A linear code $C\subseteq \Fq^n$ is called a
  \begin{enumerate}
      \item  right \emph{skew  $(\sigma,\vec{a})$-polycyclic   code}   if for each codeword  $(c_0,c_1,\ldots, c_{n-1})$ of $C,$  we have 
$$
\left( 0, \sigma(c_0),\sigma(c_1),\ldots, \sigma(c_{n-2})\right) + \sigma(c_{n-1}) \vec{a}  \in C.
$$
\item left \emph{skew  $(\sigma,\vec{a})$-polycyclic   code}   if for each codeword  $(c_0,c_1,\ldots, c_{n-1})$ of $C,$  we have 
$$
\left(\sigma(c_1),\ldots, \sigma(c_{n-1},0)\right) + \sigma(c_{0}) \vec{a}  \in C.
$$
 \end{enumerate}
    If $\vec{a}=\left(a_{_0},0,\ldots, 0,a_{\ell}, 0, \ldots, 0\right)\in \Fq^{n},$ with $a_0\neq 0 $ and $ a_{\ell}\neq 0$, for some $ 0<\ell <n $, we call a right (resp. left)    $(\sigma,\vec{a})$-polycyclic code with  associated vector $\vec{a}$  a \emph{skew $(\ell,\sigma)$-trinomial code}  of length $n$ over $\Fq$.
 \end{definition}
 In this paper we will focus on right skew polycyclic codes, however, the same result can be proved for the case of left polycyclic codes. For more details about the difference between left and right (skew) polycyclic codes see \cite{Lopez2009,Poly,ou2024linear,Ouazzou2023, Hassan2025Advances}.
 
 Let us recall that a map $T: \Fq^n \to \Fq^n,$ is called a \emph{($\sigma$-)semi-linear transformation} if it is an additive map satisfying $ T(\alpha u)=\sigma( \alpha ) T(u)$ for $\alpha \in \Fq$ and $u\in \Fq^n$. In the particular case $\sigma=\id, \ T$ is a linear transformation. 
 %We can obtain some well-known classes of codes from semi-linear transformations: 
% From the above definition we make the following transformation.
 
 \begin{remark}\label{remark1}
  Let $\vec{a}=(a_0,a_1,\ldots,a_{n-1}) \in \Fq^n$   and  $\sigma \in \aut(\Fq)$  be an automorphism of $\Fq.$ The following follows directly from the definition of (skew) polycyclic codes.
 \begin{enumerate}
 \item  Skew $(\sigma,\vec{a})$-polycyclic   codes are invariant under the $\sigma$-semi-linear transformation $  T_{\sigma,\vec{a}} ,$  called the \emph{$(\sigma,\vec{a})$-polycyclic  shift}, defined by:
 $$  T_{\sigma,\vec{a}} (v_0,v_1,\ldots,v_{n-1} )=  \left( 0, \sigma(v_0), \sigma(v_1), \ldots, \sigma(v_{n-2})\right)+ \sigma(v_{n-1})\vec{a}, $$
 for all $ v=(v_0,v_1,\ldots, v_{n-1}) \in \Fq^n.$
 \item Polycyclic  codes are invariant under the linear transformation $ T_{\vec{a}}$ defined by:
$$ T_{\vec{a}}(v_0,v_1,\ldots,v_{n-1} )=  \left( 0, v_0, v_1, \ldots, v_{n-2}\right)+ v_{n-1}\vec{a}, $$

for all $ v=(v_0,v_1,\ldots, v_{n-1}) \in \Fq^n.$
 \item  When $\sigma= \id,$ we get  $ T_{\vec{a}}= T_{\id,\vec{a}}.$
 \end{enumerate}
 \end{remark}

 Let $R_f:=\Fq[x;\sigma]/ \langle f(x)\rangle,$ with $ \langle f(x)\rangle $ being the left ideal of $\Fq[x,\sigma]$ generated by $g(x).$ In this work, we mainly work with right polycyclic codes, which we henceforth simply refer to as polycyclic codes. Under the usual identification of vectors with polynomials, 
   	\begin{equation}\label{Realization}
  	\begin{array}{rccc}
  	\Phi \ : \ & \ \Fq^n &\longrightarrow & R_f \\ & & & \\
  	  &v=(v_0, v_{_1},\ldots,v_{_{n-1}}) & \longmapsto &  v(x)=\dsum{i=0}{n-1} v_{_i} x^{i}  
  	\end{array}
  	\end{equation} 
each polycyclic code $C$ of length $n$  associated with a vector $\vec{a}$ is seen as a left $ \Fq[x;\sigma]$-submodule (or as a left ideal if $f$ is central)   in  $ R_{f}.$ 


In the following result we  recall more  characterizations of skew $(\sigma,\vec{a})$-polycyclic   codes from \cite[Theorem 4, Theorem 5]{Ouazzou2023}.

 \begin{theorem}\label{T2.3}
 A linear code  $ C\subseteq \Fq^{n} $ is a skew $(\sigma,\vec{a})$-polycyclic   code if and only if $\Phi(C)$  is a left $\Fq[x;\sigma]$-submodule of $  R_f.$ Moreover:
 \begin{enumerate}
 \item There is a monic polynomial of least degree $ g(x)\in \Fq[x;\sigma] $ such that $ g(x) $  right divides $ f(x) $
 and $ \Phi(C)=\langle g(x)\rangle. $
 \item The set $ \{ g(x),x g(x),\ldots,x^{k-1}g(x)\} $ forms a basis of $ C $ and the dimension of $ C $ is $ k=n-\deg(g). $

 \item A generator matrix $ G $ of $ C $ is given by:
 \begin{equation}\label{eq10}
 G= \left(
 \begin{array}{cccccccc}
 g_{_0} &g_{_1} &\cdots&g_{_{n-k}}& 0 &\cdots &\cdots & 0\\
 0 & \sigma(g_{_0}) & \sigma(g_{_1}) &\cdots & \sigma(g_{_{n-k}}) & 0 &\cdots & 0\\
 \vdots &\ddots &\ddots &\ddots & &\ddots & &\vdots\\
 \vdots & &\ddots &\ddots &\ddots & &\ddots &\vdots\\
 0 &\ldots & &0 & \sigma^{k-1}(g_{_0}) & \sigma^{k-1}(g_{_1}) &\ldots & \sigma^{k-1}(g_{_{n-k}})\\
 \end{array}
 \right) 
 \end{equation}
 where $k=n-\deg(g) $ and $ g(x)= \displaystyle\sum_{i=0}^{n-k} g_ix^i. $
\item There is a one-to-one correspondence between  right divisors of $f(x)$ and skew $(\sigma,\vec{a})$-polycyclic   codes of length $n$ over $\Fq$.
\end{enumerate}
 \end{theorem}




From here to the end of this paper, instead of saying ''$\Phi(C)$ is generated by $g(x)$'', we will just say $C$ is generated by $g(x)$ and write $ C= \Fq[x;\sigma] g(x).$ 
     Moreover, for any  $ \sigma \in \aut(\Fq)$ of order $ \mu $ such that $ n = \mu m$, we define the  set
$$
\mathcal{A}_{\sigma} := \left\{ (a_0, 0, \ldots, 0, a_{\mu}, 0, \ldots, 0, a_{2\mu}, \ldots, a_{(m-1)\mu}, 0, \ldots, 0) \in \mathbb{F}_q^n \mid a_i \in \mathbb{F}_{q}^{\sigma}, \ i = 0,  \ldots, (m-1)\mu \right\}.
$$
%\textcolor{purple}{ So we really are only allowed to scale in these specific coordinates? That seems surprising. What happens in the other cases?}
%\hlc[green]{yes exactly, what happned in the case of cyclic codes $f=x^n-1,$ if $n$ is a multiple of $\mu,$ we will in the same situation. Also, for constacyclic codes, $x^n-\lambda,$ $n$ should be a multiple of $\mu$ and $\lambda\in \Fq^{\sigma}.$ In the other situation the polynomial in consideration does not generate a two sided ideal. } \textcolor{purple}{So, it goes back to the same question: do we need the two-sided ideal?}
% \textcolor{blue}{We don't need two-sided ideal. Just if we the two sided ideal the we  speak about of rings isometry  in our study of isometries between skew polycyclic codes and this case we can also speak of interior isomorphism $ \varphi(a(x)):=u^{-1}(x)a(x)u(x)$, section 6 in \cite{Lobillo2025} gave some results for the case of Hamming equivalence between constacylic codes (binomial case), else we will speak about $ \Fq[x,\sigma $-module morphism. But I think there are many things that we can obtain when we consider our study in the ring of $\sigma$-linearized polynomials,$\sigma$ the Frobenius automorphism,  by exploring the linearity property of these polynomials. First idea in the case of $q$-trinomial codes, $(a_0,a_{\ell}) \sim_{(p^n,\ell)} (b_0,b_{\ell})$ if and only if $\gcrd( a_0 x^{p^n} - b_0  x, a_{\ell} x^{p^{n - \ell}} - b_{\ell}x)$ has at least one root in $\mathbb{F}_q^*,$ and we need a deep study in vector space of the root of the this $gcrd$ as we did in the first paper with commutative polynomial. But I think,  we need to focus on this question in another project and give a clear answers.   } \textcolor{purple}{Yes, let's do that in another project. I still don't quite grasp the difference between the two-sided or one-sided ideal. I thought we can always define codes, we just have to be careful with the morphisms (and the field they are defined over) then. But here you are saying we need it for the definition of code.}
% \textcolor{blue}{ Yes exactly, there is difference on the name of codes. If the modulus  polynomial $f$ is central (two- sided) then our codes are seen ideals, else the are left $\Fq[x,\sigma]$-modules. For this reason people who are working in $q$-cyclic codes with $q$-linearized polynomial  over $\mathbb{F}_{q^m}$ the assume that $m $ divides $q^n,$ to be sure that $x^{q^n}-x $ is a central polynomial and they can characterize them as Left ideals in the ring $ \mathcal{L}_{\Fq, \mathbb{F}_{q^m} }/\langle x^{q^n}-x\rangle,$ see for example \cite{martinez2017roots}. The same idea was used  in the case of skew cyclic codes they assume the the order of $\sigma$ divides $n.$ More recently for skew polycyclic codes the modulus  polynomial $f$ should be in the form of $A_{\sigma}.$  in the case of non-central modulus polynomial $f$ there you can look in these references \cite{Bag2025,Cathryn2024,Boucher2014}   }

\begin{theorem}\cite[Theorem 5.9.]{Bag2025}\\
 Let $ \sigma \in \aut(\Fq)$ be of order $ \mu $ such that $ n = \mu m .$ Then a linear code $ C \subseteq \Fq^n$ is a skew polycyclic code induced by $ \vec{a} \in \mathcal{A}_{\sigma} $  if and only if $ C $ is a left ideal of $ \mathbb{F}_q[x; \sigma] / \langle x^n - \vec{a}(x) \rangle .$
\end{theorem}

From the above theorem we can easily deduce the following result regarding skew trinomial codes.

\begin{corollary}
Let $0<\ell<n$ be an integer such that $\ell$ and $n$ are both multiples of $\mu,$ and $a, b$ two non-zeros elements of $ \Fq^{\sigma}.$ Then $C\subseteq \Fq^n$ is a skew $(\ell,\sigma)$-trinomial code associated with $x^n-ax^{\ell}-b$  if and only if $ C $ is a left ideal of $ \mathbb{F}[x;\sigma]/\langle x^n-ax^{\ell}-b \rangle. $
\end{corollary}


Studying the equivalence between two classes of codes with a cyclic structure reduces to analyzing possible isometries, with respect to the desired metric, between the ambient spaces. In the case of skew polycyclic codes, this amounts to studying isometries between 
$
R_{f_1} := \mathbb{F}_q[x, \sigma]/ \langle f_1(x)\rangle $ and $ R_{f_2} := \mathbb{F}_q[x, \sigma]/\langle f_2(x)\rangle,
$
which preserve the parameters of the codes (length, dimension, and minimum distance) as well as their algebraic structure. For skew polycyclic codes, various approaches can be employed depending on the nature of the ambient spaces under consideration. In the following remark we collect some possible approaches, out of which we will use the second in this work.

\begin{remark}
Let $R_{f_i} := \mathbb{F}_q[x, \sigma]/\langle f_i(x)\rangle$ for $i = 1,2$. 
\begin{enumerate}
    \item $R_{f_1}$ and $R_{f_2}$ are non-associative $\mathbb{F}_q^{\sigma}$-algebras (they are associative if $f_i \in \mathcal{Z}(\mathbb{F}_q[x, \sigma])$). We say that $f_1$ and $f_2$ are \emph{equivalent} if there exists an $ \mathbb{F}_{q_0}$-algebra isomorphism between $R_{f_1}$ and $R_{f_2}$ that preserves the desired distance. In the case of skew constacyclic codes, where $x^n - \lambda \in \mathcal{Z}(\mathbb{F}_q[x, \sigma])$ this was done in \cite[Section 6]{Lobillo2017}. For the commutative case $\sigma = \mathrm{id}$, see \cite{Chen2014,Chen2012,Equiv2025}.
    
    \item Another approach is to study equivalence more simply by requiring only the existence of an $\mathbb{F}_q$-vector space isomorphism $\varphi$ between $R_{f_1}$ and $R_{f_2}$ that preserves the desired distance and satisfies $\varphi(a(x)  b(x)) = \varphi(a(x)) \varphi(b(x))$. This approach was used in \cite{Boulanouar2021} to study equivalence between classes of skew constacyclic, cyclic, and negacyclic codes. We will adopt this approach in our setting; see Definition~\ref{FqIso}.
    \item We can also characterize $\sigma$-skew polycyclic codes induced by a polynomial $f(x)$ as left $\mathcal{L}_{q,\sigma}[x]/\langle f(x^{q_0}) \rangle$-modules, or as left ideals if $f \in \mathcal{Z}(\mathcal{L}_{q,\sigma}[x])$. Studying the equivalence between two such classes reduces to studying isometries between algebras with respect to the desired distance, namely between the $ \mathbb{F}_{q_0}$-algebras $\mathcal{L}_{q,\sigma}[x]/\langle f_1(x^{q_0}) \rangle$ and $\mathcal{L}_{q,\sigma}[x]/\langle f_2(x^{q_0}) \rangle$. For more details on the algebraic structure of skew cyclic codes from this point of view, see \cite{gabidulin2009rank,martinez2017roots}.
\end{enumerate}
\end{remark}

%\hlc[green]{I moved the definition of the Shur product here.}
We end this section by recalling the definition of the component-wise product of two vectors that we use throughout this paper. Given two vectors of length $n$, $x = (x_0, x_1, \ldots, x_{n-1})$ and $y = (y_0, y_1, \ldots, y_{n-1})$, the \emph{Schur product} is defined as
$$
x \star y := (x_0 y_0, x_1 y_1, \ldots, x_{n-1} y_{n-1}) .
$$







%%%%%%%%%%%%%%%%%%%%%%%%%%%%%%%%%%%%%%%%%%%%%%%%%%%%%%%%%%%%%%%%%%%%
\section{ Skew Roos-like bound for skew polycyclic codes  }
In \cite{Alfarano2021}, the authors proved a Roos-like bound for the minimum Hamming distance and the minimum rank distance of skew cyclic codes. In their work, they took  $f=x^n-1 \in \Fq[x;\sigma]$, where  $n$ is a multiple of the order of $\sigma$, which means that $x^n-1 $ is a central polynomial.
\begin{comment}
\textcolor{purple}{Do you mean a "two-sided ideal" generated by this polynomial? (same below)} \hlc[green]{yes I mean that $x^n-1$ generates a two sided  ideal}. 
\end{comment}
In this section, we prove this bound for skew \emph{poly}cyclic codes (with respect to the Hamming distance and the rank distance) without assuming that the modulus $f(x)$ is a two-sided polynomial. 

Let us first fix some notation that we use throughout this section.  To decompose a commutative polynomial into a product of linear factors, we usually search for a big extension field which contains all the roots of that polynomial. The same idea is used to decompose skew polynomials, but  in addition we need to define an automorphism over the big field extension, for more details we refer the interested reader to \cite{Lobillo2017}.   Let $\sigma\in \aut(\mathbb{F})$ be an automorphism of $ \mathbb{F},$ of order $\mu.$ 
We say that $\sigma$ has an extension $\theta$ of degree $s$ if there exists a field extension $\mathbb{F} \subseteq \mathbb{L}$ and $\theta \in \aut(\mathbb{L})$ of order $s \mu$ such that $\theta_{\mid \mathbb{F}}=\sigma$ and $ \mathbb{L}^{\theta}=\mathbb{F}^{\sigma}.$ 
Note that if $\mathbb{F}$ is a finite field, extensions of any degree always exist.\footnote{However,  if $\mathbb{F}$ is the field of rational functions over a finite field, they exist when $\mu$  order of $\sigma,$ and $s$ are coprime, see \cite[Examples 2.4 and 2.5]{Lobillo2017}.}   


Let  $f(x)$ be a skew polynomial of degree $n$, with right exponent $e ,$ such that $f(0)\ne 0.$ To construct the  extension $\theta$ of $\sigma,$ we have to consider the following two cases; 
        \begin{itemize}
            \item[•] \textbf{If $e$ is a multiple of the order  $\mu,$} i.e., $e=m\mu$ then the   $\mathbb{F}_{q_0}$-Frobenius automorphism of $ \mathbb{F}_{q_0^e}$ is an   extension of degree $m$ of $\sigma.$ 
          
           \item[•] \textbf{If  $e$ is not a multiple of $\mu,$}  we set  $e':=\lcm(e,\mu)$ and  verify that $ f(x)$ is a right divisor of $ x^{e'}-1,$ since $x^e-1$ right divides $ x^{e'}-1.$  So in this case  the   $\mathbb{F}_{q_0}$-Frobenius automorphism of $ \mathbb{F}_{q_0^{e'}}$ is  an extension of degree 
           $m =\frac{e'}{\mu}$ of $\sigma.$ 
            \end{itemize}
In the sequel, we will assume that $ e$ is a multiple of the order  $\mu,$   and the  $\mathbb{F}_{q_0}$-Frobenius automorphism of $ \mathbb{F}_{q_0^e},$ say $\theta,$ is an extension of degree $m:= \frac{e}{\mu}$ of $\sigma.$  According to  Hilbert's Theorem 90 (see \cite[Chapter VI; Theorem 6.1]{Lang2002}),  $\beta \in \mathbb{F}_{q_0^e} $ is a root of $x^e-1$ if and only if there exist $ \alpha \in \mathbb{F}_{q_0^{e}}^* $ such that  
   $\beta=\theta(\alpha) \alpha^{-1},$ i.e.,   $\beta$ is a $\sigma$-conjugate of $1$.   
   Let us now assume that  $\alpha \in \mathbb{F}_{q_0^e}$ is a \emph{normal} element, i.e., $ B:=\left\{\alpha, \theta(\alpha), \ldots, \theta^{e-1}(\alpha)\right\}$ is a basis of $\mathbb{F}_{q_0^e}$ as an $\mathbb F_{q_0}$ vector space.  For $\beta =\theta(\alpha)\alpha^{-1},$ we can prove that the set
   $$  \{ \beta, \theta(\beta), \ldots,  \theta^{e-1}(\beta)\}$$ forms a $P$-independent set and  a $P$-basis  of  the vanishing set  $V(x^e-1) $  of $x^e-1.$ It follows that the rank of $V(x^e-1) $ equals $e,$ and so we can decompose $x^e-1$ over $ \mathbb{F}_{q_0^{e}} $ as follows:     
      $$
x^e-1=\lclm(\left\{x-\theta^i(\beta) \mid 0 \leq i \leq e-1\right\}).
$$

Now we recall the idea used in \cite{Cathryn2024, Lobillo2025} to decompose the skew polynomial $x^n-a.$

\begin{remark} \label{RmqConsta}
In \cite{Cathryn2024, Lobillo2025} the authors studied 
%, the authors were  interested in Roos bound and BCH bound of
$a$-constacyclic codes of length $n$ over $\Fq,$ which correspond in our case to skew polycyclic codes associated 
    with the polynomial $f(x)=x^n-a$. 
    Their  idea was to find a decomposition of $x^n-a$ in a big extension  $ \mathbb{F}_{q_0^{\mu t}},$ where $ n=\mu t$ under the  condition that $ a=N_n^{\sigma}(\gamma)$ for some $\gamma\in \Fq.$ This latter condition ensures that $x^n-a$ has at least a root in $\Fq. $ Let $\alpha$ be a normal element in $\mathbb{F}_{q^t}$ and  $ \beta=\theta(\alpha) \alpha^{-1},$ with $\theta \in \aut(\mathbb{F}_{q^t})$ an extension of $\sigma$ of degree $t.$ Then  $\{ \gamma \beta, \gamma \theta(\beta),\ldots,  \gamma \theta^{t-1}(\beta )  \}$ forms a $P$-independent set and  a $P$-basis  of  the vanishing set  $V(x^n-a) $  of $x^n-a.$ It follows that the rank of $V(x^e-1) $ equals $n,$ and so they  decomposed $x^n-a$ over $ \mathbb{F}_{q^{t}} $ as follows:     
      $$
x^n-a=\lcm(\left\{x- \gamma \theta^i(\beta) \mid 0 \leq i \leq n-1\right\}).
$$
\end{remark}

Keeping the same notation as above,  we give the definition of the $\beta$-defining set of a given skew polycyclic code associated with a polynomial  $f\in \Fq[x;\sigma].$
\begin{definition} 
Let $g \in \Fq[x;\sigma]$ be  a right divisor of $f(x),$ denoted $ g(x) ~|_r~ f(x),$ and $ C= \Fq[x,\sigma] g(x)$ and $\widehat{C}= \mathbb{F}_{q_0^e}[x,\sigma] g(x) .$
The \emph{$\beta$-defining set} of $g$ is
$$
T_\beta(g)=\left\{ 0 \leq i \leq e-1  ~:~ x-\theta^i(\beta) ~~|_r~~ g \right\}.
$$
\end{definition}
In particular, $\lclm(\left\{x-\theta^i(\beta) ~ \mid~ i \in T_\beta(g) \right\}) ~ ~|_r~~ g$ in $ \mathbb{F}_{q_0^e}[x;\sigma].$ In the sequel we will prove that this division is in  $ \mathbb{F}_{q}[x;\sigma].$


The following lemma will be useful to prove Theorem \ref{thm:Roos} and Theorem \ref{Rank_bound}. 
\begin{lemma}[Lemma 12, \cite{Alfarano2021}]\label{Circulant_lemma}
Let $E/K$ be an extension field of finite degree $n$. Let $\alpha_1, \ldots, \alpha_{t+r} \in E$ be linearly independent elements over $K$, and  $\left\{k_0, \ldots, k_r\right\}$ $\subseteq\{0, \ldots, n-1\}$ be such that $k_r-k_0 \leq t+r-1$ and $k_{j-1}<k_j$ for $1 \leq j \leq r$. Let

$$
A_0=\left(\begin{array}{cccc}
\theta^{k_0}\left(\alpha_1\right) & \theta^{k_1}\left(\alpha_1\right) & \cdots & \theta^{k_r}\left(\alpha_1\right) \\
\vdots & \vdots & \ddots & \vdots \\
\theta^{k_0}\left(\alpha_{t+r}\right) & \theta^{k_1}\left(\alpha_{t+r}\right) & \cdots & \theta^{k_r}\left(\alpha_{t+r}\right)
\end{array}\right) .
$$
Let $b \in\{0, \ldots, n-1\}$ such that $(b, n)=1$ and
$$
A_i=\left(A_0\left|\theta^b\left(A_0\right)\right| \ldots \mid \theta^{b i}\left(A_0\right)\right)
$$
for $0 \leq i \leq r$. Then $\rank\left(A_{t-1}\right)=t+r$.
\end{lemma}


\subsection{Hamming distance}

The following theorem gives a skew Roos-like bound on Hamming distance for skew polycyclic codes.
    \begin{theorem}[Skew Roos-like bound for the Hamming distance]\label{thm:Roos}
    Let $g \in \Fq[x;\sigma]$ be such that $ g(x) ~|_r~ f(x)$, let $ C:= \Fq[x;\sigma] g(x)$ be a skew polycyclic code generated by $g(x)$ and $\widehat{C}:= \mathbb{F}_{q_0^e}[x;\sigma] g(x) $ be its extension code over $ \mathbb{F}_{q_0^e}.$
    Let $\beta:= \theta(\alpha)\alpha^{-1}$ where $\alpha$ is a normal element of $ \mathbb{F}_{q_0^e}$. 
    Suppose that there exist integers $ a, b, \delta, r, k_0, \dots , k_r$ such that $\gcd( b,e)=1,\  k_j<k_{j+1}\ \text { for } 0 \leq j \leq r-1, \ k_r-k_0 \leq \delta+r-2, $ and  
    $$ (x-{\theta^{a+ib+k_j}(\beta)}) \mid_r g(x)  ~\text{for}~ i=0,\ldots,{\delta-2}, j=0,\ldots, r-1,$$ 
    i.e., $$ \left\{ a+ib+k_j   ~:~ i=0,\ldots,{\delta-2}, j=0,\ldots, r-1 \right\} \subseteq T_{\beta}(g). $$ Then  the minimum Hamming distance $d_H(C)$ of  $C$ satisfies     $d_H(C)\geq   d_H(\widehat{C}) \geq {\delta}+{r}. $
\end{theorem}

\begin{proof}
Let $c$ be a codeword of weight $ w <\delta+r, $ then 
				$$c(x) =\dsum{h=0}{w-1} c_{h}x^{l_h}\;\mbox{ for some }\;\{ l_{0},\ldots,l_{w-1}\}\subseteq\{ 0,\ldots,n-1\}. $$
As $ (x-{\theta^{a+ib+k_j}(\beta)}) \mid_r g(x)$, we have 
				$$ 0=c(\beta^{a+ib+k_j}) = \dsum{h=0}{w-1} c_{h} N_{l_h}(\theta^{a+ib+k_j}(\beta)) = \theta^{a+ib+k_j}(\alpha)^{-1}\dsum{h=0}{w-1} c_{h} \theta^{a+ib+k_j+l_h}(\alpha)$$
It follows that  $ c=(c_1,\ldots,c_w)$ is in the kernel of the matrix $B$ where 

            $$ B= \left[A, \theta^b(A),\theta^{2b}(A),\ldots, \theta^{b(\delta-2)}(A) \right], $$ 
            and  
            $$ 
            A=\left( \theta^{k_j+l_h}(\theta^a(\beta)) \right)_{\substack{0 \leq h \leq w-1 \\ 0 \leq j \leq r-1}}.$$
According to Lemma \ref{Circulant_lemma} , we obtain  that $\rank(B)=w$, and so, $c=0.$  
Therefore, $d_H(C) \geq d_H(\widehat{C}) \geq \delta+r.$ 
\end{proof} 

%In the following remark, we provide links to existing results on the Roos-like bound in the cases of cyclic and constacyclic codes \cite{Alfarano2021, Lobillo2017, Lobillo2025}. They can be recovered as special cases of our more general result.

Then general bound in Theorem \ref{thm:Roos} includes the known bounds for special cases like cyclic and constacyclic codes. 
%\begin{remark}\label{Rmk_Gen}
 For this let  $f\in \Fq[x;\sigma]$ be a skew polynomial of degree $n$ and of right exponent $e ,$ with $f(0)\ne 0.$ 
\begin{enumerate}
    \item If $f=x^n-1$, with $n$ being a multiple of $ \mu ,$  we find the case of skew cyclic codes \cite{Alfarano2021,Lobillo2017}.
      \item If $f=x^n-\lambda $, with $n$ being a multiple of $ \mu $ and $\lambda \in \Fq^{\sigma},$  we find the case of skew constacyclic codes \cite{Lobillo2025}.
      \item  If $f=x^n-\lambda $, with $\lambda= N_n^{\sigma}(\gamma) $ for some $\gamma \in \Fq, $  we find the case of skew constacyclic  codes  \cite{Cathryn2024}.
\end{enumerate}
%\end{remark}


The above bound includes particular cases, namely the skew BCH-like bound and the Hartmann–\,Tzeng-like bound (in the Hamming distance), as summarized in the following: 
%These bounds are proved in the cases of skew cyclic codes and constacyclic codes in \cite{Alfarano2021, Lobillo2017, Lobillo2025, Cathryn2024}.

\begin{corollary}[\cite{Alfarano2021, Lobillo2017, Cathryn2024, Lobillo2025}]\label{CorBch}
Let $g \in \Fq[x,\sigma]$ be such that $ g(x) ~|_r~ f(x),$ let $ C= \Fq[x,\sigma] g(x)$ be a skew polycyclic code generated by $g(x)$ and let $\widehat{C}= \mathbb{F}_{q_0^e}[x,\sigma] g(x) $ be its extension code over $ \mathbb{F}_{q_0^e}.$ Let $\beta= \theta(\alpha)\alpha^{-1}$ where $\alpha$ is a normal element of $ \mathbb{F}_{q_0^e}$. 
\begin{enumerate}
    \item (Skew BCH-like bound) If there exist integers $a,b$ such that  
    $$ \{(x-\theta^{a+jb}(\beta)) \mid_r g(x) ~:~ j=0,\ldots, \delta-2\},$$ 
 i.e., $$ \left\{ a+jb   ~:~ j=0,\ldots,{\delta-2},  \right\} \subseteq T_{\beta}(g), $$
 with  $\gcd(e, b)=1$,  then $C$ has  minimum Hamming distance  $ d_H(C)\geq   d_H(\widehat{C}) \geq {\delta}.$
 
 \item  (Skew Hartmann-Tzeng-like bound)      If there exist integers $a,b$ such that 
 $$ \{ (x-{\theta^{a+ib+jc}(\beta)}) \mid_r g(x)  ~:~ i=0,\ldots,{\delta-2}, j=0,\ldots, {r}\},$$ 
  i.e., $$ \left\{ a+ib+jc  ~:~ i=0,\ldots,{\delta-2}, j=0,\ldots, {r} \right\} \subseteq T_{\beta}(g), $$
	with   $\gcd(e, b)=1 $ and $\gcd(e, c)<\delta$, then  the minimum Hamming distance $d_H(C)$ of  $C$ satisfies   $d_H(C)\geq   d_H(\widehat{C})  \geq {\delta}+ {r}. $ 
\end{enumerate}
\end{corollary}




\subsection{Rank distance}

We now turn to the rank distance and prove the analog bound. In this section  we consider the rank distance $ \Fq/ \mathbb{F}_{q_0}$ where $ \mathbb{F}_{q_0} =\Fq^{\sigma}$ is the fixed subfield of $ \Fq$ by $\sigma$.

\begin{theorem}[Skew Roos-like bound for the rank distance]\label{Rank_bound}
Let $g \in \Fq[x,\sigma]$ be such that $ g(x) ~|_r~ f(x),$ let $ C:= \Fq[x,\sigma] g(x)$ be a skew polycyclic code generated by $g(x)$ and let $\widehat{C}:= \mathbb{F}_{q_0^e}[x,\sigma] g(x) $ be its extension code over $ \mathbb{F}_{q_0^e}.$ Let $\beta: = \theta(\alpha)\alpha^{-1}$ where $\alpha$ is a normal element of $ \mathbb{F}_{q_0^e}$. 
    Suppose there exists integers $ a, b, \delta, r, k_0, \dots , k_r$ such that $\gcd( b,e)=1,\  k_j<k_{j+1}\ \text { for } 0 \leq j \leq r-1, \ k_r-k_0 \leq \delta+r-2, $ and  
    $$\{ (x-{\theta^{a+ib+k_j}(\beta)}) \mid_r g(x)  ~\text{ for }~ i=0,\ldots,{\delta-2}, j=0,\ldots, r-1.\}$$ 
    Then  the minimum rank distance of $C$ satisfies     $d_R(C)\geq   d_R(\widehat{C}) \geq {\delta}+ {r}. $
\end{theorem}
\begin{proof}
According to Proposition \ref{Prop_Hamming_rank}, it suffices to prove $d_H(C M) \geq \delta+r$ for all $M \in \GL(\mathbb{F}_{q_0})$. Let $M \in \GL(\mathbb{F}_{q_0})$, and assume that there is a codeword $c \in C M^{-1}$ of Hamming weight $w < \delta+r$, i.e., $\widehat{c} := cM \in C$. Let $\widehat{c}(x) := \dsum{h=0}{n-1} \widehat{c}_{h}x^{h}\in C $. For  $i = 0, \ldots, {\delta - 2}, j = 0, \ldots, r,$ we have  that  $x - {\theta^{a+ib+k_j}(\beta)}~|_r~ g(x).$ Then it right divides $\widehat{c}(x)$, and so
$$
0 = \widehat{c}(\beta^{a+ib+k_j}) = \dsum{h=0}{n-1} \widehat{c}_{h} N_{h}(\theta^{a+ib+k_j}(\beta)) = \theta^{a+ib+k_j}(\alpha)^{-1} \dsum{h=0}{n-1} \widehat{c}_{h} \theta^{a+ib+k_j+h}(\alpha) = 0.
$$
It follows that $c = \widehat{c}M^{-1}$ is in the kernel of the matrix $MB$, where
$$
B = \left[A, \theta^b(A), \theta^{2b}(A), \ldots, \theta^{b(\delta-2)}(A)\right],
$$
and 
$$
A = \left( \theta^{k_j+h}(\theta^a(\beta)) \right)_{\substack{0 \leq h \leq n-1 \\ 0 \leq j \leq r}}.
$$
Now let $S := \{l_0, \ldots, l_{w-1}\} \subseteq \{0, \ldots, n-1\}$ be the set of the non-zero components of $c$, since we assumed that $w_H(c) = w$. Set $\tilde{c} = (c_{l_0}, \ldots, c_{l_{w-1}}) \in \Fq^{w}$ and let $M_S$ be the rows of $M$ indexed by $S$.  
Then $cM = \tilde{c} M_{S}$, and so $\tilde{c}$ is in the kernel of $M_{S} B$. As $M_S \in \mathbb{M}_{w\times n}(\mathbb{F}_{q_0})$, we obtain that
$$
M_{S} B = \left[M_{S}A, \theta^b(M_{S} A), \theta^{2b}(M_{S}A), \ldots, \theta^{b(\delta-2)}(M_{S}A)\right].
$$
Next, we observe that the matrix $M_SA$ is of the form
$$
M_SA = \left( \theta^{k_j}(\theta^a(\beta_h)) \right)_{\substack{0 \leq h \leq w-1 \\ 0 \leq j \leq r}},
$$
where $\beta_h := M_{\{l_h\}} \star\alpha^{[\theta]},$ with $\star$ means component-wise product of the vectors $ M_{\{l_h\}}$ and $ \alpha^{[\theta]}, $ and  $\alpha^{[\theta]} := (\alpha, \theta(\alpha), \ldots, \theta^{n-1}(\alpha)) \in \mathbb{F}_{q_0^e}^n, $ are linearly independent over $\mathbb{F}_{q_0}$. In fact,
let $t_0, \ldots, t_{w-1} \in \mathbb{F}_{q_0}$ be such that $t_0 \beta_0 + t_1 \beta_1 + \ldots + t_{w-1} \beta_{w-1} = 0$,
then
$$
0 = \dsum{h=0}{w-1} t_h \beta_h 
= \dsum{h=0}{w-1} t_h \left(\dsum{i=0}{n-1} m_{l_h, i} \theta^{i}(\alpha) \right) 
= \dsum{h=0}{w-1} \dsum{i=0}{n-1} t_h m_{l_h, i} \theta^{i}(\alpha) 
= \dsum{i=0}{n-1} \left(\dsum{h=0}{w-1} t_h m_{l_h, i} \right) \theta^{i}(\alpha).
$$
Since $\left\{\alpha, \theta(\alpha), \ldots, \theta^{e-1}(\alpha)\right\}$ is a basis of $\mathbb{F}_{q_0^e}/\mathbb{F}_{q_0}$ and $n \leq e$,  $\alpha, \theta(\alpha), \ldots, \theta^{n-1}(\alpha)
$ are linearly independent over $\mathbb{F}_{q_0}$. It follows that $\dsum{h=0}{w-1} t_h m_{l_h, i} = 0$ for all $i = 0, \ldots, n-1$.  This forces $\left(t_0, \ldots, t_{w-1}\right) M_S = 0$. Now, since $M$ is invertible, we know $\rank(M_S) = w$, so $t_h = 0$ for all $0 \leq h \leq w-1$. Thus, $\beta_0, \ldots, \beta_{w-1}$ are linearly independent over $\mathbb{F}_{q_0}$. Therefore, by Lemma \ref{Circulant_lemma}, we have $\rank(M_S B) = w$. This forces $\tilde{c} = 0$ and hence $c = 0$. Therefore, $d_R(C) \geq d_R(\widehat{C}) \geq \delta + r$. 
\end{proof}



\begin{remark}\label{Rmq_gamma}
    In \cite[Example 5.1.5]{Cathryn2024}, K. Hechtel gave an example where the Roos bound on the rank metric for skew $ a $-constacyclic codes does not hold for $ \gamma \in (\mathbb{F}_{q_0} \setminus \mathbb{F}_{q}) $ such that $ a = N_n^{\sigma}(\gamma) .$ Following our approach, which uses the decomposition of $ x^e - 1 $ instead of $ x^n - a, $ where $e$ is the right order of $x^n-a$ (see Remark \ref{RmqConsta}), we obtain that the bound holds without any restriction. The difference between our approach and hers lies in the proof of the above theorem. Specifically, in the proof that  $ \beta_0, \beta_1, \ldots, \beta_{w-1} $ are linearly independent. However, in her argument, linear independence holds only if $ \gamma \in \mathbb{F}_{q_0} ;$ for more details, see \cite[ Theorem 5.1.4]{Cathryn2024}.

\end{remark}




As above we can recover the particular cases of the skew BCH-like bound and the skew Hartmann-Tzeng-like bound on the rank metric for skew polycyclic codes. 
\begin{comment}
\textcolor{purple}{References to the rank bounds?}
\hlc[green]{yes these references}  \cite{Alfarano2021,  Lobillo2025, Cathryn2024} \hlc[green]{ are for the  rank metric. They studied  bounds on both distances rank and Hamming.}
\end{comment}
\begin{corollary} [\cite{Alfarano2021, Cathryn2024,  Lobillo2025}] \label{CorRank_Bch}
Let $C$ be a skew polycyclic code associated with $f,$ of length $n$ over $\Fq$ generated by $  g(x) \in \mathbb{F}_{q}[x;\sigma]$  and let $ \widehat{C} $ be its extension code of the same length and generator over $ \mathbb{F}_{q_0^e}$ i.e., $\widehat{C}= \mathbb{F}_{q_0^e}[x;\theta] g$.  Let $\beta= \theta(\alpha)\alpha^{-1}$ where $\alpha$ is a normal element of $ \mathbb{F}_{q_0^e}$.
\begin{enumerate}
    \item (Skew BCH-like bound) If   $$ \{(x-\theta^{a+jb}(\beta)) \mid_r g(x) ~\text{ for }~ j=0,\ldots, \delta-2\},$$
 such that  $\gcd(e, b)=1$,  then $ d_R(C)\geq   d_R(\widehat{C}) \geq {\delta}.$ 
 
 \item  (Skew Hartmann-Tzeng-like bound)    If   $$ \{ (x-{\theta^{a+ib+jc}(\beta)}) \mid_r g(x)  ~\text{ for }~ i=0,\ldots,{\delta-2}, j=0,\ldots, {r}\},$$ 
	where   $\gcd(e, b)=1 $ and $\gcd(e, c)<\delta$, 
    then  the minimum rank distance $d_R(C)$ of  $C$ satisfies $d_R(C)\geq   d_R(\widehat{C})  \geq {\delta}+ {r}. $ 
\end{enumerate}
\end{corollary}

\subsection{Constructions of skew polycyclic codes with designed distance. }

To construct classical cyclic (constacyclic, polycyclic) codes of length $n$ over $\mathbb{F}_q$, one studies the factorization of $x^n - 1$ over $\mathbb{F}_q$, obtained by considering its roots in an extension field of $\mathbb{F}_q$. These roots form a multiplicative group of order $n$, generated by a primitive element $\zeta$. The irreducible factors of $x^n - 1$ correspond to cyclotomic cosets that are $\langle q \rangle$-orbits on $\mathbb{Z}_n$. Each coset $C_i$ defines a minimal polynomial  $
M_{\zeta^i}(x) = \prod_{k \in C_i} (x - \zeta^k), $ which is irreducible over $\mathbb{F}_q$. 


In the case of skew codes, an alternative to the classical concept of cyclotomic cosets is the notion of  \emph{$\mu$-closed sets}, introduced by Gómez-Torrecillas et al.\ in \cite{Lobillo2017} and further developed in \cite{Alfarano2021, Cathryn2024, Lobillo2025}. We adapt this concept to construct skew polycyclic codes in our setting. Let $\mathbb{Z}_e = \{0,1,\ldots,e-1\}$ be the cyclic group of order $e$, and suppose $e = m\mu$. Then,
$$
\mu \mathbb{Z}_e = \{0, \mu, \ldots, (m-1)\mu\}
$$
is a subgroup of order $m$, and the quotient group $\mathbb{Z}_e / \mu \mathbb{Z}_e$ partitions $\mathbb{Z}_e$ as
$$
\mathbb{Z}_e = \bigcup_{i=0}^{m-1} (i + \mu \mathbb{Z}_e).
$$
Let $\alpha \in \mathbb{F}_{q_0^e}$ be such that $\{\alpha, \theta(\alpha), \ldots, \theta^{e-1}(\alpha)\}$ forms a normal basis of $\mathbb{F}_{q_0^e}/\mathbb{F}_{q_0}$, and set $\beta := \alpha^{-1} \theta(\alpha)$. Then, as previously discussed, the polynomial $x^e - 1$ factors over $\mathbb{F}_{q_0^e}$ as:
$$
x^e - 1 = \lclm\left\{ (x - \beta), (x - \theta(\beta)), \ldots, (x - \theta^{e-1}(\beta)) \right\}.
$$
Grouping roots according to $\mu$-closed sets, we obtain
$$
x^e - 1 = \lclm\left\{ M_{\beta}(x), M_{\theta(\beta)}(x), \ldots, M_{\theta^{\mu - 1}(\beta)}(x) \right\},
$$
where for each $i \in \{0, \ldots, \mu - 1\}$,
$$
M_{\theta^i(\beta)}(x) := \lclm \left\{ x - \theta^{i + j\mu}(\beta) \mid 0 \leq j < m \right\}.
$$

We now recall the following result, which will be used in the construction of skew polycyclic codes with some designed distance.

\begin{lemma}(\cite[Lemma 5]{Lobillo2025}).\label{Skew_Roots}
Let $ \gamma  \in \mathbb{F}_{q_0^e} $ and $g(x)\in \mathbb{F}_{q}[x;\sigma].$ If $ \gamma   $ is a right root of $g,$ then $ \theta^{\mu}(\gamma)$ is also a right root of $g.$ 
\end{lemma}

\begin{proposition}\label{mu-closed}
Let $f\in \Fq[x;\sigma]$ be of order $e$ such that $e=m \mu,$ and  let $\left\{\alpha, \theta(\alpha), \ldots, \theta^{e-1}(\alpha)\right\}$ be a normal basis of $ \mathbb{F}_{q_0^e}/\mathbb{F}_{q_0}$. Define $\beta:=\theta(\alpha) \alpha^{-1}.$ 
\begin{enumerate}
    \item For each $ i=0,\ldots, \mu-1,$   the polynomial $ M_{\theta^i(\beta)}(x)\in \Fq[x;\sigma] $ and $ \deg(M_{\theta^i(\beta)}(x) )= m.$
   Moreover it is  the minimal polynomial of $ T_{\beta}^{(i)}:= \{ \theta^i(\alpha), \theta^{i+\mu}(\alpha), \ldots, \theta^{i+\mu(m-1)}(\alpha)\}$ over $\Fq.$ 
    \item If $ \theta^{i}(\beta)$ is a right root of $g\in \Fq[x,\sigma]$ then $ M_{\theta^i(\beta)}(x) $ right divides $ g(x)$ in $\Fq[x,\sigma].$
%    \item There are distinct integers $i_1,i_2,\ldots, i_r\in \{0,1,\ldots, \mu-1 \}$ such that  $$ f(x)=\lclm(M_{\theta^{i_1}(\beta)}(x),M_{\theta^{i_2}(\beta)}(x),\ldots, M_{\theta^{i_r}(\beta)}(x)).$$
    \end{enumerate}
\end{proposition}
\begin{proof}
\begin{enumerate}
    \item  We extend $\theta $ to $\mathbb{F}_{q_0}[x;\theta]$ in the natural way i.e., $ \theta\left( \dsum{i=0}{d-1} a_i x^i\right)= \dsum{i=0}{d-1} \theta(a_i) x^i.$ As the fixed field of $ \mathbb{F}_{q_0^e}$ by $ \theta^{\mu}$ is $ \Fq,$  it suffices  to prove that $ \theta^{\mu}( M_{\theta^i(\beta)}(x))=  M_{\theta^i(\beta)}(x).$ For each $ i=0,\ldots, \mu-1,$  we have 
    $$
    \begin{array}{rl}
        \theta^{\mu}\left( M_{\theta^i(\beta)}(x)\right)
        &= \theta^{\mu}\left( \lclm\left((x-\theta^{i}(\beta)), (x-\theta^{i+\mu}(\beta)), \ldots, (x-\theta^{i+(m-1)\mu}(\beta)) \right)  \right) \\
         & = \lclm\left((x-\theta^{i+\mu}(\beta)), (x-\theta^{i+2 \mu}(\beta)), \ldots, (x-\theta^{i+m\mu}(\beta)) \right)  \\
         &= M_{\theta^i(\beta)}(x), \ \text{since the order of $\theta$ is $e=m\mu.$} \\  
    \end{array}
     $$
     Then $ M_{\theta^i(\beta)}(x) \in \Fq[x;\sigma].$ As the set $ T_{\beta}^{(i)} $ is a subset of the $P$-independent set $ \left\{\alpha, \theta(\alpha), \ldots, \theta^{e-1}(\alpha)\right\}$, $\deg(M_{\theta^i(\beta)}(x))=m.$ Therefore, it is the minimal polynomial of $ T_{\beta}^{(i)}.$ 
     \item Let  $ \theta^{i}(\beta)$ be a right root of $g\in \Fq[x;\sigma]$ then according to  Lemma \ref{Skew_Roots}, $ \theta^{\mu}(\theta^i(\beta))=\theta^{i+\mu}(\beta)$ is also a right root of $g.$ Consequently, for any $j=0,\ldots,m-1,$ $\theta^{j\mu}(\theta^i(\beta))= \theta^{i+j\mu}(\beta) $ is a root of $g.$ Hence $ M_{\theta^i(\beta)}(x) $ right divides $ g(x)$ in $\Fq[x;\sigma],$
     since it is the minimal polynomial of $\{ \theta^i(\beta),  \theta^{\mu}(\theta^i(\beta)),  \ldots, \theta^{(m-1)\mu}(\theta^i(\beta)) \}. $
     % \item Follows from that fact that $f(x)$ right divides $ x^e-1$ and $$ x^e - 1 =\lclm\left(\left\{ M_{\beta}(x), M_{\theta(\beta)}(x), \ldots, M_{\theta^{\mu - 1}(\beta)}(x) \right\}\right).$$
    \end{enumerate}

\end{proof}
Now we recall the definition of $ \mu$-closed set and its representative  elements. 
\begin{definition}
Let $ T\subseteq \mathbb{Z}_e$ be a non-zero set. 
\begin{enumerate}
    \item We say that   $ T\subset \mathbb{Z}_e$ is \emph{$\mu$-closed} if   $i \in T$ implies $i+\mu \in T$  for every  $i \in T$. 
    \item If $T$ is $\mu$-closed  and $\{i_1,i_2,\ldots, i_l\} \subseteq T$  are integers such that
    $$T=T^{i_1} \cup \cdots \cup T^{i_l} , \text{ and } T^{i_k}\cap  T^{i_j}= \{0\} $$
    then $ S_{T}:= \{i_1,i_2,\ldots, i_l\}$ is  called the \emph{$\mu$-representative set} of $T.$ 
   \end{enumerate} 
\end{definition}
 Following  \cite[Proposition 17]{Alfarano2021}, we can easily prove the following result.

\begin{proposition}
 Let $g \in \mathbb{F}_q[x; \sigma]$ be such that $g ~|_r~ f(x)$. Let $\alpha \in \mathbb{F}_{q_0^e}$ be such that $\left\{ \alpha, \theta(\alpha), \ldots, \theta^{e-1}(\alpha) \right\}$ forms a normal basis of $\mathbb{F}_{q_0^e}/\mathbb{F}_{q_0}$. Set  $\beta := \theta(\alpha) \alpha^{-1}$, and let $ T_\beta(g) := \left\{ i \in \mathbb{Z}_e : x - \theta^i(\beta) \mid_r g \right\}$ the $\beta$-defining set of $g.$
 \begin{enumerate}
     \item If $ g \in \Fq[x;\sigma]$ then $ T_\beta(g) $ is $\mu$-closed.
     \item Conversely, if $ T_{\beta}(g) \subseteq \mathbb Z_e$ is $\mu$-closed, then $ 
    g':= \lclm( \{x-\theta^i(\beta) ~:~ i \in T_{\beta}\}) \in \Fq[x,\sigma].$
 \end{enumerate}
\end{proposition}
\begin{proof}
\begin{enumerate}
     \item  Let $  g \in \Fq[x;\sigma] $ and  $j\in T_\beta(g) ,$  then  $g(\theta^j(\beta))=0$, i.e.,  $ \gamma:=\theta^{j}(\beta)$ is a right root of $g.$  By Lemma \ref{Skew_Roots},  $ \theta^{\mu}(\gamma)=\theta^{j+\mu}(\beta) $   is also a right root of $g,$ and so $ j+\mu \in  T_{\beta}(g). $   Hence  $ T_\beta(g) $ is $\mu$-closed.
     \item Suppose that  $ T_{\beta} \subseteq \mathbb Z_e$  is $\mu$-closed.   Let $ \{i_1,i_2,\ldots, i_l\}$  be  the representative set of $T_{\beta}(g),$ and $ M_{\theta^{i_j}(\beta)}(x):= \lclm\{(x-\theta^{i_1}(\beta)), (x-\theta^{i_1+\mu}(\beta)),\ldots, (x-\theta^{i_1+(m-1)\mu}(\beta))$ for $ j=1,\ldots,r.$ It follows that $g'=\lclm\{ M_{\theta^{i_1}(\beta)}(x), M_{\theta^{i_2}(\beta)}(x),\ldots, M_{\theta^{i_r}(\beta)}(x) \}. $ So by Proposition \ref{mu-closed},  we deduce that $g'\in\Fq[x;\sigma].$  
 \end{enumerate}
\end{proof}




\begin{comment}
\textcolor{purple}{Here we need some more text, why we list this last proposition (for example), and what we have shown overall in this subsection.} \hlc[green]{I edited some typos in this part please look at it again.} \textcolor{purple}{Good, but you should either write $g(x)$ and $f(x)$ everywhere, or $g$ and $f$.}
\end{comment}
We know that if $g \in \mathbb{F}_{q}[x ; \sigma]$ is a right  divisor of $f$,  then it is in fact the minimal polynomial for its $\beta$-defining set $T_{\beta}(g)$. In the following result, we will prove that the size of the $\mu$-representative set of $T_{\beta}(g)$ has implications for when  $ C := \mathbb{F}_q[x;\sigma] g(x) $  is an MRD code. The proof of the following proposition is analogous to \cite[Proposition 26]{Alfarano2021} (the case of skew cyclic codes) and \cite[Proposition 5.2.4]{Cathryn2024} (the case of skew constacyclic codes), so we omit it.

\begin{proposition}\label{SKEW_MRD}
 Let $g \in \mathbb{F}_q[x; \sigma]$ be such that $g ~|_r~ f(x)$. 
 Suppose that the defining set $T_\beta(g)$ satisfies a skew Roos bound as in Theorem \ref{Rank_bound} for some $\delta \geq 2$ and $r \geq 0$. Then the minimum rank distance of the code $C=\mathbb{F}_{q}[x;\sigma] g$ satisfies $$\delta+r \leq d_R(C) \leq \left| S_{T_{\beta}(g)} \right|+1$$ where $  S_{T_{\beta}(g)} \subset \mathbb{Z}_{\mu} $ is the $\mu$-representative set of $T_{\beta}(g).$ 
 In particular, if $\left|  S_{T_{\beta}(g)} \right|=\delta+r-1$, then $ C$ is an MRD code.
\end{proposition}

In the following remark, we discuss how to construct an MRD code from the class of skew polycyclic codes using Proposition \ref{SKEW_MRD}, and present some existing results in the cases of skew cyclic and constacyclic codes.


Keeping the same notations as in Theorem \ref{Rank_bound}, we define
$T_{\beta}(f) := \{ i \in \mathbb{Z}_e \mid (x - \theta^i(\beta)) \mid_r f(x) \}, $ i.e., the set of exponents corresponding to the roots of $f$ in $\mathbb{F}_{q_0^e}[x;\theta]$.
According to Proposition \ref{SKEW_MRD}, we observe that the construction of MRD codes corresponds to finding a $\beta$-defining set $T \subset \mathbb{Z}_e \cap T_{\beta}(f)$ such that
$$|S_{\overline{T}}| = \delta + r - 1. $$ Here, $\overline{T} \subset \mathbb{Z}_e$ is the $\mu$-closed set generated by $T$, and $S_{\overline{T}} \subset \mathbb{Z}_{\mu}$ is its $\mu$-representative set. It follows that $$ g:= \lclm_{ \{ i\in S_{\overline{T}} \}}\left( M_{\theta^i(\beta)}(x)\right)  \in \Fq[x, \sigma] $$
is a right divisor of $f. $ Hence it is a generator of  an $[n, n-m(\delta + r - 1), \delta + r - 1]_q $ MRD code. 
\begin{remark}
In \cite[Theorem 5.2.7]{Cathryn2024}, in the case of constacyclic codes, it was proved that if the above equality holds, that is, $ S_{\overline{T}} = \delta + r - 1, $
then $S_{\overline{T}}$ must be an arithmetic progression with common difference $b$, for any value of $\mu.$ This means that
$$ S_{\overline{T}} =\{ a_0+ k b\ ~\text{with }~ a_0 \in \mathbb{Z}_{\mu} ~\text{ and }~ 0\leq k\leq \mu-1 \}.$$
This question was also addressed in \cite[Proposition 30]{Alfarano2021}, in the context of skew cyclic codes, when $\mu$ is a prime number.
\end{remark}
%\hlc[green]{In our case, with the setting of skew polycyclic codes, the result is still not clear.}

We finish this section with an example of a skew polycyclic code and its designed as well as actual minimum Hamming and rank distance.
\begin{example} 
Let $ \mathbb{F}_{2^6}= \mathbb{F}_{2}(w), $ with $ w^6 + w^4 + w^3 +  w+ 1=0 $ and  $f(x)=x^{10} + w^{40}x^9 + w^{39}x^8 + w^{12} x^6 + w^{46}x^5 + w^{42}x^4 + w^{60} x^2 + w^7x +
    w^{54}\in \mathbb{F}_{2^6}[x;\sigma]$ with $\sigma$ the Frobenius automorphism of $ \mathbb{F}_{2^6}$ then its order is $\mu=6.$ Using Magma and SageMath we found that the left exponent  of $f$ is $e:= 12.$  As $e$ is a multiple of $\mu$, the Frobenius automorphism $\theta$ of $\mathbb{F}_{2^{12}}$ is an extension of degree $ m=2 $ of $\sigma.$ Let $\gamma$ be a generator of $ \mathbb{F}_{2^{12}}$ such that $ \gamma^12 + \gamma^7 + \gamma^6 + \gamma^5 + \gamma^3 + \gamma + 1, $   and $ \alpha=\gamma^5 $ be a normal element of  $\mathbb{F}_{2^{12}}.$  Set $ \beta:=\theta(\alpha)\alpha^{-1},$ then the skew polynomial  $ x^{ 12}-1$ can be decomposed in $ \mathbb{F}_{2^{12}}[x;\sigma]$ as 
    $$  x^{ 84}-1 = \lclm \left( x-\beta, x-\theta(\beta), \ldots, x-\theta^{11}(\beta)\right)= \lclm \left( M_{\beta}(x), M_{\theta(\beta)}(x), \ldots, M_{\theta^5(\beta)}(x)\right)  $$
with
    $$\begin{array}{l}
      M_{\beta}(x)=  \lclm ( x-\beta, x-\theta^6(\beta))= x^2 + w^{38}x + w^{58} \\    
      M_{\theta(\beta)}(x)=  \lclm ( x-\theta(\beta), x-\theta^7(\beta))= x^2 + w^{13} x + w^{53}  \\
      
      M_{\theta^2(\beta)}(x)=  \lclm ( x-\theta^2(\beta), x-\theta^8(\beta))= x^2 + w^{26} x + w^{43}\\
      M_{\theta^3(\beta)}(x)=  \lclm ( x-\theta^3(\beta), x-\theta^9(\beta))= x^2 + w^{52}x + w^{23} \\ M_{\theta^4(\beta)}(x)=  \lclm ( x-\theta^4(\beta), x-\theta^10(\beta))= x^2 + w^{41} x + w^{46}\\
      M_{\theta^5(\beta)}(x)=  \lclm ( x-\theta^5(\beta), x-\theta^11(\beta))= x^2 + w^{19}x + w^{29}. 
    \end{array}
    $$
  Let now $g= x^4 + w^{52}x^3 + w^{46}x^2 + w^{23}x + w^{33 }$ be a right divisor of $f(x).$ For  $ a=2,\ b=1,\ \delta=3,\  r=1 ,$ we find that the $\beta$-defining set of $g$ is   $T_{\beta}(g)=\{ 2,3,8,9 \}$ and so $g= \lclm(x-\theta^2(\beta), x-\theta^3(\beta), x-\theta^8(\beta), x-\theta^9(\beta)) . $ Let $C$ be the skew polycyclic code induced by $f$ and generated by $g.$
  According to Theorem \ref{thm:Roos}, $C$ has minimum Hamming distance $d_H(C)\geq 3+1=4, $ and by the aid of magma we found that the exact Hamming  distance is $d_H(C)=4,$ and hence that $C$ is an $[10,6,4]_{2^6}$ code.
  
  Similarly, according to Theorem \ref{Rank_bound},  $C$  has rank distance (over $\mathbb{F}_2$) $d_R(C)\geq 3+1=4, $ and with the aid of Magma we found that the exact rank distance $d_R(C)=4,$ since it is a codeword $c=(0, 0, 1, 0, 0, 0, w^{37}, w^{57}, 0,    w^7 )$ of rank weight $4.$ Therefore $C$ is an $[10,6,4]_{2^6}$ rank-metric code.

\end{example}






\section{Equivalence  of  skew  polycyclic codes over $\Fq$}

Generalizing the usual equivalence between codes, we consider equivalence maps between two ambient spaces of skew polycyclic codes which preserve the algebraic structure of the code. This notion was first introduced in \cite{Chen2012}, to study the equivalence between two classes of classical constacyclic  codes over $\Fq.$ Recently it was extended to the case of skew constacyclic codes \cite{Boulanouar2021,Ouazzou2025}, \cite[Section 6]{Lobillo2025},  and to the case of polycyclic codes \cite{Equiv2025, Chibloun2025polycyclic}.  

\begin{definition}\label{FqIso}  
Let $f_1$ and $f_2$ be two skew polynomials from $\mathbb{F}_q[x;\sigma]$. We say that $f_1$  is \emph{Hamming} (resp. \emph{rank}) \emph{equivalent} to $f_2$, and we denote this by $ f_1 \cong_{(n,\sigma)-\text{Hamming}} f_2, $ (resp. $\cong_{(n,\sigma)-\text{Rank}}$)  if there is an $\mathbb{F}_q$-vector space isomorphism 
$$\varphi ~:~  \mathbb{F}_q[x;\sigma] /\langle f_2(x)\rangle   \quad \longrightarrow  \quad  \mathbb{F}_q[x;\sigma] /\langle 
f_1(x)\rangle,$$
such that, for all $ f(x), g(x) \in \Fq[x;\sigma] /\langle f_2(x) \rangle,$ we have 
\begin{enumerate}  
    \item $\varphi(g(x)f(x))= \varphi(g(x)) \varphi(f(x)) $.  
    \item $d_H(\varphi(f(x)), \varphi(g(x)))=d_H(f(x), g(x))$ \\ (resp. $ d_R(\varphi(f(x)), \varphi(g(x)))=d_R(f(x), g(x))$).
\end{enumerate}  
We will call such an isomorphism $\varphi$ an \emph{$\mathbb{F}_q$-morphism isometry} with respect to the Hamming (resp. rank)   distance.
\end{definition}  

We remark that, in this section,  we consider the rank distance defined over $ \Fq/ \mathbb{F}_{q'}$, where $ \mathbb{F}_{q'} $ is an arbitrary subfield of $ \Fq.$

\begin{proposition}\label{ProofEquiv}
Let $ R_{ n}$ be the set of monic skew polynomials of degree $n.$ 
\begin{comment}
 \textcolor{blue}{Monic  or not polynomial does not gives differences here, since we speak in general about the existing isometry between two classes of codes defined by polynomials. I was also thinking that we can find such isometry that gives isometry between cyclic codes ($f=x^n-1$) and polycyclic codes ( a given polynomial $ f(x)$ of degree $n$)} \textcolor{purple}{Yes, but you have to be precise with the words you use. If you restrict to monic, then it is not a ring.}
 \textcolor{blue}{yes, I didn't put attention to the structure, Thank you for mentioned this mistake. I changed the word "ring" by the "set".}
\end{comment}
Then the relations $\cong_{(n,\sigma)-Hamming}$ and  $\cong_{(n,\sigma)-Hamming}$  are equivalence relations on $ R_{n}.$ 
\end{proposition}
\begin{proof}
    See Appendix.
\end{proof}





 
% \textcolor{purple}{I like the first approach of explaining the stuff above. However, if we start with vector space isomorphisms, we should not already assume that we replace $f(X)$ with $f(\alpha x^h$, but it could be any polynomials inside. The other way around, I think there are theorems about possible forms of algebra morphisms of the form $u^{-1}(x) f(\alpha x) u(x)$. Then we have to say something about $u(x)$ because of the metric.}

%\hlc[green]{ I added the old explanation. For the other way if we consider   $ u(x) \in \Fq[x,\sigma]/\langle f(x) \rangle$ to be invertible, we need also to assume that $  \Fq[x,\sigma]/\langle f(x) \rangle $ to be a ring and so $f(x)$ to be two sided polynomial. But the in our situation we didn't consider that.}

\begin{comment}
\textcolor{purple}{Maybe something like this: As commonly done in the literature, we will focus on $\Fq$-morphism isometries of the form  $\varphi_{\alpha}(f(x)) = f(\alpha x)$, for some $\alpha \in \mathbb{F}_q^*$. Note that any such map preserves the code length $n$, the dimension $k$, and the Hamming distance; however, it only preserves the rank distance if $\alpha\in \Fq^{\sigma}$.} \textcolor{green}{Now, this is exactly what I did not want -- the subfield for the definition of rank does not necessarily need to be the stabilized subfield...}
\hlc[green]{I see what you mean. I will focus on conditions for which $\alpha\in \Fq \backslash \Fq^{\sigma}$ $ \varphi_{\alpha}(x)=f(\alpha x)$ is a rank isometry. } \textcolor{purple}{Good. But I think we still have all the results (also for the Roos bound) formulated such that the fixed subfield is the one over which the rank is defined. We should have to different subfields and also two different notations for this.} 
\hlc[green]{ I think the Roos bound does not work when considering a subfield other than the fixed subfield, since the theory is based on finding an extension that contains all the roots of our skew polynomial, and the existing results in this direction use the fixed subfield. For that, we construct the extension automorphism $\theta$ of $\sigma$, and the corresponding field extension.
} \textcolor{purple}{Ok, then we can mention it at the beginning of the previous section. However, for the equivalence it should be fine to consider two different subfields, right? Or -- if you prefer -- we write in the preliminaries section that for simplicity we restrict our study (for the whole paper) to the case where the fixed subfield is the one we take the rank over.}
\textcolor{blue}{ In the equivalence we used now the fixed subfield $\Fq^{\sigma}.$ But we can use another subfield  not necessary the fixed one, and the case we find that the number of equivalence classes for the trinomial case is 
$$N_{(n,\ell,\sigma)}:= \dfrac{ (q-1)^2}{ \lcm\left(
\frac{q'-1}{\gcd( \textcolor{red}{[n]_r},q'-1)}, \frac{q'-1}{ \gcd(\textcolor{red}{[n-\ell]_r},q'-1)}\right)} 
$$
with $q'$ is the size of the other subfield (not the fixed one,) In this case I have to edit the proofs of the rank equivalence since in the current one I used the fact that $ \alpha\in \Fq^{\sigma}.$ However the number of equivalence classes in the where we use the fixed subfieled is 
$$N_{( n,\ell, \sigma)-Rank}:= \dfrac{ (q-1)^2}{ \lcm\left(
\frac{q_0-1}{\gcd( \textcolor{red}{n},q_0-1)}, \frac{q_0-1}{ \gcd( \textcolor{red}{n-\ell},q_0-1)}\right)} 
$$
with $q_0$ is the size of the $ \Fq^{\sigma.}$
The difference between them is marked is black.
}
Note that any $\mathbb{F}_q$-vector space morphism of the form $\varphi_{\alpha}(f(x)) = f(\alpha x)$, for some $\alpha \in \mathbb{F}_q^*$  preserves the code length $n$, the dimension $k$, and the Hamming distance. 
However  if $\alpha\in \mathbb{F}_{q'} $  then it preserves the rank distance too. 
\end{comment}

As commonly done in the literature, we will focus on $\Fq$-morphism isometries of the form  $\varphi_{\alpha}(f(x)) = f(\alpha x)$, for some $\alpha \in \mathbb{F}_q^*$. Note that any such map preserves the code length $n$, the dimension $k$, and the Hamming distance; however, it only preserves the rank distance if $\alpha\in \mathbb{F}_{q'}$.\footnote{Similarly, any $ \mathbb{F}_{q'}$-vector space isomorphism $\varphi$ preserves the code length $n$, the dimension $k$, and the rank distance, but not necessarily the Hamming metric.}
\begin{comment}
\textcolor{purple}{It does not preserve the Hamming distance, that is exactly the point: only specific ones (even for $h=1$) are isometries, but in general they are not! In particular, $f(\alpha x^h), h>1$ is only a Hamming isometry if the quotient polynomials are binomials. For $h=1$ the correspond to monomial maps, so that is fine.} 
\textcolor{purple}{No, here you would have that all vector space isomorphisms with entries in the subfield preserve the rank, if the quotient polynomial does not mess it up.} 
\end{comment}
However, such maps do not necessarily preserve the algebraic structure of the codes, i.e., the skew polycyclical property  of the code. To ensure this, we must verify that
$$
\varphi(f(x)  g(x)) = \varphi(f(x))  \varphi(g(x)).
$$
Verifying this will be the main objective in the following sections. For simplicity we will first focus on trinomial codes, first in the Hamming and then in the rank metric, before we generalize the results to general polycyclic codes in the third subsection.

\begin{comment}
\hlc[green]{Do you think that we need to keep the following paragraph in blue?} \textcolor{purple}{No, let's erase it.}
\textcolor{blue}{Let for example $ f_1=x^n-a_0$ and $f_2=x^n-b_0.$
\begin{enumerate}
    \item  We consider the map
    $$\varphi: \  \mathbb{F}_q[x;\sigma]/\langle f_2(x)\rangle \longrightarrow  \mathbb{F}_q[x;\sigma]/\langle f_1(x)\rangle; \ \ x \longrightarrow  \alpha x^{n-1}, $$ for some $ \alpha \in \Fq^*.$ We can verify easily that $ \varphi$ is  $ \Fq$-vector space homomorphism, and with a simple computation we obtain that 
    $$ \varphi(x^i)=   N_{i(n-1)}^{\sigma}(\alpha)  x^{i(n-1)}=   N_{i(n-1)}^{\sigma}(\alpha)  x^{in} x^{-i}=  N_{i(n-1)}^{\sigma}(\alpha)  N_{i}^{\sigma^n}(a_0)  x^{-i}  \mod (x^n-a_0),$$
    and so
    $$ \varphi(x^n)= N_{n(n-1)}^{\sigma}(\alpha) x^{n(n-1)}=  N_{n(n-1)}^{\sigma}(\alpha) N_{n-1}^{\sigma^n}(a_0)  \mod (x^n-a_0),$$
    on the other hand  as $ x^n-b_0= b_0 \mod (x^n-b_0),$ then 
    $$  \varphi(x^n)= \varphi(b_0 )= b_0 \mod( x^n-a_0)  ,$$ 
    Next we combine the two values of $ \varphi(x^n) $ and  the equivalence gives us this equality
    $$   b_0=  N_{n(n-1)}^{\sigma}(\alpha) N_{n-1}^{\sigma^n}(a_0)   $$
To preserve the length $n$ and dimension $k$, $ \varphi$ should be and $\Fq$- vector space isomorphism. Then for $   b_0=  N_{n(n-1)}^{\sigma}(\alpha) N_{n-1}^{\sigma^n}(a_0)$ we consider 
    $$\bar \varphi: \  \mathbb{F}_q[x;\sigma]  \longrightarrow  \mathbb{F}_q[x;\sigma]/\langle f_1(x)\rangle; \ \ x \longrightarrow x^{n-1}. $$
We can verify that 
$ \bar \varphi( x^n-b_0)= \varphi(x^n)-b_0=   N_{n(n-1)}^{\sigma}(\alpha) N_{n-1}^{\sigma^n}(a_0) -b_0=0.  $
and so we can prove $\varphi$  is an $\Fq$-isomorphism. Hence  it preserves both  rank and the Hamming metric.
\noindent Let check the $\Fq$-morphism property, i.e we need to prove that $$ \varphi(a_ix^i b_jx^j)= \varphi(a_ix^i) \varphi(b_j x^j).$$ Fist we have 
$$a_i x^i  b_j x^j = a_i  \sigma^i(b_j)  x^{i+j}.$$
Then 
$$ \varphi( a_i x^i  b_j x^j ) = a_i  \sigma^i(b_j)  N_{(i+j)(n-1)}^{\sigma}(\alpha)  N_{i+j}^{\sigma^n}(a_0)  x^{-(i+j)} , . $$
In the other hand, 
$$ \varphi( a_i x^i)  \varphi( b_j x^j ) = a_i  N_{i(n-1)}^{\sigma}(\alpha) N_i^{\sigma^n}(a_0)  \sigma^{-i}(b_j)  \sigma^{-i}(N_{j(n-1)}^{\sigma}(\alpha))   \sigma^{-i}(N_j^{\theta}(a_0))  x^{-(i+j)}. $$
\noindent  For $\alpha=1,$ we need  for  the $\Fq$-morphism property  that 
$$   \sigma^i(b_j N_j^{\theta}(a_0)) = \sigma^{-i}(b_j N_j^{\theta}(a_0)),$$
which is not generally true unless   $ \sigma^{2i}(a) =  a$, for all $a\in \Fq$ and $ i =0,\ldots, n-1$. Therefore, $ \varphi $ is not a ring $\Fq$-morphism in general.
In conclusion $\varphi $ is not an $\Fq$-morphism isometry.  
\item  If we consider the map
    $$\varphi: \  \mathbb{F}_q[x;\sigma]/\langle f_2(x)\rangle \longrightarrow  \mathbb{F}_q[x;\sigma]/\langle f_1(x)\rangle; \ \ x \longrightarrow   x^{n-1}+x^{n-2}.$$
    We can prove that if $b_0= N^{\sigma^n}_{n-1}(a_0)+ N^{\sigma^n}_{n-2}(a_0),$ then $\varphi$  is an $\Fq$-vector space isomorphism. Hence it preserves the rank metric. But again is not an $ \Fq$-morphism !!!!!
\item More general if $ \varphi(x)= x^{n-1}+x^{n-2}+ \ldots+ x+1 .$ We can prove that if  $$ b_0=   N^{\sigma^n}_{n-1}(a_0)+ N^{\sigma^n}_{n-2}(a_0) + \ldots+N_1^{\sigma^n}(a_0)+1 ,$$ 
then $\varphi$  is an $\Fq$-vector space isomorphism. Hence it preserves the rank metric. But again is not an $ \Fq$-morphism.
    \end{enumerate}
If  $ f_1= x^n-b_{\ell} x^{\ell}-b_0$   and $ f_2= x^n-a_{\ell} x^{\ell}-a_0,$  the situation is more complicated. So for these reasons, in this article we restrict our study to the case where $\varphi(x) = \alpha x$, and we say that $f_1$ and $f_2$ are $(n, \sigma)$-equivalent. In future work, we will further investigate how to preserve the algebraic structure of such codes. 
We begin with the case of trinomial codes, and then generalize to arbitrary polycyclic codes. Finally, we conclude the section by considering the restriction $\alpha \in \mathbb{F}_q^{\sigma}$.}
\end{comment}

\subsection{ Hamming equivalence  of  skew  trinomial codes}

Now, we give the following special case of equivalence,  which we refer to as  {Hamming $(n,\ell,\sigma)$-equivalence}, between two classes of $(\ell,\sigma)$-trinomial codes. 

 \begin{definition}\label{Def_Iso}
Let $a_0, a_{\ell}, b_0, b_{\ell}$ be nonzero elements of $\mathbb{F}_q$, and let $\ell$ be an integer such that $0 < \ell < n$. We say that $(a_0, a_{\ell})$ and $(b_0, b_{\ell})$ are { Hamming $(n,\ell,\sigma)$-equivalent} 
in $\mathbb{F}_q^* \times \mathbb{F}_q^*$, denoted by 
$$
(a_0, a_{\ell}) \sim_{(n,\ell,\sigma)} (b_0, b_{\ell}),
$$
if there exists an $\alpha \in \mathbb{F}_q^*$ such that the following map:

  \begin{equation}
  	\begin{array}{cccc}
  	\varphi_{\alpha}:& \mathbb{F}_q[x;\sigma] /\langle x^n-b_{\ell}x^{\ell}-b_0 \rangle  &\longrightarrow & \mathbb{F}_q[x;\sigma] /\langle x^n-a_{\ell}x^{\ell}-a_0 \rangle, \\ 
  	& f(x) & \longmapsto &  f(\alpha x),
  	\end{array}
  	\end{equation} 
is an  $\mathbb{F}_q$-morphism isometry with respect to the Hamming distance,  i.e., 
$$\varphi_{\alpha}( f(x) g(x))= \varphi_{\alpha}( f(x)) \varphi_{\alpha}( g(x)), \quad \text{and } \quad d_H(\varphi_{\alpha}(f(x)), \varphi_{\alpha}(g(x)))=d_H(f(x), g(x)), $$ 
for all $ f(x), g(x) \in \Fq[x;\sigma] /\langle x^n-b_{\ell}x^{\ell}-b_0 \rangle.$
\end{definition}






\begin{remark}\label{Remark.3}
\begin{enumerate}
 \item  Note that the $(n,\ell,\sigma)$-equivalence relation in the above definition generalizes the $(n,\sigma)$-equivalence of constacyclic codes studied in \cite[Section 4]{Ouazzou2025}, which was denoted by $ \lambda \sim_{(n,\sigma)} \mu .$
 \item  If $\sigma=\id$ we recover  the case of $(n,\ell)$-equivalence for  $\ell$-trinomial   codes studied in \cite{Equiv2025}.

\end{enumerate}
\end{remark}


In the following theorem we give essential characterizations of  the Hamming $(n,\ell,\sigma)$-equivalence   between two classes of skew $(\ell,\sigma)$-trinomial codes of length $n$ over $\mathbb{F}_q$. 


\begin{theorem}\label{Th.1}
 Let $0<\ell<n$ be an integer,  $(a_0, a_{\ell})$ and $(b_0, b_{\ell})$ be elements of $\mathbb{F}_q^{*} \times \mathbb{F}_q^{*}$,   $\xi$ be  a primitive element of $\Fq$   and  $ \sigma$ be an automorphism of $\Fq$. The following statements are equivalent:

\begin{enumerate}
    \item $(a_0, a_{\ell}) \sim_{(n,\ell,\sigma)} (b_0, b_{\ell}),$ i.e.,  $(a_0, a_{\ell}) $ and  $ (b_0, b_{\ell})$ are Hamming \textcolor{purple}{$(n,\ell,\sigma)$} equivalent.  

\item   There exists $\alpha\in \Fq^*$ such that $$ a_0 N_n^{\sigma}(\alpha)  = b_0   \quad  \text{and} \quad   a_{\ell} N^{\sigma}_{n-\ell}(\sigma^{\ell}(\alpha))=  b_{\ell}. $$
\item  There exists $\alpha\in \Fq^*$ such that $$ (a_0, a_{\ell}) \star
\left( N_n^{\sigma}(\alpha), N_{n-\ell}^{\sigma}(\sigma^{\ell} (\alpha)) \right)= (b_0,b_{\ell}) .$$
    
 \item   $ (a_0,a_{\ell})^{-1}\star (b_0,b_{\ell})\in H_{\ell,\sigma},$ where   $H_{\ell,\sigma}$ is the cyclic subgroup of $\Fq^*\times \Fq^*$ generated by $\left(N^{\sigma}_{n}(\xi), N_{n-\ell}^{\sigma}(\sigma^{\ell} (\xi))   \right).$
\end{enumerate}
The equivalence between (1) and (4) implies that the number of $(n,\ell,\sigma)$-equivalence classes is $$N_{(n,\ell,\sigma)}:= \dfrac{ (q-1)^2}{ \lcm\left(
\frac{q-1}{\gcd([n]_r,q-1)}, \frac{q-1}{ \gcd([n-\ell]_r,q-1)}\right)} .$$
\end{theorem}


\begin{proof}\;
\begin{itemize}
    \item[(1) $\Rightarrow$ (2):]
   Suppose that $(a_0, a_{\ell}) \sim_{(n,\ell,\sigma)} (b_0, b_{\ell})$. Then by Definition \ref{Def_Iso}, there exists $\alpha \in \mathbb{F}_q^*$ such that the map
    $$
    \varphi_{\alpha} : \Fq[x;\sigma] /\langle x^n - b_{\ell} x^{\ell} - b_0 \rangle \to \Fq[x;\sigma] /\langle x^n - a_{\ell} x^{\ell} - a_0 \rangle, \quad f(x) \ \mapsto \ f(\alpha x)
    $$
    is an  $\mathbb{F}_q$-morphism isometry.  It follows that
$$ 
\varphi_{\alpha}(x^i)= (\alpha x)^i= N_i^{\sigma}(\alpha) x^{ i}\ \text{mod}_r (x^n-a_{\ell} x^{\ell}-a_0), \ \forall  \ i=0,1, \ldots, n-1.
$$

As $\varphi_{\alpha}$ is an $\Fq$-morphism and  $ \varphi(x^n-b_{\ell}x^{\ell}-b_0)=0 \ \text{mod}_r (x^n- a_{\ell} x^{\ell}-a_0) ,$ we have
\begin{equation}
\varphi( x^n)=   b_{\ell}  N_{\ell}^{\sigma}(\alpha) x^{\ell} +b_0.
 \end{equation}

On the other hand,
\begin{equation}
\varphi( x^n)=  N_n^{\sigma}(\alpha) x^n= N_n^{\sigma}(\alpha) ( a_{\ell}x^{\ell}+a_0) = N_n^{\sigma}(\alpha) a_{\ell} x^{\ell}+ N_n^{\sigma}(\alpha) a_0.
\end{equation}


Comparing term by term, we obtain  that $$ a_0 N_n^{\sigma}(\alpha)  = b_0, \quad    \text{and} \quad   a_{\ell} N_n^{\sigma}(\alpha) =  N_{\ell}^{\sigma}(\alpha) b_{\ell}. $$ It follows that
$$ a_0 N_n^{\sigma}(\alpha)  = b_0   \quad  \text{and} \quad   a_{\ell} N_n^{\sigma}(\alpha)  N_{\ell}^{\sigma}(\alpha^{-1})=  b_{\ell}. $$
 On the other hand we have 
\begin{equation}\label{Eq_formulae}
\begin{array}{rl}
    N_n^{\sigma}(\alpha) N_{\ell}^{\sigma}(\alpha^{-1})  & =  N_{\ell+(n-\ell)}^{\sigma}(\alpha) N_{\ell}^{\sigma}(\alpha^{-1})\\
     &=   \sigma^{\ell}(N_{n-\ell}^{\sigma}(\alpha)) N^{\sigma}_{\ell}(\alpha) N_{\ell}^{\sigma}(\alpha^{-1})  \\
     &=  \sigma^{\ell}(N^{\sigma}_{n-\ell}(\alpha))\\
     & = N^{\sigma}_{n-\ell}(\sigma^{\ell}(\alpha))
\end{array}
\end{equation}
Hence we obtain that 
$$ a_0 N_n^{\sigma}(\alpha)  = b_0   \quad  \text{and} \quad   a_{\ell} N^{\sigma}_{n-\ell}(\sigma^{\ell}(\alpha))=  b_{\ell}. $$
 

 \item[(2) $\Rightarrow$ (3):] Immediate.
 \item[(3) $\Rightarrow$ (4):] Suppose that

$$  (b_0, b_{\ell} )= \left( N_{n}^{\sigma}(\alpha) a_0, N_{n-\ell}^{\sigma}(\sigma^{\ell}(\alpha) a_{\ell}\right)=\left(N_{n}^{\sigma}(\alpha), N_{n-\ell}^{\sigma}(\sigma^{\ell}(\alpha))\right)\star (a_0, a_{\ell}).$$ 
It follows that
$$  (a_0,a_{\ell})^{-1}\star (b_0, a_{\ell})=  \left( N_{n}^{\sigma}(\alpha), N_{n-\ell}^{\sigma}( \sigma^{\ell}(\alpha)) \right)= \left( N_{n}^{\sigma}(\xi^j),
N_{n-\ell}^{\sigma}(\sigma^{\ell}(\xi^j))\right)= 
\left(N_{n}^{\sigma}(\xi), N_{n-\ell}^{\sigma}(\sigma^{\ell}(\xi))\right)^j .$$
Then $ (a_0,a_{\ell})^{-1}\star (b_0, a_{\ell}) $ belongs to the cyclic subgroup  $H_{\ell,\sigma}$ generated by $  \left(N_{n}^{\sigma}(\xi), N_{n-\ell}^{\sigma}(\sigma^{\ell}(\xi))\right) $  as a subgroup of $ \Fq^{*}\times \Fq^{*}.$

    \item[(4) $\Rightarrow$ (1):] 
     Suppose that $ (a_0,a_{\ell})^{-1}\star (b_0, b_{\ell}) $ is an element of the cyclic subgroup  $H_{\ell,\sigma}$ generated by 
     $  (N_{n}^{\sigma}(\xi), N_{n-\ell}^{\sigma}(\sigma^{\ell}(\xi)))  $  as a subgroup of $ \Fq^{*}\times \Fq^{*}.$ Then there exists an integer $ h$ such that 
     $$ (a_0,a_{\ell})^{-1}\star (b_0, b_{\ell})= \left(N_{n}^{\sigma}(\xi), N_{n-\ell}^{\sigma}(\sigma^{\ell}(\xi))\right)^h = \left( N_{n}^{\sigma}(\xi^h), N_{n-\ell}^{\sigma}(\sigma^{\ell}(\xi^h))\right).$$
    For $ \alpha=  \xi^{h}, $ we obtain  that $  a_i N_{n-i}^{\sigma}( \sigma^{i}(\alpha))= b_i, $ for any $ i\in \{ 0,\ell\}.$
Now, let us consider the map $ \tilde{\varphi}_{\alpha}$  
  \begin{equation}
        \begin{array}{cccc}
        \tilde{\varphi}_{\alpha} : & \mathbb{F}_q[x;\sigma]   & \longrightarrow & \mathbb{F}_q[x;\sigma] /\langle x^n - a_{\ell} x^{\ell} - a_0 \rangle, \\ 
        & f(x) & \longmapsto & f(\alpha x).
        \end{array}
    \end{equation} 
	$ \tilde{\varphi}_{\alpha}$ is a surjective $\Fq$-vector space homomorphism. Indeed, for  all $0\leq j \leq n-1$, $x^{j} = \tilde{\varphi}_{\alpha}(N_{j}^{\sigma}(\alpha^{-1}) x^{j}).$ 
	Moreover, 
 $$ 
    \begin{array}{rl}
    \tilde{\varphi}_{\alpha} (x^n - b_{\ell} x^{\ell} - b_0) & = N^{\sigma}_n(\alpha) x^n - N^{\sigma}_{\ell}(\alpha) b_{\ell} x^{\ell} - b_0 \\
    & = N^{\sigma}_n(\alpha) x^n - N^{\sigma}_{\ell}(\alpha)  N_{n-\ell}^{\sigma}( \sigma^{\ell}(\alpha)) a_{\ell} x^{\ell} - N^{\sigma}_n(\alpha) a_0 \\
    & = N^{\sigma}_n(\alpha) x^n -  \sigma^{\ell}(N_{n-\ell}^{\sigma}( \alpha))  N^{\sigma}_{\ell}(\alpha)  a_{\ell} x^{\ell} - N^{\sigma}_n(\alpha) a_0 \\
    & = N^{\sigma}_n(\alpha) x^n - N_{\ell+ n-\ell}^{\sigma}( \alpha)  a_{\ell} x^{\ell} - N^{\sigma}_n(\alpha) a_0 \\
     & = N^{\sigma}_n(\alpha) x^n - N^{\sigma}_n(\alpha)  a_{\ell} x^{\ell} - N^{\sigma}_n(\alpha) a_0 \\
    & = N^{\sigma}_n(\alpha) (x^n - a_{\ell} x^{\ell} - a_0) \\
    & = 0 \mod (x^n - a_{\ell} x^{\ell} - a_0).
    \end{array}
    $$ 
	So $\langle x^n - b_{\ell} x^{\ell} - b_0\rangle \subseteq \ker\tilde{\varphi}_{\alpha}$. And for all $f(x)\in \ker\tilde{\varphi}_{\alpha}$,
    $$f(\alpha x)= 0~ (\text{mod } x^{n}- a_{\ell} x^{\ell} - a_0).$$ Therefore, there exists $g(x)\in \Fq[x;\sigma]$ such that $f(\alpha x)=g(x)(x^{n}- a_{\ell} x^{\ell} - a_0)$. Thus, 
 $$
  \begin{array}{rl}
 f(x)&=g(\alpha^{-1}x)( N^{\sigma}_n(\alpha^{-1}) x^{n}- N^{\sigma}_{\ell}(\alpha^{-1}) a_{\ell} x^{\ell} - a_0) \\
 & =N^{\sigma}_n(\alpha^{-1}) g(\alpha^{-1}x)(x^{n}- N^{\sigma}_{n-\ell}(\alpha) a_{\ell} x^{\ell} - N^{\sigma}_n(\alpha) a_0)\\
 &= N^{\sigma}_n(\alpha^{-1}) g(\alpha^{-1}x)( x^n-  b_{\ell} x^{\ell} - b_0), \ \text{since  $a_k N^{\sigma}_{n-k}( \sigma^k(\alpha)) =b_k, \ \forall k \in \{0, \ell\}.$} \\ 
 \end{array}$$
	So $\ker\tilde{\varphi}_{\alpha} \subseteq \langle x^{n}- a_{\ell} x^{\ell} - a_0\rangle $ and hence $$\ker\tilde{\varphi}_{\alpha} = \langle x^{n}-  b_{\ell} x^{\ell} - b_0 \rangle.$$
	Therefore  by the  first isomorphism  theorem, the map 
      \begin{equation}
        \begin{array}{cccc}
        \varphi_{\alpha} : & \mathbb{F}_q[x;\sigma]/\langle x^n - b_{\ell} x^{\ell} - b_0 \rangle  & \longrightarrow & \mathbb{F}_q[x;\sigma] /\langle x^n - a_{\ell} x^{\ell} - a_0 \rangle, \\ 
        & f(x) & \longmapsto & f(\alpha x).
        \end{array}
    \end{equation} 
	is an  isomorphism. As  the weights of $f(x)$ and $f(\alpha x)$ are the same,  $\varphi_{\alpha}$  is an isometry.
    
    We now need  to verify that $\varphi_{\alpha} $ is a $\Fq$-morphism.
    Let $ f(x)=\dsum{i=0}{n-1} f_ix^i$ and $g(x)=  \dsum{i=0}{n-1} g_ix^i$  in $ \Fq[x;\sigma]/\langle x^n- b_{\ell}x^{\ell}-b_0 \rangle.$ On the one hand we have 

           $$
           \begin{array}{rl}
f(x)g(x) &=\dsum{j=0}{n-1}\dsum{i=0}{j} f_i \sigma^i\left(g_{j-i}\right)  x^j+ 
\dsum{j=0}{n-1} \dsum{i=j+1}{n-1} f_i \sigma^i\left(g_{n-i+j}\right)  x^{j+n} \\
 
\end{array}
$$
as $x^{j+n}= \sigma^{j}(b_{\ell}) x^{j+\ell} + \sigma^{j}(b_{0}) x^{j+\ell} \mod \ (x^n-b_{\ell}x^{\ell} -b_0).$ Then
$$
f(x)g(x) = \sum_{j=0}^{n-1} \left( \sum_{i=0}^{j} f_i \sigma^i(g_{j-i}) \right) x^j +
\sum_{j=0}^{n-1} \left( \sum_{i=j+1}^{n-1} f_i \sigma^i(g_{n-i+j}) \sigma^j(b_\ell)x^{j+\ell} + \sum_{i=j+1}^{n-1} f_i \sigma^i(g_{n-i+j}) \sigma^j(b_0)x^j \right).
$$
It follows that
$$
f(x)g(x) = \sum_{j=0}^{n-1} \left( \sum_{i=0}^{j} f_i \sigma^i(g_{j-i}) + \sum_{i=j+1}^{n-1} f_i \sigma^i(g_{n-i+j}) \sigma^j(b_0) \right) x^j +
\sum_{j=0}^{n-1} \left( \sum_{i=j+1}^{n-1} f_i \sigma^i(g_{n-i+j}) \sigma^j(b_\ell) \right) x^{j+\ell}.
$$

Next, as $\varphi_\alpha(x^j) = N_j^\sigma(\alpha)x^j$ and $\varphi_\alpha(x^{j+\ell}) = N_{j+\ell}^\sigma(\alpha)x^{j+\ell},$ we have  
\begin{equation}\label{eq11}
\begin{array}{rl}
\varphi_\alpha(f(x)g(x)) &= \dsum{j=0}{n-1} \left( \dsum{i=0}{j} f_i \sigma^i(g_{j-i})  N_j^\sigma(\alpha) +
 \dsum{i=j+1}{n-1} f_i \sigma^i(g_{n-i+j}) \sigma^j(b_0)  N_j^\sigma(\alpha) \right)x^j \\

&+\dsum{j=0}{n-1} \left( \dsum{i=j+1}{n-1} f_i \sigma^i(g_{n-i+j}) \sigma^j(b_\ell) N_{j+\ell}^\sigma(\alpha) \right) x^{j+\ell}.
\end{array}
\end{equation}

On the other hand, we   compute $\varphi_\alpha(f(x))\varphi_\alpha(g(x)).$
$$
\begin{array}{rl}
 \varphi_\alpha(f(x))\varphi_\alpha(g(x))    &=  \left( \dsum{i=0}{n-1} f_i N_i^\sigma(\alpha)x^i \right) \left( \dsum{j=0}{n-1} b_j N_j^\sigma(\alpha)x^j \right) \\
    & = \dsum{j=0}{n-1} \left( \dsum{i=0}{j} f_i N_i^\sigma(\alpha) \sigma^i(g_{j-i}N_{j-i}^\sigma(\alpha)) \right) x^j \\
    &  +\dsum{j=0}{n-1} \left( \dsum{i=j+1}{n-1} f_i N_i^\sigma(\alpha) \sigma^i(g_{n-i+j}N_{n-i+j}^\sigma(\alpha)) \right)x^{j+n}.
\end{array}
$$
As $x^{j+n} \equiv \sigma^j(a_\ell)x^{j+\ell} + \sigma^j(a_0)x^j \mod (x^n - a_\ell x^\ell - a_0)$,  we have 
\begin{equation}\label{eq12}
\begin{array}{rl}
\varphi_\alpha(f(x))\varphi_\alpha(g(x)) &= \dsum{j=0}{n-1} \left( \dsum{i=0}{j} f_i\sigma^i(g_{j-i})  N_i^\sigma(\alpha) \sigma^i(N_{j-i}^\sigma(\alpha))\right)x^j \\
&+ \dsum{j=0}{n-1} \left( \dsum{i=j+1}{n-1} f_i  \sigma^i(g_{n-i+j})  \sigma^j(a_0) N_i^\sigma(\alpha)\sigma^i(N_{n-i+j}^\sigma(\alpha))  \right)x^j 
\\
& +\dsum{j=0}{n-1} \left( \dsum{i=j+1}{n-1} a_i  \sigma^i(g_{n-i+j}) N_i^\sigma(\alpha) \sigma^i(N_{n-i+j}^\sigma(\alpha))\sigma^j(a_\ell) \right)x^{j+\ell}.\\
\end{array}
\end{equation}

By \cite[Proposition 2.1]{Cherchem2016},  
$
N_{i+j}^\sigma(\alpha) = N_i^\sigma(\alpha)\sigma^i(N_j^\sigma(\alpha)).
$
Therefore, 
$$ N_i^\sigma(\alpha) \sigma^i(N_{j-i}^\sigma(\alpha))= N_j^{\sigma}(\alpha) .$$ 
So, for any $k\in \{ 0, \ell\}$, we have 
%$$ 
%\begin{array}{rl}
% N_i^\sigma(\alpha)\sigma^i(N_{n-i+j}^\sigma(\alpha))  \sigma^j(a_0)   & = N_{n+j}^{\sigma}(\alpha)  \sigma^j(a_0)  \\
%  &=  \sigma^{ j} ( N^{\sigma}_{n}(\alpha)) N^{\sigma}_j(\alpha)  \sigma^{ j}(a_0) \\
%  &= \sigma^{ j} ( N^{\sigma}_{n}(\alpha))  N^{\sigma}_j(\alpha) \sigma^{ -j} (b_0)\sigma^{ j} ( N^{\sigma}_{n}(\alpha)), \ \ 
%  \text{since } \ 
%  a_{0}  N^{\sigma}_{n}(\alpha)  = b_{0}\\
%  &= \sigma^{ j} ( N^{\sigma}_{n}(\alpha)) N^{\sigma}_j(\alpha)  \sigma^{ j} (b_0) \sigma^{ -j} ( N^{\sigma}_{n}(\alpha))\\
% &= \sigma^{ j} (b_0)  N^{\sigma}_j(\alpha)  \\
%\end{array}
%$$

$$ 
\begin{array}{rl}
& N_i^\sigma(\alpha)\sigma^i(N_{n-i+j}^\sigma(\alpha))  \sigma^j(a_{k})   \\
 & = N_{n+j}^{\sigma}(\alpha)  \sigma^j(a_{k})  \\
  &=  \sigma^{ j} ( N^{\sigma}_{n}(\alpha)) N^{\sigma}_j(\alpha)  \sigma^{ j}(a_{k}) \\
  &= \sigma^{ j} ( N^{\sigma}_{n}(\alpha))  N^{\sigma}_j(\alpha) \sigma^{ j} (b_{k}) \sigma^{ -j} ( N^{\sigma}_{n-k}(\sigma^{k}(\alpha))), \ \ 
  \text{since } \ 
  a_{k}  N^{\sigma}_{n-k}(\sigma^{k}(\alpha))  = b_{k}\\
  &= \sigma^{ j} ( N^{\sigma}_{n}(\alpha)) N^{\sigma}_{j}(\alpha)   \sigma^{ k-j} ( N^{\sigma}_{n-k}(  \alpha)) \sigma^{ j} (b_{k}) \\
&=  \sigma^{ j} ( N^{\sigma}_{n}(\alpha)) N^{\sigma}_{j}(\alpha)   \sigma^{ k-j} ( N^{\sigma}_{n-k}(  \alpha))     N_{k-j}^{\sigma}(\alpha)  N_{k-j}^{\sigma}(\alpha^{-1}) \sigma^{ j} (b_{k})   \\ 
&=  \sigma^{ j} ( N^{\sigma}_{n}(\alpha)) N^{\sigma}_{j}(\alpha)   N_{n-j}^{\sigma}(\alpha) N_{k-j}^{\sigma}(\alpha^{-1}) \sigma^{ j} (b_{k})   \\ 
&=  \sigma^{ j} ( N^{\sigma}_{n}(\alpha)) N^{\sigma}_{j}(\alpha)  \sigma^{-j}(N_n^{\sigma}(\alpha)) N_{-j}^{\sigma}(\alpha) N_{k-j}^{\sigma}(\alpha^{-1}) \sigma^{ j} (b_{k})   \\ 
&=   N^{\sigma}_{j}(\alpha)  N_{-j}^{\sigma}(\alpha) N^{\sigma}_{-j}(\alpha^{-1}) \sigma^{-j}(N^{\sigma}_{k}(\alpha^{-1})) \sigma^{ j} (b_{k})   \\ 
&=   N^{\sigma}_{j}(\alpha)  \sigma^{-j}(N^{\sigma}_{k}(\alpha^{-1})) \sigma^{ j} (b_{k})   \\ 
 & =N^{\sigma}_{j}(\alpha)  \sigma^{j}(N^{\sigma}_{k}(\alpha)) \sigma^{ j} (b_{k})   \\ 
 & =N^{\sigma}_{k+j}(\alpha)  \sigma^{ j} (b_{k})   
\end{array}
$$
By comparing equations  (\ref{eq11}) and  (\ref{eq12}), we get 
$$
\varphi_\alpha(f(x)g(x)) = \varphi_\alpha(f(x))\varphi_\alpha(g(x)).
$$

Hence we have shown that $\varphi_\alpha$ is an $\mathbb{F}_q$-morphism which preserves the Hamming weight,  so
$(a_0, a_{\ell}) \sim_{(n,\ell,\sigma)} (a_0, a_{\ell}).$ Hence,  (4) $\Rightarrow$ (1).
\end{itemize}
%\hl{By the equivalence between (1) and (4),} \hlc[pink]{This equivalence is not established yet. This part of the proof is trying to establish $(4)\rightarrow (1)$} 
%\textcolor{blue}{This part is just to count the number of equivalence classes. See above that xe proved that $(a_0, a_{\ell}) \sim_{(n,\ell,\sigma)} (a_0, a_{\ell}),$ which means that  (4) $\Rightarrow$ (1). Then we have the equivalence}

Now, by the equivalence between (1) and (4) we deduce that the number of $(n,\ell,\sigma)$-equivalence classes on $\Fq^*\times \Fq^*$  corresponds to the order of the group $ \left( \Fq^*\times \Fq^* \right)/H_{\ell,\sigma},$ which equals 
 $$ N_{(n,\ell,\sigma)}= \dfrac{ (q-1)^2}{ \lcm \left(\frac{q-1}{ \gcd([n]_r,q-1)}, \frac{q-1}{ \gcd(p^{r\ell}[n-\ell]_r,q-1)}\right)}=
 \dfrac{ (q-1)^2}{ \lcm \left(\frac{q-1}{ \gcd([n]_r,q-1)}, \frac{q-1}{ \gcd([n-\ell]_r,q-1)}\right)}
 .$$ 

\end{proof}


\begin{remark}\label{Remark_TwoSided}
The reader may wonder what happens if we assume that the polynomials
$ x^n - a_{\ell}x^{\ell} - a_0 $ and $ x^n - b_{\ell}x^{\ell} - b_0 $
are elements of the center $ Z(\mathbb{F}_q[x; \sigma]) = \mathbb{F}_q^{\sigma}[x^{\mu}] $,
where $ \mu $ is the order of $ \sigma .$ In this case  $ \Fq[x;\sigma]/\langle  x^n-a_{\ell}x^{\ell}-a_0\rangle $ and $ \Fq[x;\sigma]/\langle  x^n-b_{\ell}x^{\ell}-b_0\rangle $ are $ \Fq^{\sigma}$-algebras (associative), and the definition of the $(n,\sigma)$-equivalence remains the same as in Definition \ref{Def_Iso}.  Moreover, all the statement of the Theorem \ref{Th.1}, remain the same except for the last one on the number of $(n,\ell,\sigma)$-equivalence classes which will be given by
    $$ \dfrac{ (\vert \Fq^{\sigma}\vert -1)^2}{ \lcm \left(\frac{\vert \Fq^{\sigma}\vert-1}{ \gcd([n]_r,\vert \Fq^{\sigma}\vert-1)}, \frac{\vert \Fq^{\sigma}\vert-1}{ \gcd([n-\ell]_r,\vert \Fq^{\sigma}\vert-1)}\right)}= \dfrac{ (q_0-1)^2}{ \lcm \left(\frac{q_0-1}{ \gcd([n]_r,q_0-1)}, \frac{q_0-1}{ \gcd([n-\ell]_r,q_0-1)}\right)}
 .$$ 
\end{remark}

%\vspace{.5cm}
Using the equivalence between the assertions $(1)$ and $(4)$ of Theorem \ref{Th.1} we derive the following characterization regarding the equivalence between the class of skew $(\ell,\sigma)$-trinomial codes  associated with  $x^n - a_{\ell}x^{\ell} - a_0$ and the class  associated with $x^n - x^{\ell} - 1.$
\begin{corollary}  \label{Equiv_1}
Let $\sigma$ be and automorphism of $ \Fq,$ $\ell$ be an integer such that  $0<\ell<n$,  and  $(a_0, a_{\ell})$ be an element of $\mathbb{F}_q^{*} \times \mathbb{F}_q^{*}$. 
Then the following statements are equivalent.
    \begin{enumerate}
        \item  The class of skew  $(\ell,\sigma)$-trinomial codes  associated with  $x^n - a_{\ell}x^{\ell} - a_0$ is $(n,\ell,\sigma)$-equivalent (Hamming equivalence) to the class of   skew $(\ell,\sigma)$-trinomial codes  associated with $x^n - x^{\ell} - 1.$
        \item  There exists $ \alpha \in \Fq^*$ such that $(a_0, a_{\ell})\star\left(N^{\sigma}_{n}(\alpha), N^{\sigma}_{n-\ell}(\sigma^{\ell}(\alpha)) \right)= (1, 1)$. 
        \item  There exists  an $(n,\sigma)$-th root  $\alpha \in \Fq^* $  of $a_0$     such that $   a_0 =  N^{\sigma}_{\ell}(\alpha) a_\ell $.
        \end{enumerate}
\end{corollary}
\begin{proof}\;
\begin{itemize}
    \item[(1) $\Rightarrow$ (2):] Follows from the equivalence of  (1) and (3) of Theorem \ref{Th.1}.
 \item[(2) $\Rightarrow$ (3):]  
 If $(a_0, a_{\ell})\star\left(N_n^{\sigma}(\alpha), N_{n-\ell}^{\sigma}(\sigma^{\ell}(\alpha))\right) = (1, 1)$, then 
$$ a_0 N_n^{\sigma}(\alpha)=1 \quad \quad 
\text{and} \quad  
 \quad a_\ell N_{n-\ell}^{\sigma}(\sigma^{\ell}(\alpha))=1.
$$
 It follows that
 $$  a_0 = N_n^{\sigma}(\alpha^{-1}) \quad \text{ and } \quad  a_\ell N_{n-\ell}^{\sigma}(\sigma^{\ell}(\alpha^{-1}) ) N_n^{\sigma}(\alpha^{-1}) = a_0.
  $$
 $$ \iff a_0 = N_n^{\sigma}(\alpha^{-1}) \quad \text{ and } \quad  a_{\ell} N_{\ell}^\sigma(\alpha^{-1})= a_0 $$
 
  i.e., $\beta:=\alpha^{-1}$ is an $(n,\sigma)$-th root of $a_0$ and $  a_{\ell} N_{\ell}^\sigma(\beta)= a_0  $.
 

 \item[(3) $\Rightarrow$ (1):] Suppose that $\alpha $ is an $(n,\sigma)$-th root of $a_0$ such that $ a_0 =  N^{\sigma}_{\ell}(\alpha) a_\ell.$

 $$ N_{n}^{\sigma}(\alpha)  = a_0\ \ \text{ and } \ \   
   a_{\ell}=  N^{\sigma}_{\ell}(\alpha^{-1})a_0= N^{\sigma}_{\ell}(\alpha^{-1}) N_{n}^{\sigma}(\alpha). $$ 
From (\ref{Eq_formulae}), we have $ N^{\sigma}_{\ell}(\alpha^{-1}) N_{n}^{\sigma}(\alpha)=  N_{n-\ell}^{\sigma}(\sigma^{\ell}(\alpha)).$
Then 
$$ N_{n}^{\sigma}(\alpha)  = a_0\ \ \text{ and } \ \   
   a_{\ell}= N_{n-\ell}^{\sigma}(\sigma^{\ell}(\alpha)). $$ 
 It follows that  
$$(a_0,a_{\ell} )\star \left(N_{n}^{\sigma}(\alpha^{-1}), N_{n-\ell}^{\sigma}(\sigma^{\ell}(\alpha^{-1})) \right)= (1,1) .$$
 By the second statement of Theorem \ref{Th.1}, we conclude that  $ (a_0, a_{\ell}) \sim_{(n,\ell,\sigma)}(1,1) ,$ and the result holds.  

   
\end{itemize}
 
 \end{proof}

 




  In the following result, we describe the associated polynomials of possible skew $(\ell,\sigma)$-trinomial codes based on the $(n,\ell,\sigma)$-equivalence relation defined above in Definition \ref{Def_Iso}.
\begin{theorem} \label{Equiv_2}
Let $n, \ell$ be two integers such that $0<\ell<n,$ let $\xi$ be  a primitive element of $\Fq$, and $\sigma$ be an automorphism of $ \Fq.$  Set  $ 
 d:=\gcd \left(\frac{q-1}{ d_0}, \frac{q-1}{ d_{\ell}}\right)$  and $ d_i:=\gcd([n-i]_r,q-1)$, for $ \ i\in \{0,\ell\}$.  Then the class of skew  $(\ell,\sigma)$-trinomial  codes associated with $x^n-a_{\ell}x^{\ell}-a_0$ is $(n,\ell,\sigma)$-equivalent (Hamming equivalence) to the class of skew  $(\ell,\sigma)$-trinomial  codes associated with  $x^n-\xi^{j} x^{\ell}- \xi^{i+h[n]_r},$ for some $i\in\{0,1,\ldots, d_0-1\}$,  $j\in \{0,1,\ldots, d_{\ell}-1\}$  and $ h\in \{0,\ldots, d-1\}$.  
\end{theorem}

\begin{proof}
\begin{enumerate}
    \item  If $ d = 1 $, then the cyclic group $ H_{\ell,\sigma} $ generated by $ \left( N_n(\xi), N_{n-\ell}(\sigma^{\ell}(\xi))\right) $ is isomorphic to the group $ \langle  N_n(\xi) \rangle \times \langle N_{n-\ell}(\sigma^{\ell}(\xi)) \rangle $ and has order $ \frac{(q-1)^2}{d_0 d_{\ell}} $. By Theorem \ref{Th.1}, the number of $ (n, \ell) $-equivalence classes is $ d_0 d_{\ell} $. Therefore, we can partition $ \Fq^* \times \Fq^* $ as follows

$$
\Fq^* \times \Fq^* =  \bigcup_{i=0}^{d_0-1} \bigcup_{j=0}^{d_{\ell}-1} (\xi^i, \xi^j) H_{\ell,\sigma}.
$$
Then any pair $(a_0,a_{\ell})$ is $(n,\ell,\sigma)$-equivalent to one of the pairs   $(\xi^i, \xi^j), $ for  $i=0,1,\ldots, d_0-1$, and  $j=0,1,\ldots, d_{\ell}-1$.
 \item 
 If $d\neq 1, $ the number of  $(n,\ell,\sigma)$-equivalence classes is $ d d_0 d_{\ell} ,$ and so we partition $ \Fq^* \times \Fq^* $ as follows:
$$  \Fq^{*}\times \Fq^{*} = \bigcup_{h=0}^{d-1}  \bigcup_{i=0}^{d_0-1} \bigcup_{j=0}^{d_{\ell}-1} (\xi^{i+h[n]_r} ,\xi^{j}) H_{\ell,\sigma}, $$
which implies the second statement, similarly to the first case.
\end{enumerate}
  
  
 
\end{proof} 


In the following proposition, we deduce additional characterizations of the $(n,\ell,\sigma)$-equivalence between skew $(\ell,\sigma)$-trinomial codes, linking this $(n,\ell,\sigma)$-equivalence to that of skew  constacyclic codes \cite{Ouazzou2025,Boulanouar2021,Lobillo2025}. 


%\hl{What about the proof of this proposition?}
%\textcolor{blue}{ The proof of this proposition is the same as the above one where we need to use the $\Fq-$algebra isometries  between $ \Fq[x; \sigma]/\langle x^{n-i}-b_i$ and $\Fq[x; \sigma]/\langle x^{n-i}-a_i,$ for $i=0,1$. We studied this case before in our last paper about the equivalence of skew constacyclic codes.  So I think we don't need to repeat  the proof here.  }
\begin{proposition}\label{Eq_Consta}
Let $ (a_0, a_{\ell}) $ and $ (b_0, b_{\ell}) $ be elements of $ \Fq^{*} \times \Fq^{*} $ such that $ (a_0, a_{\ell}) \sim_{(n, \ell,\sigma)} (b_0, b_{\ell}) $. Then:
\begin{enumerate}
    \item For each $ i \in \{0, \ell\} $, $ a_i^{-1} b_i \in \langle N_{n-i}^{\sigma}(\sigma^i(\xi)) \rangle $, where $ \xi $ is a primitive element of $ \Fq^* $.
    \item For each $ i \in \{0, \ell\} $, $ (a_i^{-1} b_i)^{d_i} = 1 $, where $ d_i := \frac{q-1}{\gcd([n-i]_r, q-1)} $.
    \item For each $ i \in \{0, \ell\} $, $ a_i \sim_{(n-i,\sigma)} b_i $.
    \item For each $ i \in \{0, \ell\} $ the class of skew  $(a_i,\sigma)$-constacyclic codes of length $n-i$ is Hamming equivalent to the class of skew   $(b_i,\sigma)$-constacyclic codes of length $n-i$  over $\Fq.$
\end{enumerate}
\end{proposition}
\begin{proof}
\begin{enumerate}
\item  As $ (a_0, a_{\ell}) \sim_{(n, \ell,\sigma} (b_0, b_{\ell})$, by the second  statement of Theorem \ref{Th.1}, 
there exists $\alpha\in \Fq^*$ such that $$ a_0 N_n^{\sigma}(\alpha)  = b_0   \quad  \text{and} \quad   a_{\ell} N^{\sigma}_{n-\ell}(\sigma^{\ell}(\alpha))=  b_{\ell}. $$
This gives
$$ a_0^{-1}  b_0 = N_n^{\sigma}(\alpha)  \quad  \text{and} \quad   a_{\ell}^{-1}   b_{\ell}=  N^{\sigma}_{n-\ell}(\sigma^{\ell}(\alpha)) . $$
Which means that  
$$ a_0^{-1}  b_0 \in   \langle  N_n^{\sigma}(\xi) \rangle  \quad  \text{and} \quad   a_{\ell}^{-1}   b_{\ell} \in  
\langle   N^{\sigma}_{n-\ell}(\sigma^{\ell}(\xi))  \rangle. $$

\item Let   $i\in \{0,\ell\}$ then the order of $ \langle   N^{\sigma}_{n-i}(\sigma^{i}(\xi))  \rangle =
\langle   \xi^{p^{ri}[n]_r}  \rangle $ is  $$d_i=\dfrac{q-1}{\gcd(p^{ri}[n-i]_r, q-1)}= \dfrac{q-1}{\gcd([n-i]_r, q-1)} , \text{ since $ \gcd(p^{ri},q-1)=1.$ }$$  It follows that  $(b_i a_i^{-1})^{d_i} =1 ,\ i\in \{0,\ell\}.$
\item By the second  statement of Theorem \ref{Th.1}, 
there exists $\alpha\in \Fq^*$ such that $$ a_0 N_n^{\sigma}(\alpha)  = b_0   \quad  \text{and} \quad   a_{\ell} N^{\sigma}_{n-\ell}(\sigma^{\ell}(\alpha))=  b_{\ell}. $$  As in the proof of \cite[Theorem 6]{Ouazzou2025}. For $i\in \{0,\ell\},$ we verify
 that the  map  $\varphi_i$ given by 
\begin{equation}
  	\begin{array}{cccc}
  	\varphi_{i}:& \mathbb{F}_q[x] /\langle x^{n-i}-b_0 \rangle  &\longrightarrow & \mathbb{F}_q[x] /\langle x^{n-i}-a_0 \rangle, \\
  	& f(x) & \longmapsto &  f(\sigma^i(\alpha) x). 
  	\end{array}
  	\end{equation} 
    is an $\Fq$-morphism isometry with respect to the Hamming  metric.
    \item Follows straight-forwardly from the third statement.
\end{enumerate}
\end{proof}

\subsection{Rank equivalence  of  skew  trinomial codes}
From the beginning of this section, we know that an $\mathbb{F}_{q}$-morphism is an isometry of the form $\varphi_{\alpha}$, where $\alpha \in \mathbb{F}_{q'}$, preserves the rank distance. Not really, we are just checking maps of the form $\varphi_\alpha$ for simplicity. Moreover, $\varphi_\alpha$ preserves the distance, not $\alpha$ itself. Recall here that the rank of $v=(v_0,v_1,\ldots,v_{n-1})\in \Fq^n$ is the dimension of the $\mathbb{F}_{q'}$-vector space $ \langle v_0,v_1,\ldots,v_{n-1}\rangle_{\mathbb{F}_{q'}}.$ We  now give the definition of rank equivalence in this case, along with a complete characterization, as in the case of Hamming equivalence.


\begin{definition}\label{DefFqSigma}
Let $a_0, a_{\ell}, b_0, b_{\ell}$ be nonzero elements of $\mathbb{F}_q$, and let $\ell$ be an integer such that $0 < \ell < n$. We say that $(a_0, a_{\ell})$ and $(b_0, b_{\ell})$ are {$( n,\ell,\sigma, \mathbb{F}_{q'})$-equivalent} in $\mathbb{F}_q^* \times \mathbb{F}_q^*$ with respect to the rank metric, and we  denote
$$ (a_0, a_{\ell}) \sim_{( n,\ell,\sigma, \mathbb{F}_{q'})} (b_0, b_{\ell}),$$
if there exists an $\alpha \in \mathbb{F}_{q'}^* $ such that the map

  \begin{equation}
  	\begin{array}{cccc}
  	\varphi_{\alpha}:& \mathbb{F}_q[x;\sigma] /\langle x^n-b_{\ell}x^{\ell}-b_0 \rangle  &\longrightarrow & \mathbb{F}_q[x;\sigma] /\langle x^n-a_{\ell} x^{\ell}-a_0 \rangle, \\ 
  	& f(x) & \longmapsto &  f(\alpha x),
  	\end{array}
  	\end{equation} 
is an  $\mathbb{F}_q$-morphism isometry with respect to the rank metric. 
\end{definition}

We give now the following result regarding the rank equivalence between two classes of skew trinomial codes. Because the proof is analogous to the one of Theorem \ref{Th.1}---with the small change that we require $\alpha \in \mathbb{F}_{q'}\subseteq \Fq$---we omit it here. 
\begin{theorem}\label{Th.RankEquiv}
Let $0<\ell<n$ be an integer,  $(a_0, a_{\ell})$ and $(b_0, b_{\ell})$ be elements of $\mathbb{F}_q^{*} \times \mathbb{F}_q^{*}$,  and   $\zeta$ be  a primitive element of $\mathbb{F}_{q'}$. 
Then the following statements are equivalent:

\begin{enumerate}
    \item   $ (a_0, a_{\ell})  \sim_{( n,\ell,\sigma, \mathbb{F}_{q'})} (b_0, b_{\ell}) .$ 
    \item There exists $\alpha \in  \mathbb{F}_{q'}^*$ such that the map

  \begin{equation}
  	\begin{array}{cccc}
  	\varphi_{\alpha}:& \mathbb{F}_q[x;\sigma] /\langle x^n-b_{\ell}x^{\ell}-b_0 \rangle  &\longrightarrow & \mathbb{F}_q[x;\sigma] /\langle x^n-a_{\ell} x^{\ell}-a_0 \rangle, \\ 
  	& f(x) & \longmapsto &  f(\alpha x),
  	\end{array}
  	\end{equation} 
    is an  $\mathbb{F}_q$-morphism isometry with respect to the rank metric.
       
     \item  There exists $\alpha\in  \mathbb{F}_{q'}^* $ such that $ (a_0,a_{\ell})\star (N_n^{\sigma}(\alpha), N_{n-\ell}^{\sigma}(\sigma^{\ell}(\alpha)))= (b_0,b_{\ell}) $.

      \item  $ (a_0, a_{\ell})^{-1} \star(b_0,b_{\ell})\in H_{\ell,  \mathbb{F}_{q'}},$ where   $ H_{\ell, \mathbb{F}_{q'}}$ is the cyclic subgroup of $  \mathbb{F}_{q'}^*\times  \mathbb{F}_{q'}^*$ generated by $( N_n^{\sigma}(\zeta), N_{n-\ell}^{\sigma}(\sigma^{\ell}(\zeta)) ) .$
\end{enumerate}
The equivalence between (1) and (4) implies that the number of $( n,\ell,\sigma, \mathbb{F}_{q'})$-equivalence classes is $$ 
N_{( n,\ell,\sigma, \mathbb{F}_{q'})}=\dfrac{ (q-1)^2}{ \lcm(\frac{q'-1}{ \gcd([n]_r,q'-1)}, \frac{q'-1}{ \gcd([n-\ell]_r,q'-1)})} .$$
\end{theorem}
%\begin{proof}
%\end{proof}



In the case where we calculate the rank over the fixed subfield $ \mathbb{F}_q^{\sigma} = \mathbb{F}_{q_0}, $ we obtain the following characterizations of rank equivalence. In this case, we provide more characterizations compared to Theorem~\ref{Th.RankEquiv}, using the right roots of our skew trinomial polynomials.

\begin{corollary} \label{Th.1FqSigma} 
Let $0<\ell<n$ be an integer,  $(a_0, a_{\ell})$ and $(b_0, b_{\ell})$ be elements of $\mathbb{F}_q^{*} \times \mathbb{F}_q^{*}$,  and  
$\zeta$ be  a primitive element of $\mathbb{F}_{q_0}:=\Fq^{\sigma}$. 
Then the following statements are equivalent:
\begin{enumerate}
    \item   $ (a_0, a_{\ell})  \sim_{( n,\ell, \Fq^{\sigma})} (b_0, b_{\ell}) ,$ i.e.,  $(a_0, a_{\ell}) $ and $ (b_0, b_{\ell})$ are rank equivalent. 
    \item There exists an $\alpha \in (\mathbb{F}_q^{\sigma})^*$ such that the map

  \begin{equation}
  	\begin{array}{cccc}
  	\varphi_{\alpha}:& \mathbb{F}_q[x;\sigma] /\langle x^n-b_{\ell}x^{\ell}-b_0 \rangle  &\longrightarrow & \mathbb{F}_q[x;\sigma] /\langle x^n-a_{\ell} x^{\ell}-a_0 \rangle, \\ 
  	& f(x) & \longmapsto &  f(\alpha x),
  	\end{array}
  	\end{equation} 
    
is an  $\mathbb{F}_q$-morphism isometry with respect to the rank metric.

    \item  There exists $\alpha\in (\Fq^{\sigma})^* $ that is  a common right  root of the polynomials  $\displaystyle a_i x^{n-i } - b_i \in \Fq[x;\sigma], i \in \{0,\ell\}. $ 

    \item   There exists $\alpha\in (\Fq^{\sigma})^* $ that is a right root of  $\gcrd(a_0 x^n - b_0, a_{\ell} x^{n - \ell} - b_{\ell}) \in \Fq[x;\sigma].$
       \item There exists $\alpha\in (\Fq^{\sigma})^* $ that is a right root of $\gcrd(x^n - b_0 a_0^{-1}, x^{n - \ell} - b_{\ell} a_{\ell}^{-1})\in \Fq[x;\sigma] .$ 
     \item  There exists $\alpha\in (\mathbb{F}_q^{\sigma})^* $ such that $ (a_0,a_{\ell})\star ( \alpha^n,   \alpha^{n-\ell})= (b_0,b_{\ell}) $.

     
      \item  $ (a_0, a_{\ell})^{-1} \star(b_0,b_{\ell})\in H_{\ell,\Fq^{\sigma}},$ where   $ H_{\ell,\Fq^{\sigma}}$ is the cyclic subgroup of $ (\Fq^{\sigma})^*\times (\Fq^{\sigma})^*$ generated by $(\zeta^n,\zeta^{n-\ell}) .$
\end{enumerate}
The equivalence between (1) and (7) implies that the number of $(n,\ell,\Fq^{\sigma})$-equivalence classes is $$ 
N_{(n,\ell,\Fq^{\sigma})}=\dfrac{ (q-1)^2}{ \lcm(\frac{q_0-1}{ \gcd(n,q_0-1)}, \frac{q_0-1}{ \gcd(n-\ell,q_0-1)})} .$$
\end{corollary}
\begin{proof}
     See Appendix.
\end{proof}


In the following result we give characterizations of rank equivalence for the class of  $(\ell,\sigma)$-trinomial codes  associated with  $x^n - a_{\ell}x^{\ell} - a_0$ to the class of $(\ell,\sigma)$-trinomial codes  associated with  $x^n - x^{\ell} - 1.$ The proof is similar to the case of Hamming equivalence (Corollary \ref{Equiv_1}), so we omit it.


\begin{corollary} \label{Caract_(1,1)}
Let $\sigma$ be and automorphism of $ \Fq,$ $\ell$ be an integer such that  $0<\ell<n$,  and  $(a_0, a_{\ell})$ be an element of $\mathbb{F}_q^{*} \times \mathbb{F}_q^{*}$. Then the following statements are equivalent.
    \begin{enumerate}
        \item  The class of skew  $(\ell,\sigma)$-trinomial codes  associated with  $x^n - a_{\ell}x^{\ell} - a_0$ is  $(n,\ell,\sigma, \mathbb{F}_{q'})$-equivalent (rank equivalence) to the class of   skew  $(\ell,\sigma)$-trinomial codes  associated with $x^n - x^{\ell} - 1.$
        \item  There exists $ \alpha \in \mathbb{F}_{q'}^*  $ such that $(a_0, a_{\ell})\star\left(N^{\sigma}_{n}(\alpha), N^{\sigma}_{n-\ell}(\sigma^{\ell}(\alpha)) \right)= (1, 1)$. 
        \item  There exists  an $(n,\sigma)$-th root  $ \alpha \in \mathbb{F}_{q'}^* $  of $a_0$  such that 
        $   a_0 =  N^{\sigma}_{\ell}(\alpha) a_\ell.$ 
        \end{enumerate}
\end{corollary}



 In the following result, we describe the associated polynomials of possible skew $(\ell,\sigma)$-trinomial codes based on the  $( n,\ell,\sigma, \mathbb{F}_{q'})$-equivalence relation defined above in Definition \ref{DefFqSigma}.
\begin{theorem}\label{RankClasses}
Let $n, \ell$ be two integers such that $0<\ell<n,$ let $\xi$ be  a primitive element of $\Fq$, $ \displaystyle{ \zeta:=\xi^{\frac{q-1}{q'-1}}}$  a primitive element of $\mathbb F_{q'}$ and $\sigma$ be an automorphism of $ \Fq.$  Set  $n':=\frac{q-1}{q'-1} ,$
  $  d:=\gcd \left(\frac{q'-1}{ d_0}, \frac{q'-1}{ d_{\ell}}\right)$  and $ d_i:=\gcd([n-i]_r,q'-1)$, for $ \ i\in \{0,\ell\}$. Then   the class of skew  $(\ell,\sigma)$-trinomial  codes associated with $x^n-a_{\ell}x^{\ell}-a_0$ is $(n,\sigma,\mathbb{F}_{q'})$-equivalent (rank equivalence) to the class of skew  $(\ell,\sigma)$-trinomial  codes associated with  $x^n-\xi^{j} x^{\ell}- \xi^{i+hn'},$ for some $i\in\{0,1,\ldots, d_0-1\}$,  $j\in \{0,1,\ldots, d_{\ell}-1\}$, and $ h\in \{0,\ldots, {n'}^2 d-1\}.$
  \end{theorem}
  \begin{proof}
      The proof here is the same as that one of Theorem \ref{Equiv_2}, the unique difference between them is in the number of equivalence classes. For the complete proof  see Appendix.
  \end{proof}

In the following proposition, we link the rank equivalence of skew trinomial codes to the one of skew constacyclic codes. The proof is again similar to the Hamming case (Proposition \ref{Eq_Consta}), details can be found in the Appendix.
\begin{proposition}\label{RankEquiv_Consta}
Let $ (a_0, a_{\ell}) $ and $ (b_0, b_{\ell}) $ be elements of $ \Fq^{*} \times \Fq^{*} $ such that $ (a_0, a_{\ell}) \sim_{(n, \ell, \mathbb{F}_{q'})} (b_0, b_{\ell}) $. Then:
\begin{enumerate}
    \item There exists $\alpha\in  \mathbb{F}_{q'}^{*} ,$ such that for each $ i \in \{0, \ell\} $, $ a_i^{-1} b_i \in
    \langle N_{n-i}^{\sigma}(\sigma^i(\alpha)) \rangle.$
    \item For each $ i \in \{0, \ell\} $, $ (a_i^{-1} b_i)^{d_i} = 1 $, where $ d_i := \dfrac{q'-1}{\gcd([n-i]_r, q'-1)} $.
    \item For each $ i \in \{0, \ell\} $, the map 
    \begin{equation*}
  	\begin{array}{cccc}
  	\varphi_{i}:& \mathbb{F}_q[x;\sigma] /\langle x^{n-i}-b_i\rangle  &\longrightarrow & \mathbb{F}_q[x;\sigma] /\langle
    x^{n-i}-b_i\rangle, \\ 
  	& f(x)& \longmapsto &  f(\sigma^i(\alpha) x)  
  	\end{array}
  	\end{equation*} 
is an $\mathbb{F}_q$-morphism isometry  with respect to the rank  distance.
    \item For each $ i \in \{0, \ell\} $ the class of skew  $(a_i,\sigma)$-constacyclic codes of length $n-i$ is { rank equivalent} to the class of skew   $(b_i,\sigma)$-constacyclic codes of length $n-i$  over $\Fq.$ 
\end{enumerate}
\end{proposition}
%\begin{proof}
%    It is similar to the proof of Proposition \ref{Eq_Consta} with a small difference. For more details see the Appendix section.
%\end{proof}
  
  
 


 
\begin{remark}\label{ConstacyclicFq_sigma}
In the case where we calculate the rank over the fixed subfield $ \mathbb{F}_q^{\sigma} = \mathbb{F}_{q_0}$ (i.e., we take $\alpha \in \Fq^{\sigma}$ such that $ \varphi_{\alpha}(x)= \alpha x$ is an $\Fq$-morphism isometry with respect to the rank metric (Definition \ref{DefFqSigma})), we have $ N_n(\sigma^i(\alpha))= \alpha^n$ for any $i$, since $ \alpha$ will be fixed by any power of $\sigma.$ In the case described here we simply get that $[n-i]_r$ is equal to $n-i.$ For more details see Appendix Proposition \ref{PropFq_sigma}.  
\end{remark}




\subsection{Generalization to skew polycyclic codes}
 In this subsection we generalize the results on {Hamming and rank}  equivalence to general skew  polycyclic codes. We start with the generalized definition of {Hamming} equivalence.

 %for which we will use the notation $\vec{a}(x) := a_0 + a_1 x + \cdots + a_{n-1} x^{n-1}$. 
 %We use the notation  $\vec{a}(x)$ to distinguish between  the  equivalence of right polycyclic codes and left polycyclic codes.
 
 %\textcolor{purple}{How so? I thought we always just consider right polycyclic codes.}
 %\hlc[green]{ I think this sentence has no sense. I had in mind to consider left polycyclic codes where I used in pervious works $\lleft{a} $ to indicate their associated vector. So I removed that sentence now.}
 
 %\hl{Why not use the more common notation $a(x)= a_0 + a_1 x + \cdots + a_{n-1} x^{n-1}$?}
 %\textcolor{blue}{ I used this notation to give difference between equivalence of right polycyclic codes and left polycyclic codes. I plan to focus on a project where the equivalence between two different classes ( left and right ) will be studied.}
 \begin{definition}
Let $ \vec{a} = (a_0, a_1, \ldots, a_{n-1}) $ and $ \vec{b} = (b_0, b_1, \ldots, b_{n-1}) $ be elements in $\mathbb{F}_q^n .$
We say that $ \vec{a} $ and $ \vec{b} $ are {Hamming $ (n,\sigma)$-equivalent}, and we denote this by 
$$ \vec{a} \sim_{(n,\sigma)} \vec{b}, $$ 
if there exists \( \alpha \in \mathbb{F}_q^* \) such that the following map

  \begin{equation}
  	\begin{array}{cccc}
  	\varphi_{\alpha}:& \mathbb{F}_q[x;\sigma] /\langle x^n-\vec{b}(x) \rangle  &\longrightarrow & \mathbb{F}_q[x;\sigma] /\langle x^n-\vec{a}(x) \rangle, \\ 
  	& f(x) & \longmapsto &  f(\alpha x),
  	\end{array}
  	\end{equation} 
is an $\mathbb{F}_q$-morphism isometry with respect to the  Hamming distance.    Moreover, we can easily verify that ``$ \sim_{(n,\sigma)}$'' is an equivalence relation.
 \end{definition}

\begin{remark}
 \begin{enumerate}
     \item  If $ \vec{a} = (\lambda, 0, \ldots, 0) $ and $ \vec{b} = (\mu, 0, \ldots, 0) $, then we recover the case of $(n,\sigma)$-equivalence for skew constacyclic codes studied in \cite{Ouazzou2025}.
     \item  If $ \vec{a} = (a_0, 0, \ldots,a_{\ell},0,\ldots, 0) $ and  $ \vec{b} = (b_0, 0, \ldots,b_{\ell},0,\ldots, 0) $, then we recover the case of $(n,\ell,\sigma)$-equivalence for  skew  $(\ell,\sigma)$-trinomial  codes studied in Section 3.
     \item If $\sigma=\id$, we recover  the case of $n$-equivalence for   polycyclic   codes studied in \cite{Equiv2025}.
     \end{enumerate}
 \end{remark}

 In the following lemma, we show that this notion of equivalence automatically implies that the vectors $\vec{a}$ and $\vec{b}$ have zero entries in exactly the same positions. This implies that skew  $(\ell,\sigma)$-trinomial code families can only be equivalent to other skew $(\ell,\sigma)$-trinomial code families.
 
 \begin{lemma}
    Let $ \vec{a} = (a_0, a_1, \ldots, a_{n-1}) $ and $ \vec{b} = (b_0, b_1, \ldots, b_{n-1}) $ be elements in $\mathbb{F}_q^n $ such that  $ \vec{a} \sim_{(n,\sigma)}\vec{b} $. Then $ a_i \neq 0 $ if and only if $ b_i \neq 0 $, for any $ 0\leq  i\leq n-1,$  and so $\vec{a}$ and $\vec{b}$ have the same Hamming weight.
 \end{lemma}
\begin{proof}
Suppose that $ \vec{a} \sim_n \vec{b},$ then there is $\alpha \in \Fq^*$ such that $ \varphi_{\alpha}$ is an $\mathbb{F}_q$-morphism isometry between $ \Fq[x;\sigma] /\langle x^n-\vec{b}(x) \rangle $  and $ \Fq[x;\sigma] /\langle x^n-\vec{a}(x) \rangle$. Then, as in the proof of [Theorem \ref{Th.1}, (1) $\Rightarrow$ (2)], we obtain that 
$$ b_i = N_{n-i}^{\sigma}(\alpha) a_i, \quad \forall i=0, \ldots, n-1.$$
Hence the result holds.
    
\end{proof}

We now generalize Theorem \ref{Th.1} to the general skew polycyclic case: 

\begin{theorem}\label{Th_general}
Let $ \vec{a}=(a_0,a_1,\ldots, a_{n-1}), \vec{b}=(b_0,b_1,\ldots, b_{n-1})\in \mathbb{F}_q^{n}$ have non-zero entries in the same $m$ positions, i.e., $a_{i_j}$ and $b_{i_j}$ are non-zero  for $0\leq i_0<\dots<i_{m-1}\leq n-1$. Moreover, let $\xi $ be a primitive element of $\Fq$. Then the following statements are equivalent:
\begin{enumerate}
\item $ \vec{a} \sim_{(n,\sigma)} \vec{b}$,  i.e., $ \vec{a} $ and  $ \vec{b}$ are  Hamming $ (n,\sigma)$-equivalent. 

\item   There exists $\alpha\in \Fq^*$ such that 
 $$  a_{i_j} N^{\sigma}_{n-i_j}(\sigma^{i_j}(\alpha)) = b_{i_j}, \   \text{for  each $ j \in \{0,1,\ldots,m-1\}$,}$$   where the  $ a_{i_j}$'s and $  b_{i_j}$'s are the non-zero components of $\vec{a}$ and $\vec{b}.$
 
 \item  There exists $\alpha\in \Fq^*$ such that 
 $$ (a_{i_0}, a_{i_1},\ldots, a_{i_{m-1}}) \left( N^{\sigma}_{n-i_0}(\sigma^{i_0}(\alpha)), N^{\sigma}_{n-i_1}(\sigma^{i_1}(\alpha)) , \ldots, N^{\sigma}_{n-i_{m-1}}(\sigma^{i_{m-1}}(\alpha))\right)=  ( b_{i_0}, b_{i_1},\ldots, b_{i_{m-1}}),$$
 where $m$ is the number of the non-zero entries  of $\vec{a}.$
\item $ (a_{i_0}, a_{i_1},\ldots, a_{i_{m-1}})^{-1}\star ( b_{i_0}, b_{i_1},\ldots, b_{i_{m-1}}) \in H,$ where $H$ is the cyclic subgroup of $ (\Fq^*)^m$ generated by 
$  \left( N^{\sigma}_{n-i_0}(\sigma^{i_0}(\xi)), N^{\sigma}_{n-i_1}(\sigma^{i_1}(\xi)) , \ldots, N^{\sigma}_{n-i_{m-1}}(\sigma^{i_{m-1}}(\xi))\right),$  and   $j=0,\ldots,m-1.$
\end{enumerate}
In particular,the number of $(n,\sigma)$-equivalence classes is
$$ N_{(n,\sigma)}=\dfrac{ (q-1)^m}{ \lcm_{ j\in \{i_0,i_1,\ldots, i_{m-1} \} }\left(\frac{q-1}{ \gcd([n-j]_r,q-1)}\right)}. $$
\end{theorem}

\begin{proof}\;
\begin{itemize}
    \item[(1) $\Rightarrow$ (2):] 
    Suppose that $ \vec{a} \sim_{(n,\sigma)} \vec{b} $, then there is $\alpha \in \mathbb{F}_q^*$ such that

$$\varphi_{\alpha} : \Fq[x;\sigma] /\langle x^n - \vec{b}(x) \rangle \to \Fq[x;\sigma] /\langle x^n - \vec{a}(x) \rangle, \quad f(x) \mapsto f(\alpha x) $$
is an $\mathbb{F}_q$-vector space isometry.  It follows that

$$ 
\varphi_{\alpha}(x^{k}) = (\alpha x)^k = N^{\sigma}_k(\alpha) x^{k} \mod (x^n - \vec{a}(x)), \quad \forall  \ k = 0, 1, \ldots, n-1.
$$

As $\varphi_{\alpha}$ is an $\mathbb{F}_q$-algebra isometry and $ \varphi(x^n - \vec{b}(x)) = 0 \mod (x^n - \vec{a}(x))$, we have
\begin{equation}
\varphi_{\alpha}(x^n) =\varphi_{\alpha} (\vec{b}(x)) = b_0 + N_1^{\sigma}(\alpha) b_1 x + \ldots + N_{n-1}^{\sigma}(\alpha)b_{n-1} x^{n-1} \mod (x^n - \vec{a}(x)).
\end{equation}

On the one hand,
\begin{equation}
\varphi(x^n) = N^{\sigma}_n(\alpha) x^n = N^{\sigma}_n(\alpha) ( a_0 + a_1 x + \ldots + a_{n-1} x^{n-1}) \mod (x^n - \vec{a}(x)).
\end{equation}

Comparing term by term, we deduce that   
$$ a_i N_{n}^{\sigma}(\alpha) = N_i^{\sigma}(\alpha)b_i, \ \text{ for any }  i \in \{0, 1, \ldots, n-1\},$$ 
and so 
$$ a_i N_{n}^{\sigma}(\alpha)  N_i^{\sigma}(\alpha^{-1}) =b_i, \ \text{ for any }  i \in \{0, 1, \ldots, n-1\}.$$ 
On the other hand, for any $  i \in \{0, 1, \ldots, n-1\} , $ we have 
$$ \begin{array}{rl}
    N_n^{\sigma}(\alpha) N_{i}^{\sigma}(\alpha^{-1})  & =  N_{i+(n-i)}^{\sigma}(\alpha) N_{i}^{\sigma}(\alpha^{-1})\\
     &=   \sigma^{i}(N_{n-i}^{\sigma}(\alpha)) N^{\sigma}_{i}(\alpha) N_{i}^{\sigma}(\alpha^{-1})  \\
     &=  \sigma^{i}(N^{\sigma}_{n-i}(\alpha))\\
     & = N^{\sigma}_{n-i}(\sigma^{i}(\alpha)).
\end{array}
$$

Hence,
$$ a_i N_{n-i}^{\sigma}( \sigma^i(\alpha))  =b_i, \ \text{ for any }  i \in \{0, 1, \ldots, n-1\}.$$ 

As  $ a_{i_j}$'s and $  b_{i_j}$'s are the non-zeros components of $\vec{a}$ and $\vec{b},$  we have
$$ a_{i_j} N_{n-{i_j}}^{\sigma}(\sigma^{i_j}(\alpha)) = b_{i_j},\  \ j=0,\ldots,m-1.$$

 \item[(2) $\Rightarrow$ (3):]  is immediate. 

 \item[(3) $\Rightarrow$ (4):]
Suppose that there is $\alpha\in \Fq^{*}$ such that 
$$  (b_{i_0}, b_{i_1},\ldots, b_{i_{m-1}} )=  \left( N_{n-i_0}^{\sigma}(\sigma^{i_0}(\alpha)), 
N^{\sigma}_{n-i_1}(\sigma^{i_1}(\alpha)), \ldots, N^{\sigma}_{n-i_{m-1}}(\sigma^{i_{m-1}}(\alpha))\right)\star ( a_{i_0},  a_{i_1}, \ldots, a_{i_{m-1}}).$$ 
For $\alpha= \xi^h,$ we obtain 
$$  ( a_{i_0}, a_{i_1},\ldots, a_{i_{m-1}})^{-1}\star (b_{i_0}, b_{i_1},\ldots, b_{i_{m-1}} ) =   \left( N_{n-i_0}^{\sigma}(\sigma^{i_0}(\xi)), 
N^{\sigma}_{n-i_1}(\sigma^{i_1}(\xi)), \ldots, N^{\sigma}_{n-i_{m-1}}(\sigma^{i_{m-1}}(\xi))\right)^h .$$
It follows that $  ( a_{i_0},  a_{i_1}, \ldots, a_{i_{m-1}})^{-1}\star ( b_{i_0},  b_{i_1}, \ldots, b_{i_{m-1}}) $ belongs to the cyclic subgroup  $H$ of $ (\Fq^{*})^m$ generated by $ \left( N_{n-i_0}^{\sigma}(\sigma^{i_0}(\xi)), 
N^{\sigma}_{n-i_1}(\sigma^{i_1}(\xi)), \ldots, N^{\sigma}_{n-i_{m-1}}(\sigma^{i_{m-1}}(\xi))\right).$
   
    \item[(4) $\Rightarrow$ (1):]  
     Suppose that $ (a^{n-{i_0}}, a^{n-{i_1}} , \ldots, a^{n-i_{m-1}})^{-1}\star (b^{n-{i_0}}, b^{n-{i_1}} , \ldots, b^{n-i_{m-1}}) $ is an element of the cyclic subgroup  $H$ of $ (\Fq^{*})^m$ generated by
$$ \left( N_{n-i_0}^{\sigma}(\sigma^{i_0}(\xi)), 
N^{\sigma}_{n-i_1}(\sigma^{i_1}(\xi)), \ldots, N^{\sigma}_{n-i_{m-1}}(\sigma^{i_{m-1}}(\xi))\right) .$$
Then there exists an integer $ h$ such that 
     $$  ( a_{i_0},  a_{i_1}, \ldots, a_{i_{m-1}})^{-1}\star ( b_{i_0},  b_{i_1}, \ldots, b_{i_{m-1}}) =  \left( N_{n-i_0}^{\sigma}(\sigma^{i_0}(\xi)), 
N^{\sigma}_{n-i_1}(\sigma^{i_1}(\xi)), \ldots, N^{\sigma}_{n-i_{m-1}}(\sigma^{i_{m-1}}(\xi))\right)^h $$
    For $ \alpha= \xi^{h}, $ we obtain  that $  a_j N_{n-j}^{\sigma}(\sigma^j(\alpha))= b_j, $ for any $ j\in \{ i_0,i_1,\ldots, i_{m-1}\}.$
As in the proof of Theorem \ref{Th.1}, we verify that  $ \varphi_{\alpha}, $ defined as  
  \begin{equation}
        \begin{array}{cccc}
        \varphi_{\alpha} : & \Fq[x;\sigma]/\langle x^n - \vec{b}(x) \rangle,  & \longrightarrow & \Fq[x;\sigma] /\langle x^n - \vec{a}(x) \rangle, \\ 
        & f(x) & \longmapsto & f(\alpha x),
        \end{array}
    \end{equation} 
    is an $\Fq$-vector space  isometry with respect to the Hamming distance.
To verify that $\varphi_{\alpha}$ is an $\mathbb{F}_q$-morphism, the key idea is to use the fact that  
$$ 
a_{i_k} N^{\sigma}_{n-i_k}(\sigma^{i_k}(\alpha)) = b_{i_k}, \quad \forall k \in \{0,1,\ldots,m-1\}, \quad \text{and} \quad N_{i+j}^\sigma(\alpha) = N_i^\sigma(\alpha)\sigma^i(N_j^\sigma(\alpha)),
$$  
to  obtain 
$$ 
N_i^\sigma(\alpha)\sigma^i(N_{n-i+j}^\sigma(\alpha)) \sigma^j(a_{k}) = N^{\sigma}_{k+j}(\alpha) \sigma^{j}(b_{k}).
$$  

From this, we can conclude that  
$$ 
\varphi_{\alpha}(f(x)g(x)) = \varphi_{\alpha}(f(x))\varphi_{\alpha}(g(x)).
$$  

We gave full details of the computation in the proof of Theorem \ref{Th.1},  so we do  not repeat the same process here.
\end{itemize}

\noindent By the equivalence between (1) and (4), we deduce that the number of $(n,\sigma)$-equivalence classes  on $(\Fq^*)^m$ corresponds to the order of the group $ (\Fq^*)^m/H,$ which equals 
$$ N_{(n,\sigma)}=\dfrac{ (q-1)^m}{ \lcm_{ j\in \{i_0,i_1,\ldots, i_{m-1} \} }\left(\frac{q-1}{ \gcd(p^{r j}[n-j]_r,q-1)}\right)}=\dfrac{ (q-1)^m}{ \lcm_{ j\in \{i_0,i_1,\ldots, i_{m-1} \} }\left(\frac{q-1}{ \gcd([n-j]_r,q-1)}\right)}. $$ 


\end{proof}

Similarly to the case of skew $(\ell,\sigma)$-trinomial codes, Theorem \ref{Th_general} implies the following results regarding the equivalence of skew polycyclic codes. The proofs are analogous to the skew $(\ell,\sigma)$-trinomial case. 

\begin{corollary}\label{Equiv_General2}
As before let $ \vec{a}=(a_0,a_1,\ldots, a_{n-1})$ and $ \vec{b}=(b_0,b_1,\ldots, b_{n-1})$ be elements of $\mathbb{F}_q^{n}$  
of the same weight $m$  and denote by $ a_{i_j}$ and $  b_{i_j}$  the non-zero components of $\vec{a}$ and $\vec{b}$. 
\begin{enumerate}
    \item     The class of skew polycyclic  codes associated with $ x^n-\dsum{j=0}{m-1} a_{i_j} x^{i_j} $ is $(n,\sigma)$-equivalent (Hamming equivalence) to the class of  skew polycyclic  codes associated with the vector  $x^n-\dsum{j=0}{m-1}  x^{i_j}$  if and only if there exists $\alpha \in \Fq^*$ such that
    $$ ( a_{i_0}, a_{i_1},\ldots, a_{i_{m-1}}) 
    \left( N_{n-i_0}^{\sigma}(\sigma^{i_0}(\alpha)),  N^{\sigma}_{n-i_1}(\sigma^{i_1}(\alpha)), \ldots, N^{\sigma}_{n-i_{m-1}}(\sigma^{i_{m-1}}(\alpha))\right) = (1, 1,\ldots, 1 ).$$  
   

   
     \item    Let  $ d:= \gcd_{ j\in \{i_0,i_1,\ldots, i_{m-1} \} }\left(\frac{q-1}{ \gcd([n-j]_r,q-1)}\right) $, then the  class of  skew polycyclic  codes associated with  $ x^n-\dsum{j=0}{m-1} a_{i_j} x^{i_j} $ is $(n,\sigma)$equivalent (Hamming equivalence) to the class of  skew polycyclic  codes associated with the polynomial  $ x^n-\dsum{j=1}{m-1} \xi^{k_j} x^{i_j}- \xi^{k_0+h[n]_r} ,$ with $k_j=0,1,\ldots, \gcd(n-i_j,q-1) , \ \ h=0,\ldots, d-1.$ 
\end{enumerate}
\end{corollary} 

\hspace{0.3cm}

Now we return to the case of rank equivalence, and we start by the following definition. 
\begin{definition} Let $ \vec{a} = (a_0, a_1, \ldots, a_{n-1}) $ and $ \vec{b} = (b_0, b_1, \ldots, b_{n-1}) $ be elements in $\mathbb{F}_q^n .$ We say that $\vec{a}$ and $\vec{b}$ are {$( n,\sigma, \mathbb{F}_{q'})$-equivalent} in $\mathbb{F}_q^* \times \mathbb{F}_q^*$ with respect to the rank metric, and we  denote
$$ \vec{a} \sim_{( n,\sigma, \mathbb{F}_{q'})} \vec{b},$$
if there exists an $\alpha \in \mathbb{F}_{q'}^* $ such that the map
  \begin{equation}
  	\begin{array}{cccc}
  	\varphi_{\alpha}:& \mathbb{F}_q[x;\sigma] /\langle x^n-\vec{b}(x) \rangle  &\longrightarrow & \mathbb{F}_q[x;\sigma] /\langle x^n- x^n-\vec{a}(x) \rangle, \\ 
  	& f(x) & \longmapsto &  f(\alpha x),
  	\end{array}
  	\end{equation} 
is an  $\mathbb{F}_q$-morphism isometry with respect to the rank metric. 
\end{definition}

\begin{theorem}
Let $ \vec{a}=(a_0,a_1,\ldots, a_{n-1}), \vec{b}=(b_0,b_1,\ldots, b_{n-1})\in \mathbb{F}_q^{n}$ have non-zero entries in the same $m$ positions, i.e., $a_{i_j}$ and $b_{i_j}$ are non-zero  for $0\leq i_0<\dots<i_{m-1}\leq n-1$. Moreover, let $ \mathbb{F}_{q'}$ be a subfield of $\Fq$ and 
$\zeta $ its primitive element. Then the following statements are equivalent:

\begin{enumerate}
    \item $ \vec{a} \sim_{(n,\sigma,\mathbb{F}_{q'})} \vec{b} $ i.e, $\vec{a} $ and $ \vec{b}$ are rank equivalent.
    \item There exists an $\alpha \in \mathbb{F}_{q'}^*$ such that the map

  \begin{equation}
  	\begin{array}{cccc}
  	\varphi_{\alpha}:& \mathbb{F}_q[x;\sigma] /\langle x^n-\vec{b}(x) \rangle  &\longrightarrow & \mathbb{F}_q[x;\sigma] /\langle x^n-\vec{a}(x) \rangle, \\ 
  	& f(x) & \longmapsto &  f(\alpha x),
  	\end{array}
  	\end{equation} 
    
is an  $\mathbb{F}_q$-morphism isometry.

   
     \item  There exists $\alpha\in \mathbb{F}_{q'}^* $ such that 
     $$ (a_{i_0}, a_{i_1},\ldots, a_{i_{m-1}})\star ( N_{n-i_0}^{\sigma}(\sigma^{i_0}(\alpha)), N_{n-i_1}^{\sigma}(\sigma^{i_1}(\alpha)),\ldots, N_{n-i_{m-1}}^{\sigma}(\sigma^{i_{m-1}}(\alpha)) ) =   ( b_{i_0}, b_{i_1},\ldots, b_{i_{m-1}}).$$
     
\item  $ (a_{i_0}, a_{i_1},\ldots, a_{i_{m-1}})^{-1}\star( b_{i_0}, b_{i_1},\ldots, b_{i_{m-1}}) \in H_{n, \mathbb{F}_{q'}},$ where $H_{n,\mathbb{F}_{q'}}$ is the cyclic subgroup of $\mathbb{F}_{q'}^{* m}$ generated by 
$( N_{n-i_0}^{\sigma}(\sigma^{i_0}(\zeta)), N_{n-i_1}^{\sigma}(\sigma^{i_1}(\zeta)),\ldots, N_{n-i_{m-1}}^{\sigma}(\sigma^{i_{m-1}}(\zeta)) ).$
\end{enumerate}
The equivalence between (1) and (4) implies that the number of $(n,\sigma,\mathbb{F}_{q'})$-equivalence classes is 
$$ 
N_{(n,\sigma,\mathbb{F}_{q'})}= \dfrac{ (q-1)^m}{ \lcm_{ j\in \{i_0,i_1,\ldots, i_{m-1} \} }\left(\frac{q'-1}{ \gcd([n-j]_r,q'-1)}\right)}. $$ 
\end{theorem}

When the rank is calculated over the fixed subfield $\Fq^{\sigma}$, we obtain additional characterizations of the rank equivalence equivalence.
\begin{corollary}
Let $ \vec{a}=(a_0,a_1,\ldots, a_{n-1}), \vec{b}=(b_0,b_1,\ldots, b_{n-1})\in \mathbb{F}_q^{n}$ have non-zero entries in the same $m$ positions, i.e., $a_{i_j}$ and $b_{i_j}$ are non-zero  for $0\leq i_0<\dots<i_{m-1}\leq n-1$. Moreover, let $\zeta $ be a primitive element of $\Fq^{\sigma}$. Then the following statements are equivalent:

\begin{enumerate}
    \item $ \vec{a} \sim_{( n,\Fq^{\sigma})} \vec{b} .$
    \item There exists an $\alpha \in (\mathbb{F}_q^{\sigma})^*$ such that the map

  \begin{equation}
  	\begin{array}{cccc}
  	\varphi_{\alpha}:& \mathbb{F}_q[x;\sigma] /\langle x^n-\vec{b}(x) \rangle  &\longrightarrow & \mathbb{F}_q[x;\sigma] /\langle x^n-\vec{a}(x) \rangle, \\ 
  	& f(x) & \longmapsto &  f(\alpha x),
  	\end{array}
  	\end{equation} 
    
is an  $\mathbb{F}_q$-morphism isometry with respect to the rank metric.

    \item  There exists $\alpha\in (\Fq^{\sigma})^* $ that is  a common right  root of the polynomials  $\displaystyle a_i x^{n-i } - b_i \in \Fq[x;\sigma], i \in \{0,\ldots,m-1\}. $ 

    \item   There exists $\alpha\in (\Fq^{\sigma})^* $ that is a right root of  $\gcrd_{\{0 \leq j \leq m-1\}}( a_{i} x^{n - i} - b_{i}) \in \Fq[x;\sigma].$
    
     \item  There exists $\alpha\in (\mathbb{F}_q^{\sigma})^* $ such that 
     $$ (a_{i_0}, a_{i_1},\ldots, a_{i_{m-1}})\star(\alpha^{n-{i_0}}, \alpha^{n-{i_1}} , \ldots, \alpha^{n-{i_{m-1}}})=  ( b_{i_0}, b_{i_1},\ldots, b_{i_{m-1}}).$$
\item  $ (a_{i_0}, a_{i_1},\ldots, a_{i_{m-1}})^{-1}\star( b_{i_0}, b_{i_1},\ldots, b_{i_{m-1}}) \in H_{n,\Fq^{\sigma}},$ where $H_{n,\Fq^{\sigma}}$ is the cyclic subgroup of $ (\Fq^{\sigma})^{* m}$ generated by 
$ ( N_{n-i_0}^{\sigma}(\sigma^{i_0}(\zeta)), N_{n-i_1}^{\sigma}(\sigma^{i_1}(\zeta)),\ldots, N_{n-i_{m-1}}^{\sigma}(\sigma^{i_{m-1}}(\zeta)) ) .$
\end{enumerate}
The equivalence between (1) and (7) implies that the number of $(n,\Fq^{\sigma})$-equivalence classes is 
$$ 
N_{(n,\Fq^{\sigma})}= \dfrac{ (q-1)^m}{ \lcm_{ j\in \{i_0,i_1,\ldots, i_{m-1} \} }\left(\frac{q_0-1}{ \gcd(n-j,q_0-1)}\right)}. $$ 
\end{corollary}

Similarly to  the case of skew $(\ell,\sigma)$-trinomial codes, Theorem \ref{Th_general} implies the following results regarding the rank equivalence of skew polycyclic codes. The proofs are analogous to the skew $(\ell,\sigma)$-trinomial case in Theorem \ref{RankClasses}.

\begin{corollary}
As before let $ \vec{a}=(a_0,a_1,\ldots, a_{n-1})$ and $ \vec{b}=(b_0,b_1,\ldots, b_{n-1})$ be elements of $\mathbb{F}_q^{n}$  
of the same weight $m$  and denote by $ a_{i_j}$ and $  b_{i_j}$  the non-zeros components of $\vec{a}$ and $\vec{b}$. Let $ \mathbb{F}_{q'}$ be a subfield of $\Fq, $ and $n':=\frac{q-1}{q'-1} .$  Then,
\begin{enumerate}
    \item    The class of skew polycyclic  codes associated with   $ x^n-\dsum{j=0}{m-1} a_{i_j} x^{i_j} $ is $(n,\sigma,\mathbb{F}_{q'})$-equivalent (rank equivalence) to the class of  skew polycyclic  codes associated with the vector  $x^n-\dsum{j=0}{m-1}  x^{i_j}$  if and only if there exists $\alpha \in \mathbb{F}_{q'}^*$ such that
    $$ ( a_{i_0}, a_{i_1},\ldots, a_{i_{m-1}})\star( N_{n-i_0}^{\sigma}(\sigma^{i_0}(\alpha)), N_{n-i_1}^{\sigma}(\sigma^{i_1}(\alpha)),\ldots, N_{n-i_{m-1}}^{\sigma}(\sigma^{i_{m-1}}(\alpha)) ) = (1, 1,\ldots, 1 ).$$  
        
     \item  Let   $ d:= \gcd_{ j\in \{i_0,i_1,\ldots, i_{m-1} \} }\left(\frac{q'-1}{ \gcd([n-j]_r,q'-1)}\right)$, then the  class of  skew polycyclic  codes associated with  $ x^n-\dsum{j=1}{m-1} a_{i_j} x^{i_j}-a_0 $ is  $(n,\sigma,\mathbb{F}_{q'})$-equivalent (rank equivalence) to the class of  skew polycyclic  codes associated with the polynomial  $ x^n-\dsum{j=1}{m-1} \xi^{k_j} x^{i_j}- \xi^{k_0+h n'} ,$ with $k_j=0,1,\ldots, \gcd(n-i_j,q_0-1) , \ \ h=0,\ldots, n'^m d-1.$ 
\end{enumerate}
\end{corollary} 











\section{Examples for code equivalence}

In this section, we provide some examples related to the Hamming and rank equivalence studied in the previous section. In particular, we characterize and analyze the possible number of Hamming equivalence classes and the representative element of each class, for skew  trinomial  codes of length $n$  over $\mathbb{F}_4$, $\mathbb{F}_8$, and $\mathbb{F}_9$. For rank equivalence, we characterize the number and the representative elements of each equivalence class in the case of trinomial codes of length $n$ over $\mathbb{F}_{2^s}$. Similar methods can be used to construct more examples for different alphabets.
 
\subsection{Hamming metric}

\begin{example}[Skew $(\ell,\sigma)$-trinomial  codes  over $\mathbb{F}_4$] 

Let $q = 4$, and $\xi$ be a primitive element of $\mathbb{F}_4$. In this example, we are interested in skew $(\ell,\sigma)$-trinomials of length $n$ over $\mathbb{F}_4$. Since the degree of the extension field $\mathbb{F}_4/\mathbb{F}_2$ is $2$, we have $\text{Aut}(\mathbb{F}_4) = \{\text{id}, \tau\}$, where $\tau$ is the Frobenius automorphism of $\mathbb{F}_4$, i.e., $\tau(a) = a^2$ for all $a\in \mathbb{F}_4$. 

For $\sigma=\tau,$ we consider skew $(\ell,\sigma)$-trinomial  codes   associated with  a skew polynomial of the form $f(x) =x^{n}-a_{\ell}x^{\ell}-a_0\in \mathbb{F}_4[x,\sigma],$
 where $\ell$ is an integer such that $0<\ell <n$.  Let $\vec{a}(x):=   a_{\ell}x^{\ell}+a_0 ,$  the Hamming weight of $ \vec{a}(x) $ is  $2.$
 As $[i]_1= \dfrac{2^i-1}{2-1}=2^i-1,$  by  Theorem  \ref{Th.1}, the number of Hamming $(n,\ell, \sigma)$-equivalence classes is given by  
 
 $$N:=\dfrac{3^2}{ \lcm\left(\dfrac{3}{\gcd( [n]_1 ,3 )}, \dfrac{3}{\gcd( [n-\ell]_1,3)}  \right)}    
      =\dfrac{9}{ \lcm\left( \frac{3}{\gcd( 2^{n}-1 ,  2^2-1 )}, \frac{3}{\gcd( 2^{n-\ell}-1 ,  2^2-1 )}\right) }  
 $$ 
 Now we analyze the value of $N$  based on the parity of $n$ and $\ell$.
 \begin{enumerate}
\item \textbf{If $n$ and $\ell$ are odd,} then $ \gcd( 2^{n}-1 ,  2^2-1 )=2^{\gcd(n,2)}-1=1 $ and  $ \gcd( 2^{n-\ell}-1 ,  2^2-1 )= 3.$ 
Then the number of $(n,\sigma)-$equivalence classes is $N=3,$ and by Theorem \ref{Equiv_2}, we have 
     $$d_0=\gcd( [n]_1, 3)= 1,\ \ d_{\ell}= \gcd([n-\ell]_1, 3)= 3 \text{ and} \ \  d=\gcd\left( \frac{3}{d_0}, \frac{3}{d_{\ell}}\right)=1.$$
 It follows that each  skew $(\ell,\sigma)$-trinomial  code associated with a skew  polynomial $f$ of the form  $f =x^{n}- a_{\ell} x^{\ell}-a_0$ is equivalent to a skew $(\ell,\sigma)$-trinomial  code  associated with one of the following  skew polynomials: 
 $$f_{0}= x^{n} - x^{\ell} - 1 , \quad f_{1}= x^{n} - \xi x^{\ell} - 1 \quad  f_{2}= x^{n} - \xi^2 x^{\ell} - 1 .  $$
 
     \item \underline{\textbf{If $n$ and $\ell$ are even,}} 
     then $ \gcd( 2^{n}-1 ,  2^2-1 )=2^{\gcd(n,2)}-1=3 $ and  $ \gcd( 2^{n-\ell}-1 ,  2^2-1 )= 3.$ 
Then the number of $(n,\sigma)-$equivalence classes is $N=9.$ In this case the $(n,\sigma)$-equivalence relation has no advantage in the sense that equivalence classes are singleton sets and the number of polynomials to be considered to construct skew polycyclic codes is not reduced.  
%\hl{What exactly does that mean?}
%\textcolor{blue}{ That means that each equivalence classes class contains one element and so it hsn't no advantages in reducing the number of possible polynomials to construct skew polycyclic codes. }

  \item \underline{\textbf{If $n$  is odd and $\ell$ is even,}}   then
  $ \gcd( 2^{n}-1 ,  2^2-1 )=2^{\gcd(n,2)}-1=1 $ and
  $ \gcd( 2^{n-\ell}-1 ,  2^2-1 )= 1.$  Then the number of $(n,\sigma)-$equivalence classes is $N=3,$ 
and by Theorem \ref{Equiv_2}, we have 
     $$d_0=\gcd( [n]_1, 3)= 1,\ \ d_{\ell}= \gcd([n-\ell]_1, 3)= 1 \text{ and} \ \  d=\gcd\left( \frac{3}{d_0}, \frac{3}{d_{\ell}}\right)=3$$
    
 It follows that each skew $(\ell,\sigma)$-trinomial  code associated with a skew  polynomial $f$ of the form  $f =x^{n}- a_{\ell} x^{\ell}-a_0,$ is equivalent to a skew $(\ell,\sigma)$-trinomial  code  associated with one of the following  skew polynomials: 
$$f_{0}= x^{n} - x^{\ell} - 1 , \quad f_{1}= x^{n} -  x^{\ell} -   \xi^{[n]_1},  \quad  f_{2}= x^{n} -   x^{\ell} -  \xi^{2[n]_1}.  $$

 \item \underline{\textbf{If $n$  is even and  $\ell$ is odd,}} then,  as above, we obtain that  each skew $(\ell,\sigma)$-trinomial  code associated with a skew  polynomial $f$ of the form  $f =x^{n}- a_{\ell} x^{\ell}-a_0,$ is equivalent to a skew $(\ell,\sigma)$-trinomial  code  associated with one of the following  skew polynomials:
 $$f_{0}= x^{n} - x^{\ell} - 1 , \quad f_{1}= x^{n} - x^{\ell} - \xi,  \quad  f_{2}= x^{n} - x^{\ell} - \xi^2 .  $$
   
 \end{enumerate}
 \end{example}

\begin{example}[Skew $(\ell,\sigma)$-trinomial  codes  over $\mathbb{F}_{8}$] 

Let $q = 8$, and $\xi$ be a primitive element of $\mathbb{F}_8$. In this example, we are interested in skew $(\ell,\sigma)$-trinomial codes of length $n$ over $\mathbb{F}_8$. Since the degree of the extension field $\mathbb{F}_{8}/\mathbb{F}_2$ is 3, we have $\text{Aut}(\mathbb{F}_{8}) = \{\text{id}, \tau,  \tau^{2}\}$, where $\tau$ is the Frobenius automorphism of $\mathbb{F}_{8}$, i.e., $\tau(a) = a^2$ for all $a\in \mathbb{F}_{8}$. Then we consider a number of cases:
\begin{enumerate}
    \item \textbf{ Case 1: }  For $ \sigma=\tau,$ we consider skew $(\ell,\sigma)$-trinomial  codes     associated with  a skew polynomial of the form $f(x) =x^{n}-a_{\ell}x^{\ell}-a_0\in \mathbb{F}_{8}[x,\sigma],$ where $\ell $ is an integer such that $0<\ell <n$.  Let $\vec{a}(x):=   a_{\ell}x^{\ell}+a_0$  with Hamming weight   $2.$
As $ [i]_1=\dfrac{2^{i}-1}{2-1}=2^i-1$ is odd, by Theorem  \ref{Th.1}, the number of  Hamming $(n,\ell, \sigma)$-equivalence classes is given by  
 
 $$N:=\dfrac{7^2}{ \lcm\left(\frac{7}{\gcd( [n]_1 ,7 )}, \frac{7}{\gcd( [n-\ell]_1,7)}  \right)}    
      =\dfrac{49}{ \lcm\left( \frac{7}{\gcd( 2^n-1 ,2^3-1 )}, \frac{7}{ \gcd( 2^{n-\ell}-1 ,2^3-1 )}\right)} .
 $$
 
 
 \begin{enumerate}
\item \underline{\textbf{If $ \gcd(n, 3)=\gcd(n-\ell, 3)=1$,}}  then the number of $(n,\sigma)$-equivalence classes is   $  N=7.$ By Theorem \ref{Equiv_2}, we have 
     $$  d_0= \gcd([n]_1, 3)=1, \ \ \ d_{\ell}= \gcd( [n-\ell]_1, 7)= 1 \text{ and} \ \  d=\gcd\left( \frac{ 7}{d_0}, \frac{ 7}{d_{\ell}}\right)= 7.$$
 It follows that each  skew $(\ell,\sigma)$-trinomial  code associated with a skew  polynomial $f$ of the form  $f =x^{n}- a_{\ell} x^{\ell}-a_0,$ is equivalent to a  skew $(\ell,\sigma)$-trinomial  code  associated with one of the following  skew polynomials: 
  $$f_{i}(x)= x^{n} - x^{\ell} - \xi^{ i [n]_1 }  \ \ \text{for } i\in \{0,\ldots, 6\}.  $$ 
 
     \item \underline{\textbf{If  $  \gcd(n, 3)= \gcd(n-\ell, 3)=3 $,}} then the number of $(n,\sigma)$-equivalence classes is   $  N=49,$ and the $(n,\sigma)$-equivalence relation has no advantage.

   \item \underline{\textbf{If  $  \gcd(n, 3)=1$  and  $\gcd(n-\ell, 3)=3 $,}}  then  the number of equivalence classes is  $N=7.$ 
   In this case, $ d_{0}= \gcd([n]_1, 7)= 1,\ \ d_{\ell}= \gcd( [n-\ell]_1, 7)= 7 .$ 
   It follows that each  skew $(\ell,\sigma)$-trinomial  code associated with a skew  polynomial $f$ of the form  $f =x^{n}- a_{\ell} x^{\ell}-a_0$ is equivalent to a  skew $(\ell,\sigma)$-trinomial  code  associated with one of the following  skew polynomials: 
  $$f_{i}(x)= x^{n} - \xi^i x^{\ell} - 1 \ \ \text{for } i\in \{0,\ldots, 6\}.  $$ 
  \item  \textbf{\underline{If  $  \gcd(n, 3)=3$  and  $\gcd(n-\ell, 3)=1 $, } } then, as above, we obtain that each 
  each  skew $(\ell,\sigma)$-trinomial  code associated with a skew  polynomial $f$ of the form  $f =x^{n}- a_{\ell} x^{\ell}-a_0,$ is equivalent to a  skew $(\ell,\sigma)$-trinomial  code  associated with one of the following  skew polynomials: 
  $$f_{i}(x)= x^{n} -  x^{\ell} - \xi^i \ \ \text{for } i\in \{0,\ldots, 6\}.  $$ 

    \end{enumerate}

\item \textbf{Case 2:} For $\sigma=\tau^2,$  we consider skew $(\ell,\sigma)$-trinomial   codes   associated with  a skew polynomial of the form $f(x) =x^{n}-a_{\ell}x^{\ell}-a_0\in \mathbb{F}_{8}[x,\sigma]$, where $\ell $ is an integer such that $0<\ell <n$.  Let $\vec{a}(x):=   a_{\ell}x^{\ell}+a_0 $  with  Hamming weight 2.   By Theorem  \ref{Th.1}, the number of Hamming $(n,\ell, \sigma)$-equivalence classes is given by  
 $$N:=\dfrac{7^2}{ \lcm\left(\frac{7}{\gcd( [n]_2 ,7 )}, \frac{7}{\gcd( [n-\ell]_2,7)}  \right)}.
 $$
 In this case, we performed computations  using Magma software, where we calculated $\operatorname{gcd}([n]_2, 7)$ for $n$ ranging from $1$ to $100,000,000$. We observed two distinct cases: 
$$
\gcd([n]_2, 7) = \begin{cases}
 1 , & \text{if $\gcd(n,3)=1$ }\\
7, & \text{otherwise }\\
  \end{cases} 
$$     

 \begin{enumerate}

 \item \underline{\textbf{If $\gcd(n,3)= \gcd(n,\ell)=1$,}} then $\displaystyle{  \gcd([n]_2, 7)=  \gcd([n-\ell]_2, 7)= 1}$, and so the number of  $(n, \sigma)$-equivalence classes is $ 7$. So the possible trinomials are: 
     $$f_{i}(x)= x^{n} - x^{\ell} - \xi^{ i [n]_2 }  \ \ \text{for  }\  i\in \{0,\ldots, 6\}.  $$ 

    \item \underline{ \textbf{If $\gcd(n,3)= \gcd(n,\ell)=3$,}} then $\gcd([n]_2, 7)= \gcd([n-\ell]_2, 7) = 7$.  In this case the number of possible classes is $N=49,$ and the equivalence relation has no advantage in this case.
    \item \underline{\textbf{If  $\gcd(n,3)= 1 $ and $ \gcd(n,\ell)=3$,}}  then the number of equivalence classes is $N=7,$ and the possible trinomials are 
    $$   x^n-\xi^i x^{\ell}- 1,  \ \ \text{for  }\  i \in \{0,\ldots, 6\}. $$

    \item  \underline{\textbf{If  $  \gcd(n,3)= 3 $ and $ \gcd(n,\ell)=1$,}} then the number of equivalence classes is $N=7,$ and the possible trinomials are 
    $$x^n- x^{\ell}- \xi^i,  \ \ \text{for  }\  i \in \{0,\ldots, 6\}. $$
   

     \end{enumerate}
 \end{enumerate}
 \end{example}



 \begin{example} [Skew $(\ell,\sigma)$-trinomial  codes  over $\mathbb{F}_9$]
Let $q = 9$, and $\xi$ be a primitive element of $\mathbb{F}_9$. In this example, we are interested in skew $(\ell,\sigma)$-trinomial   codes of length $n$ over $\mathbb{F}_9$. Since the degree of the extension field $\mathbb{F}_9/\mathbb{F}_3$ is $2$, we have $\text{Aut}(\mathbb{F}_9) = \{\text{id}, \tau\}$, where $\tau$ is the Frobenius automorphism of $\mathbb{F}_9$, defined as $\tau(a) = a^3$ for all $a \in \mathbb{F}_9$. 

For $\sigma=\tau,$ we consider skew $(\ell,\sigma)$-trinomial  codes   associated with  a skew polynomial of the form $f(x) =x^{n}-a_{\ell}x^{\ell}-a_0\in \mathbb{F}_4[x,\sigma],$
 where $\ell $ is  an integer such that $0<\ell <n$.  Let $\vec{a}(x):=   a_{\ell}x^{\ell}+a_0 ,$  with Hamming weight 2.
 As $[i]_1= \dfrac{3^i-1}{3-1},$ by Theorem  \ref{Th.1}, the number of the Hamming $(n,\ell, \sigma)$-equivalence classes is given by  
 
 $$N:=\dfrac{8^2}{ \lcm\left(\dfrac{8}{\gcd( [n]_1 ,8 )}, \dfrac{8}{\gcd( [n-\ell]_1,8)}  \right)}    
 $$ 
  In this case, we performed computations  using Magma software, where we calculated $\operatorname{gcd}([n]_1, 8)$ for $n$ ranging from $1$ to $100,000,000$. We observed the following cases:
$$
\gcd([n]_1, 8) = \begin{cases}
 1 , & \text{if $n$ is odd.}\\
4, &  \text{if $n \equiv  2 \mod(4)$.}\\
8, &  \text{if $n \equiv  0 \mod(4)$.}\\
  \end{cases} 
$$ 

\begin{enumerate}
    \item  \underline{\textbf{If $n$ and $\ell$ are odd,}} then $\gcd \left( [n]_1 , 8\right) =\gcd \left( [n-\ell]_1 , 8\right)= 1$, and the number of $(n,\ell,\sigma)$-equivalence classes is $ N= 8.$  Then by Theorem \ref{Equiv_2}, we have 
     $$d_0=\gcd( [n]_1, 8)= 3,\ \ d_{\ell}= \gcd( 3^{\ell}[n-\ell]_1, 3)= 1 \text{ and} \ \  d=\gcd\left( \frac{8}{d_0}, \frac{8}{d_{\ell}}\right)=8.$$

     It follows that each  skew $(\ell,\sigma)$-trinomial  code associated with a skew  polynomial $f$ of the form  $f =x^{n}- a_{\ell} x^{\ell}-a_0$ is equivalent to a  skew $(\ell,\sigma)$-trinomial  code  associated with one of the following  skew polynomials: 
  $$f_{i}(x)= x^{n} - x^{\ell} - \xi^{ i [n]_1 }  \ \ \text{for}, i\in \{0,\ldots, 7\}.  $$

    


 \item  \underline{\textbf{If $n \equiv n-\ell \equiv 2 \pmod{4},$}} then $\gcd \left( [n]_1 , 8\right) =\gcd \left( [n-\ell]_1 , 8\right)= 4$, and the number of $(n,\ell,\sigma)$-equivalence classes is $ N= 16.$  Then by Theorem \ref{Equiv_2}, we have 
     $$d_0=\gcd( [n]_1, 8)= 4,\ \ d_{\ell}= \gcd( [n-\ell]_1, 3)= 4 \text{ and} \ \  d=\gcd\left( \frac{8}{d_0}, \frac{8}{d_{\ell}}\right)=4.$$

     It follows that each  skew $(\ell,\sigma)$-trinomial  code associated with a skew  polynomial $f$ of the form  $f =x^{n}- a_{\ell} x^{\ell}-a_0,$ is equivalent to a  skew $(\ell,\sigma)$-trinomial  code  associated with one of the following  skew polynomials: 
  $$f_{i}(x)= x^{n} - \xi^i x^{\ell} - \xi^{ j+h [n]_1 }  \ \ \text{for}, i,j,h\in \{0,\ldots, 3\}.  $$

\item \underline{\textbf{ If $n \equiv n-\ell \equiv 0 \pmod{4}$,}} then $\gcd \left( [n]_1 , 8\right) =\gcd \left( [n-\ell]_1 , 8\right)= 8$, and the number of $(n,\ell,\sigma)$-equivalence classes is $ N= 64.$ Hence  the equivalence relation has no advantage in this case.
\end{enumerate}
One can also discuss other cases like \textbf{ $ n$ is odd} and\textbf{ $ n-\ell\equiv 2 \mod(4)$} using a similar approach as above to determine the possible skew polynomials.
\end{example}








\subsection{Rank metric}


\begin{example}[Skew $(\ell,\sigma)$-trinomial codes over $ \mathbb{F}_{2^s}$]

Let $q = 2^s$, and $\xi$ be a primitive element of $\mathbb{F}_{2^s}$. In this example, we are interested in rank equivalence of  skew $(\ell,\sigma)$-trinomial codes of length $n$ over $\mathbb{F}_{2^s}$. The degree of the extension   $\mathbb{F}_{2^s}/\mathbb{F}_2$ is $s$, and  $\text{Aut}(\mathbb{F}_{2^s}) = \{\text{id}, \tau, \ldots, \tau^{s-1}\}$, where $\tau$ is the Frobenius automorphism of $\mathbb{F}_{2^s}$, i.e., $\tau(a) = a^2$ for all $a\in \mathbb{F}_{2^s}$. 

Let $\sigma=\tau^r,$ then the fixed subfield of $ \Fq$ by $ \sigma$ is $ \mathbb{F}_{2^{\gcd(r,s)}}.$
We consider skew $(\ell,\sigma)$-trinomial  codes   associated with  a skew polynomial of the form $f(x) =x^{n}-a_{\ell}x^{\ell}-a_0\in \mathbb{F}_{2^s}[x;\sigma],$ where $\ell$ is an integer such that $0<\ell <n$.  Let $\vec{a}(x):=   a_{\ell}x^{\ell}+a_0 .$ According to Corollary \ref{Th.1FqSigma}, the number of $( n,\ell, \Fq^{\sigma} )$-equivalence classes is given by  
 
 $$ N:=\dfrac{(2^s-1)^2}{ \lcm\left(\dfrac{2^{\gcd(r,s)}-1}{\gcd( n , 2^{\gcd(r,s)}-1)}, \dfrac{2^{\gcd(r,s)}-1}{\gcd( n-\ell, 2^{\gcd(r,s)}-1)}  \right)} . 
 $$ 
 Now we analyze the value of $N$  based on the parity of $n,\ell, r$ and $s.$
 \begin{enumerate}
\item \textbf{If $\gcd(r,s)=1$,} then $ N= (2^s-1)^2$ and each $( n,\ell, \Fq^{\sigma})$-equivalence  class contains one element.
\item \textbf{If $\gcd(r,s)= l\neq 1$ and $n,\ \ell$ are even,} then 

$$ N=\dfrac{(2^s-1)^2}{ \lcm\left(\dfrac{2^{\ell}-1}{1}, \dfrac{2^{\ell}-1}{1}  \right)}= \frac{(2^s-1)^2}{2^{\ell}-1}.$$
By Theorem \ref{RankClasses}, we have 
     $$d_0=\gcd( n, 2^{\ell}-1)= 1,\ \ d_{\ell}= \gcd( n-\ell, 2^{\ell}-1)= 1 \text{ and} \ \  d=\gcd\left( \frac{ 2^{\ell}-1}{d_0}, \frac{2^{\ell}-1}{d_{\ell}}\right)=2^{\ell}-1.$$
 It follows that each  skew $(\ell,\sigma)$-trinomial  code associated with a skew  polynomial $f$ of the form  $f =x^{n}- a_{\ell} x^{\ell}-a_0$ is rank equivalent to a skew $(\ell,\sigma)$-trinomial  code  associated with one of the following  skew polynomials: 
  $$f_i(x)= x^n- x^{\ell}- \xi^{\ell n'},$$
  for $\ell \in \{ 0,\ldots, d n'^2-1\},$ with $ n':= \frac{2^s-1}{2^l-1}.$
  

\item \textbf{If $\gcd(r,s)= \ell\neq 1$ and $n$ and $ \ell$ are  odd,} then
$$ N=\dfrac{(2^s-1)^2}{ \lcm\left(\dfrac{2^{l}-1}{ \gcd(n, 2^{l}-1) }, \dfrac{2^{l}-1}{1}  \right)}.$$
\begin{enumerate}
    \item If $ \gcd(n, 2^{\ell}-1)=1  $ then the number of $( n,\ell, \Fq^{\sigma})$-equivalence classes is $ \displaystyle{N =\frac{(2^s-1)^2}{2^{\ell}-1}.} $
By Theorem \ref{RankClasses}, we have 
     $$d_0=\gcd( n, 2^{\ell}-1)= 1,\ \ d_{\ell}= \gcd( n-\ell, 2^{\ell}-1)= 1 \text{ and} \ \  d=\gcd\left( \frac{ 2^{\ell}-1}{d_0}, \frac{2^{\ell}-1}{d_{\ell}}\right)=
     2^{\ell}-1.$$
 It follows that each  skew $(\ell,\sigma)$-trinomial  code associated with a skew  polynomial $f$ of the form  $f =x^{n}- a_{\ell} x^{\ell}-a_0$ is rank equivalent to a skew $(\ell,\sigma)$-trinomial  code  associated with one of the following  skew polynomials: 
  $$f_i(x)= x^n- x^{\ell}- \xi^{\ell n'},$$
  for $\ell \in \{ 0,\ldots, d n'^2-1\},$ with $ n':= \frac{2^s-1}{2^{\ell}-1}.$
  \item If $ 2^{\ell}-1  $ is a multiple of $n$, then the number of $( n,\ell, \Fq^{\sigma})-$equivalence classes is $ \displaystyle{N =\frac{(2^s-1)^2}{2^{\ell}-1}.} $
By Theorem \ref{RankClasses}, we have 
     $$d_0=\gcd( n, 2^{\ell}-1)= 2^{\ell}-1,\ \ d_{\ell}= \gcd( n-\ell, 2^{\ell}-1)= 1 \text{ and} \ \  d=\gcd\left( \frac{ 2^{\ell}-1}{d_0}, \frac{2^{\ell}-1}{d_{\ell}}\right)=1.$$
     
 It follows that each  skew $(\ell,\sigma)$-trinomial  code associated with a skew  polynomial $f$ of the form  $f =x^{n}- a_{\ell} x^{\ell}-a_0$ is rank equivalent to a skew $(\ell,\sigma)$-trinomial  code  associated with one of the following  skew polynomials: 
  $$f_i(x)= x^n- x^{\ell}- \xi^{i+\ell n'},$$
  for $i \in \{ 0,\ldots,d_0-1\} ;\quad \ell \in \{ 0,\ldots,n'^2-1\},$ with $ n':= \frac{2^s-1}{2^{\ell}-1}.$
\end{enumerate}

 \end{enumerate}
\end{example}




\section*{Conclusion}
In this paper, we investigated bounds on the (rank and Hamming) minimum distance of and equivalences of classes of skew polycyclic codes over a finite field $\mathbb{F}_q$, associated with a skew polynomial $f(x) \in \mathbb{F}_q[x;\sigma]$. First, we proved the Roos-like bound for both the Hamming and the rank metric for this class of codes. Our bound generalized many known bounds on various types of cyclic codes, both in the skew or commutative setting. Moreover, we analyzed constructions for such code with a designed distance, based on the concept of $\mu$-closed sets as an analog of cyclotomic cosets in the commutative case. 

Next, we focused on the equivalence of ambient spaces for skew polycyclic codes with respect to both metrics. We provided general conditions under which two families of skew polycyclic codes are equivalent. This was used to determine the exact number of equivalence classes. In the particular case of skew trinomial codes, we showed which code families are equivalent to the ones defined by the polynomial $x^n-x^\ell -1$, and more generally, how any equivalence class of a given skew trinomial can be determined. Finally, we gave some examples illustrating applications of the theoretical results on equivalence developed in this paper.

In future work, we will consider more general isometries to reduce the number of equivalence classes and use them for a more refined classification of polycyclic codes. Moreover, we plan to introduce the concept of ``skew cyclotomic cosets'' to facilitate the factorization of $x^n-1$ and $x^n-a$ and trinomial polynomials in $\mathbb{F}_q[x;\sigma]$, which may give an efficient way to classify skew polycyclic and skew constacyclic  codes of length $n$ over a finite field $\Fq$. 



\printbibliography


\section{Appendix}
In this section, for clarity we provide some proofs regarding rank metric equivalence that were previously omitted due to their similarity to the Hamming case. 

\subsection{ Proof of Proposition \ref{ProofEquiv} }
\begin{proof}
The proof is the same in both cases; therefore, for simplicity, we denote $\cong_{(n,\sigma)}$
  without specifying the type of equivalence used.
\begin{enumerate}
\item  $f_1(x) \cong_{n,\sigma} f_1(x)  ,$ since 
%$\varphi:=\id ,$ 
the identity  automorphism of  $\mathbb{F}_q[x;\sigma] /\langle 
f_1(x)\rangle$  is  an $\Fq$-morphism isometry.
\item  If $ f_1(x) \cong_{n,\sigma} f_2(x) $ 
then there exists  $ \varphi_{\alpha}$ an $\Fq$-morphism isometry
 $$\varphi: \mathbb{F}_q[x,\sigma] /\langle f_2(x) \rangle \longrightarrow \mathbb{F}_q[x,\sigma] /\langle f_1(x) \rangle.$$
 
As $ \varphi$ preserves the Hamming (resp. rank) weight, then for all $f\in \mathbb{F}_q[x,\sigma] /\langle f_2(x) \rangle, $   $$ w(f)=w(\varphi(\varphi^{-1}(f)))=w(\varphi^{-1}(f)).$$
Hence, 
$$
\varphi^{-1}: \mathbb{F}_q[x,\sigma] /\langle f_1(x) \rangle \longrightarrow \mathbb{F}_q[x,\sigma] /\langle f_2(x) \rangle.
$$
is also an  $\Fq$-morphism isometry. Hence, $ f_2(x) \cong_{n,\sigma} f_1(x).$

\item Suppose that $ f_1(x) \cong_{n,\sigma} f_2(x) $  and $ f_2(x) \cong_{n,\sigma} f_3(x),$  then there exist two $\Fq$-morphism isometries
$$
\begin{aligned}
& \varphi_1: \mathbb{F}_q[x,\sigma] /\left\langle f_2(x) \right\rangle \longrightarrow \mathbb{F}_q[x,\sigma] /\langle 
f_1(x) \rangle \\
& \varphi_2: \mathbb{F}_q[x,\sigma] /\langle f_3(x)  \rangle \longrightarrow \mathbb{F}_q[x,\sigma] /\langle f_2(x)  \rangle .
\end{aligned}
$$
Therefore
$$
\varphi_{1}  \circ \varphi_{2} : \mathbb{F}_q[x,\sigma] /\langle f_3(x)  \longrightarrow 
\mathbb{F}_q[x,\sigma] /\langle f_1(x)  \rangle
$$
is an $\Fq$-morphism isometry. Hence, $ f_1(x) \cong_{n,\sigma} f_3(x). $
\end{enumerate}
 \end{proof}
 
 \subsection{Proof of Corollary \ref{Th.1FqSigma}}
 \begin{proof}\;
\begin{itemize}
    \item[(1) $\Rightarrow$ (2):] Immediately follows from the Definition \ref{DefFqSigma}.
 \item[(2) $\Rightarrow$ (3):] Suppose that we have an $\Fq$-morphism  isometry $\varphi_{\alpha}$ given by 
$\varphi_{\alpha}(x) = \alpha x$ for some $ \alpha \in (\mathbb{F}_q^{\sigma})^*$. It follows that

$$ 
\varphi_{\alpha}(x^i)= (\alpha x)^i= N_i^{\sigma}(\alpha) x^{ i} =\alpha^i x^i,$$
since  $\alpha \in (\mathbb{F}_q^{\sigma})^* $. Therefore, $ N_i^{\sigma}(\alpha)=\alpha^i\ \forall  \ i=0,1, \ldots, n-1.$

As $\varphi_{\alpha}$ is an $\Fq$-morphism isometry and  $ \varphi_{\alpha}(x^n-b_{\ell}x^{\ell}-b_0)=0 \ \text{mod}_r (x^n- a_{\ell} x^{\ell}-a_0) ,$ we have
\begin{equation}
\varphi( x^n)=   b_{\ell}  \alpha^{\ell} x^{\ell} +b_0.
 \end{equation}

On the other hand,
\begin{equation}
\varphi( x^n)=  N_n^{\sigma}(\alpha) x^n= \alpha^n ( a_{\ell}x^{\ell}+a_0) = \alpha^n  a_{\ell} x^{\ell}+ \alpha^n a_0.
\end{equation}


Comparing term by term, we obtain  that $ a_0 \alpha^n  = b_0   $ and $  a_{\ell} \alpha^n =  \alpha^{\ell} b_{\ell}. $ As $ \alpha \in (\mathbb{F}_q^{\sigma})^*$, we have 
$$ a_0 \alpha^n= a_0 N_n^{\sigma}(\alpha)   = b_0   \quad  \text{and} \quad   a_{\ell} \alpha^{n-\ell} = a_{\ell}N_{n-\ell}^{\sigma}(\alpha) =  b_{\ell}, $$
which means that $\alpha$ is a common right root of the polynomials $a_0 x^n - b_0$ and $a_{\ell} x^{n-\ell} - b_{\ell}$.

    \item[(3) $\Rightarrow$ (4)\;] and (4) $\Rightarrow$ (5) are immediate.
 \item[(5) $\Rightarrow$ (6):]
    Let $ \alpha \in (\mathbb{F}_q^{\sigma})^*$ be a right root of the polynomial $\gcrd(x^{n} - b_0 a_0^{-1}, x^{n-\ell} - b_{\ell} a_{\ell}^{-1})$. Then $\alpha$ is a common root of the polynomials $a_0 x^n - b_0$ and $a_{\ell} x^{n-\ell} - b_{\ell}$. It follows that  $a_i N^{\sigma}_{n-i}(\alpha)=a_i \alpha^{n-i} = b_i,\ \text{for any} \ i\in \{ 0,\ell\}$, i.e.,
$$  (b_0, b_{\ell} )= (\alpha^n a_0, \alpha^{n-\ell} a_{\ell})=(\alpha^n, \alpha^{n-\ell})\star(a_0, a_{\ell} ).$$ 

\item[(6) $\Rightarrow$ (7):] Suppose that

$$  (b_0, b_{\ell} )= \left( \alpha^n a_0, \alpha^{n-\ell} a_{\ell}\right)=\left(\alpha^n, \alpha^{n-\ell})\right)\star (a_0, a_{\ell}).$$ 
It follows that
$$  (a_0,a_{\ell})^{-1}\star (b_0, a_{\ell})=  \left( \alpha^n, \alpha^{n-\ell} \right)= \left( \zeta^{jn}, \zeta^{j(n-\ell)}\right)= 
\left( \zeta^{n}, \zeta^{(n-\ell)}\right)^j .$$
Then $ (a_0,a_{\ell})^{-1}\star (b_0, a_{\ell}) $ belongs to the cyclic subgroup  $H_{\ell,\Fq^{\sigma}}$ generated by $  
\left( \zeta^{n}, \zeta^{(n-\ell)}\right) $  as a subgroup of $ (\Fq^{\sigma})^{*}\times (\Fq^{\sigma})^{*}.$

\item[(7) $\Rightarrow$ (1):]
     Suppose that $ (a_0,a_{\ell})^{-1}\star (b_0, b_{\ell}) $ is an element of the cyclic subgroup  $H_{\ell,\sigma}$ generated by 
     $  (\zeta^n, \zeta^{n-\ell})  $  as a subgroup of $ (\Fq^{\sigma})^{*}\times (\Fq^{\sigma})^{*}.$ Then there exists an integer $ h$ such that 
     $$ (a_0,a_{\ell})^{-1}\star (b_0, b_{\ell})=   \left( \zeta^{nh}, \zeta^{(n-\ell)h}\right).$$
    For $ \alpha=  \zeta^{h}, $ we obtain  that $  a_i \alpha^{n-i}= b_i, $ for any $ i\in \{ 0,\ell\}.$
As in the proof of theorem \ref{Th.1}, we can prove that $ \varphi_{\alpha}$ given by
  \begin{equation}
        \begin{array}{cccc}
        \tilde{\varphi}_{\alpha} : & \mathbb{F}_q[x;\sigma]/\langle x^n - b_{\ell} x^{\ell} - b_0 \rangle,   & \longrightarrow & \mathbb{F}_q[x;\sigma] /\langle x^n - a_{\ell} x^{\ell} - a_0 \rangle, \\ 
        & f(x) & \longmapsto & f(\alpha x),
        \end{array}
    \end{equation} 
    is an $\Fq$-morphism isometry with respect to the rank metric, with particularity that $ \alpha \in (\mathbb{F}_q^{\sigma})^* $ and $ N_i^{\sigma}(\alpha)= \alpha^i.$ 
\end{itemize}
By the equivalence between (1) and (7), we deduce that the number of $(n,\ell,\Fq^{\sigma})$-equivalence classes on $\Fq^*\times \Fq^*$ corresponds to the order of the group $ \left( \Fq^*\times \Fq^* \right)/H_{\ell,\Fq^{\sigma}},$ which equals 
$$ N_{(n,\ell,\Fq^{\sigma})}=\dfrac{ (q-1)^2}{ \lcm(\frac{q_0-1}{ \gcd(n,q_0-1)}, \frac{q_0-1}{ \gcd(n-\ell,q_0-1)})} . $$
 \end{proof}

 

\subsection{Proof of Theorem  \ref{RankClasses}}
\begin{proof}\;
\begin{enumerate}
    \item  If $ d = 1 $, then the cyclic group $ H_{\ell,\sigma} $ generated by $ \left( \zeta^n, \zeta^{n-\ell}\right) $ is isomorphic to the group $ \langle  N_n^{\sigma}(\zeta) \rangle \times \langle N_{n-\ell}^{\sigma}(\sigma^{\ell}(\zeta))  \rangle $ and has order $ \frac{(q'-1)^2}{d_0 d_{\ell}} $. By Theorem \ref{Th.1}, the number of $ (n,\ell,\sigma, \mathbb{F}_{q'})$-equivalence classes is $ \left(\frac{q-1}{q'-1}\right)^2 d_0 d_{\ell}= n'^2 d_0 d_{\ell} $. 
    Therefore, we can partition $ \Fq^* \times \Fq^* $ as follows

$$
\Fq^* \times \Fq^* = \bigcup_{l =0}^{n'^2} \bigcup_{i=0}^{d} \bigcup_{j=0}^{d_{\ell}-1} ( \xi^{i+ln'}\xi^i, \xi^j ) H_{\ell, \mathbb{F}_{q'} }.
$$
Then any pair $(a_0,a_{\ell})$ is $( n,\ell,\sigma, \mathbb{F}_{q'})$-equivalent to one of the pairs   $ ( \xi^{i+ln' },  \xi^{j } ) , $ for  $i=0,1,\ldots, d_0-1$; \   $j=0,1,\ldots, d_{\ell}-1;  $ and $ l=0,\ldots, n'^2-1.$
 \item 
 If $d\neq 1, $ the number of  $(n,\ell,\sigma, \mathbb{F}_{q'})$-equivalence classes is $ n'^2 d d_0 d_{\ell} ,$ and so we partition $ \Fq^* \times \Fq^* $ as follows:
$$  \Fq^{*}\times \Fq^{*} = 
\bigcup_{h=0}^{d n'^2 -1}  \bigcup_{i=0}^{d_0-1} \bigcup_{j=0}^{d_{\ell}-1} (\xi^{i +hn'} ,\xi^{j}) H_{\ell,\mathbb{F}_{q'}}, $$
which implies the second statement, similarly to the first case.
\end{enumerate}
\end{proof}\;

\subsection{Proof of Proposition \ref{RankEquiv_Consta}}
\begin{proof}
\begin{enumerate}
\item  As $ (a_0, a_{\ell}) \sim_{(n,\ell,\sigma, \mathbb{F}_{q'})} (b_0, b_{\ell})$, by the second  statement of Theorem \ref{Th.1}, 
there exists $\alpha\in \Fq^*$ such that $$ a_0 N_n^{\sigma}(\alpha)  = b_0   \quad  \text{and} \quad   a_{\ell} N^{\sigma}_{n-\ell}(\sigma^{\ell}(\alpha))=  b_{\ell}. $$
This gives
$$ a_0^{-1}  b_0 = N_n^{\sigma}(\alpha)  \quad  \text{and} \quad   a_{\ell}^{-1}   b_{\ell}=  N^{\sigma}_{n-\ell}(\sigma^{\ell}(\alpha)) . $$
Which means that  
$$ a_0^{-1}  b_0 \in   \langle  N_n^{\sigma}(\xi) \rangle  \quad  \text{and} \quad   a_{\ell}^{-1}   b_{\ell} \in  
\langle   N^{\sigma}_{n-\ell}(\sigma^{\ell}(\xi))  \rangle. $$

\item Let   $i\in \{0,\ell\}$ then the order of $ \langle   N^{\sigma}_{n-i}(\sigma^{i}(\xi))  \rangle =
\langle   \xi^{p^{ri}[n]_r}  \rangle $ is  $$d_i=\dfrac{q-1}{\gcd(p^{ri}[n-i]_r, q-1)}= \dfrac{q-1}{\gcd([n-i]_r, q-1)} , \text{ since $ \gcd(p^{ri},q-1)=1.$ }$$  It follows that  $(b_i a_i^{-1})^{d_i} =1 ,\ i\in \{0,\ell\}.$
\item by the second  statement of Theorem \ref{Th.1}, 
there exists $\alpha\in \Fq^*$ such that $$ a_0 N_n^{\sigma}(\alpha)  = b_0   \quad  \text{and} \quad   a_{\ell} N^{\sigma}_{n-\ell}(\sigma^{\ell}(\alpha))=  b_{\ell}. $$  As in the proof of \cite[Theorem 6]{Ouazzou2025}. For $i\in \{0,\ell\},$ we verify
 that the  map  $\varphi_i$ given by 
\begin{equation}
  	\begin{array}{cccc}
  	\varphi_{i}:& \mathbb{F}_q[x] /\langle x^{n-i}-b_0 \rangle  &\longrightarrow & \mathbb{F}_q[x] /\langle x^{n-i}-a_0 \rangle, \\
  	& f(x) & \longmapsto &  f(\sigma^i(\alpha) x).
  	\end{array}
  	\end{equation} 
    is an $\Fq$-morphism isometry with respect to the Hamming  metric.
    \item Follows from the third statement.
\end{enumerate}
\end{proof}

\subsection{More details in relation to Remark \ref{ConstacyclicFq_sigma}}
We prove here an alternative version of Proposition \ref{RankEquiv_Consta}, when we calculate the rank over the fixed subfield.
 \begin{proposition}\label{PropFq_sigma}
Let $ (a_0, a_{\ell}) $ and $ (b_0, b_{\ell}) $ be elements of $ \Fq^{*} \times \Fq^{*} $ such that $ (a_0, a_{\ell}) \sim_{(n, \ell,\Fq^{\sigma)}} (b_0, b_{\ell}) $. Then:
\begin{enumerate}
    \item There exists $\alpha\in (\Fq^{\sigma})^{*} ,$ such that for each $ i \in \{0, \ell\} $, $ a_i^{-1} b_i \in \langle \alpha^{n-i} \rangle.$
    \item For each $ i \in \{0, \ell\} $, $ (a_i^{-1} b_i)^{d_i} = 1 $, where $ d_i := \dfrac{q_0-1}{\gcd(n-i, q_0-1)} $.
    \item For each $ i \in \{0, \ell\} $, the map 
    \begin{equation*}
  	\begin{array}{cccc}
  	\varphi_{i}:& \mathbb{F}_q[x;\sigma] /\langle x^{n-i}-b_i\rangle  &\longrightarrow & \mathbb{F}_q[x;\sigma] /\langle
    x^{n-i}-b_i\rangle, \\ 
  	& f(x)& \longmapsto &  f(\alpha x)  
  	\end{array}
  	\end{equation*} 
is an $\mathbb{F}_q$-morphism isometry  with respect to the rank  distance.
    
    \item For each $ i \in \{0, \ell\} $ the class of skew  $(a_i,\sigma)$-constacyclic codes of length $n-i$ is { rank equivalent} to the class of skew   $(b_i,\sigma)$-constacyclic codes of length $n-i$  over $\Fq.$ 
\end{enumerate}
\end{proposition}
\begin{proof}
\begin{enumerate}
\item  As $ (a_0, a_{\ell}) \sim_{(n, \ell,\Fq^{\sigma)}} (b_0, b_{\ell})$, by the third statement of Theorem \ref{Th.1FqSigma}, there
exists  $\alpha\in \Fq^{\sigma}$ which is a common right root of $\displaystyle a_i x^{n-i } - b_i \in \Fq[x;\sigma], i \in \{0,\ell\}. $ It follows that $$ a_0 \alpha^n= a_0 N_n^{\sigma}(\alpha)   = b_0   \quad  \text{and} \quad   a_{\ell} \alpha^{n-\ell} = a_{\ell}N_{n-\ell}^{\sigma}(\alpha) =  b_{\ell}. $$ 
This gives
$$ a_0^{-1}  b_0 = \alpha^n  \quad  \text{and} \quad   a_{\ell}^{-1}   b_{\ell}=  \alpha^{n-\ell} . $$
Which means that  
$$ a_0^{-1}  b_0 \in   \langle \alpha^{n} \rangle  \quad  \text{and} \quad   a_{\ell}^{-1}   b_{\ell} \in   \langle \alpha^{n-\ell} \rangle. $$

\item  Let $\zeta $ be a primitive element of $ \Fq^{\sigma}$ then $\alpha= \zeta^i\in \Fq^{\sigma}$ for some $ 0\leq i\leq q_0-1.$  
Then for each  $i\in \{0,\ell\}$ the order of $ \langle \alpha^{n-i} \rangle $ is  $d_i=\dfrac{q_0-1}{\gcd(n-i, q_0-1)}.$ It follows that  $(b_i a_i^{-1})^{d_i} =1 ,\ i\in \{0,\ell\}.$
\item By Theorem \ref{Th.1FqSigma}, there is  a common root $\alpha$  of $a_0  x^n-b_0$ and $a_1 x^{n-1}-b_1$. Then for 
 any  $i\in \{0,1\}$,  we have $ a_0  \alpha^{n-i}=b_0.$ As in the proof of \cite[Theorem 6]{Ouazzou2025}, we verify
 that the  map  $\varphi_i$ given by 
\begin{equation}
  	\begin{array}{cccc}
  	\varphi_{i}:& \mathbb{F}_q[x] /\langle x^{n-i}-b_0 \rangle  &\longrightarrow & \mathbb{F}_q[x] /\langle x^{n-i}-a_0 \rangle, \\
  	& f(x) & \longmapsto &  f(\alpha x).
  	\end{array}
  	\end{equation} 
    is an $\Fq$-morphism isometry with respect to the rank metric.
    \item Follows from the third statement.
\end{enumerate}
\end{proof}

\end{document} 


\textcolor{blue}{Yes I got your point.   The following are some first remarks }
\begin{enumerate}

\item \textbf{The case of skew constacylic codes.}
\begin{enumerate}

    \item  To be sure that $ \varphi$ sends a skew polycyclic code  $C=\langle g(x)\rangle$ associated with $f_2(x)$ to another polycycic code  and $ C'=\langle \varphi(g(x))\rangle$  associated with $f_1,$  We need that $\varphi$ possesses  "$ \Fq$-morphism property", I mean $ \varphi$ satisfy    $ \varphi(a(x)b(x))= \varphi(a(x))\varphi(b(x))$
    for any $a(x), b(x) \in \Fq[x;\sigma] /\langle f_1(x) \rangle .$   But the verification of the "$ \Fq$- morphism property" is  the worst part, even for  the case where $\varphi(x)=\alpha x. $ \textcolor{purple}{Yes, of course. In the end, this is exactly what our results are all about. I would just like to have a good justification why we can even restrict to just $\alpha x$ and in the Hamming metric you can easily do this with the distance. It would be nice to do the same for the rank metric. (Maybe I was not clear enough -- I did not mean just to look at $x\mapsto x^{n-1}$ etc., but all of these with possible scalar coefficients.) Anyway, have a look at my comment in the chat -- I guess there are known results about the structure of ring morphisms that we implicitly use here, right?}
    \textcolor{blue}{yes if we consider morphisms of petit algebras we will obtain many results. But here we need to consider the class of $\Fq^{\sigma}$-linear rank metric codes. I am still a bit confusing about the structure here is ait additive over $\Fq$ or not..., Anyway I have some results I will add them as soon possible.  }
    
    Let for example $ f_1=x^n-a_0$ and $f_2=x^n-b_0$ are skew polynomials, and
    $$\varphi: \  \mathbb{F}_q[x;\sigma]/\langle f_2(x)\rangle \longrightarrow  \mathbb{F}_q[x;\sigma]/\langle f_1(x)\rangle; \ \ x \longrightarrow x^{n-1}. $$
    We can verify that $ \varphi$ is  $ \Fq$-linear, and if we assume  $ \varphi$ is an $\Fq$-morphism. We observe that 
    $$ \varphi(x^i)=  \varphi(x)^i= x^{i(n-1)}= x^{in} x^{-i}= N_{i}^{\sigma^n}(a_0)  x^{-i}  \mod (x^n-a_0),$$
    But 
    $$ \varphi(x^n)=  \varphi(x)^n= x^{n(n-1)}= N_{n-1}^{\sigma^n}(a_0)  \mod (x^n-a_0),$$
    on the other hand  as $ x^n-b_0= b_0 \mod (x^n-b_0)$
    $$  \varphi(x^n)= \varphi(b_0 )= b_0 \mod( x^n-a_0)  ,$$ 
    Next we combine the two value and  the equivalence gives us this equality
    $$   b_0= N_{n-1}^{\sigma^n}(a_0)   $$

    
    
    \item To preserve the length and dimension $ \varphi$ should be and $\Fq$- vector space isomorphism.
    As above let $ f_1=x^n-a_0$ and $f_2=x^n-b_0$ are skew polynomials, and
    $$\bar \varphi: \  \mathbb{F}_q[x;\sigma]  \longrightarrow  \mathbb{F}_q[x;\sigma]/\langle f_1(x)\rangle; \ \ x \longrightarrow x^{n-1}. $$
We can verify that 
$ \bar \varphi( x^n-b_0)= \varphi(x^n)-b_0= N_{n-1}^{\sigma^n}(a_0) -b_0=0.  $
So we can prove that $\varphi$  is an $\Fq$-isomorphism. Hence  it preserves the rank metric in the case of constacyclic codes.

\item Let check the $\Fq$-morphism property, i.e we need to prove that
$$ \varphi(a_ix^i b_jx^j)= \varphi(a_ix^i) \varphi(b_j x^j)$$

Fist we have 
$$a_i x^i  b_j x^j = a_i  \sigma^i(b_j)  x^{i+j}.$$
Then 
$$ \varphi( a_i x^i  b_j x^j ) = a_i  \sigma^i(b_j)  N_{i+j}^{\theta}(a_0)  x^{-(i+j)} , \quad \text{ with } \theta=\sigma^n. $$

In the other hand 
$$ \varphi( a_i x^i)  \varphi( b_j x^j ) = a_i \sigma^{-i}(b_j)  N_i^{\theta}(a_0) \sigma^{-i}(N_j^{\theta}(a_0))  x^{-(i+j)} , \quad \text{ with } \theta=\sigma^n. $$

To prove obtain the $\Fq$-morphism property we need that 
$$   \sigma^i(b_j N_j^{\theta}(a_0)) = \sigma^{-i}(b_j N_j^{\theta}(a_0)),$$

which is not generally true unless   $ \sigma^{2i}(a) =  a$, for all $a\in \Fq$ and $ i =0,\ldots, n-1$. Therefore, $ \varphi $ is not a ring $\Fq$-morphism in general.
\item In conclusion if $ x^n-a_0  $ and $ x^n-b_0$ $ \varphi(x)=x^{n-1}$ is not an $\Fq$-morphism isometry. It preserves the rank weight but dosen't preserve the (((ideal structure))) of codes, I mean that

If $g(x) $ right divides $ f_2(x)$ we are not able to say that $\varphi(g(x))= g(x^{n-1}) $ right divides $ f_1(x). $ Even though, $ \deg(g(x)) = \deg(\varphi(g(x)),$ since $ \varphi$ is an $ \Fq$-vector space isomorphism.  

    \item  if 
    $$ \varphi(x)= x^{n-1}+x^{n-2} ,$$
   
    $$ \varphi ( x^n) = \varphi(x)^n= ( x^{n(n-1)}+x^{n(n-2)} )=  N^{\sigma^n}_{n-1}(a_0)+ N^{\sigma^n}_{n-2}(a_0) \mod (x^n-a_0) $$
    and 
    $$ \varphi (x^n)= \varphi(b_0)=b_0  $$
   It follows that  the equivalence gives us this equality
    $$ b_0= N^{\sigma^n}_{n-1}(a_0)+ N^{\sigma^n}_{n-2}(a_0)$$
    As above we can verify that 
    $$\bar \varphi: \  \mathbb{F}_q[x;\sigma]  \longrightarrow  \mathbb{F}_q[x;\sigma]/\langle f_1(x)\rangle; \ \ x \longrightarrow x^{n-1}. $$
We can verify that 
$ \bar \varphi( x^n-b_0)= \varphi(x^n)-b_0=  N^{\sigma^n}_{n-1}(a_0)+ N^{\sigma^n}_{n-2}(a_0) -b_0=0.  $
So we can prove that $\varphi$  is an $\Fq$-isomorphism. Hence it preserves the rank metric.

But also is not an $ \Fq$-morphism !!!!!

\item More general if $ \varphi(x)= x^{n-1}+x^{n-2}+ \ldots+ x+1 $ then 
$$ \varphi ( x^n) = \varphi(x)^n= ( x^{n(n-1)}+x^{n(n-2)}+ \ldots+ x^n+1 )=  N^{\sigma^n}_{n-1}(a_0)+ N^{\sigma^n}_{n-2}(a_0) + \ldots+N_1^{\sigma^n}(a_0)+1 \mod (x^n-a_0) $$
    and 
    $$ \varphi (x^n)= \varphi(b_0)=b_0  $$
    It follows that  the equivalence gives us this equality
    $$ b_0=   N^{\sigma^n}_{n-1}(a_0)+ N^{\sigma^n}_{n-2}(a_0) + \ldots+N_1^{\sigma^n}(a_0)+1 $$





The same problem here $\varphi$ is not an $ \Fq$-morphism !!!!!
    
    
    
    \end{enumerate}    
    \item  \textbf{The case of skew trinomial codes}
    \begin{enumerate}
        \item Let $ f_1= x^n-b_{\ell} x^{\ell}-b_0$   and $ f_2= x^n-a_{\ell} x^{\ell}-a_0,$  and 
            $$\varphi: \  \mathbb{F}_q[x;\sigma]/\langle f_2(x)\rangle \longrightarrow  \mathbb{F}_q[x;\sigma]/\langle f_1(x)\rangle; \ \ x \longrightarrow x^{n-1}. $$
  
     Also here the problem is that $ \varphi$  is not an $ \Fq$-morphism !!!!!   
    \end{enumerate}
   


        
        
    \end{enumerate}
\textcolor{blue}{End of the comment.}


Note that any $\mathbb{F}_q$-vector space isomorphism of the form $\varphi(f(x)) = f(\alpha x^{h})$, for some $\alpha \in \mathbb{F}_q^*$ and $h \in \{1, \ldots, n - 1$\}, preserves the code length $n$, the dimension $k$, and the Hamming distance ( \textcolor{purple}{Not that straight-forwardly} \textcolor{blue}{
I assumed that there exist  $ \varphi$ an $\mathbb{F}_q$-vector space isomorphism and   of the form $\varphi(f(x)) = f(\alpha x^{h}).$ Then since it is isomorphism, hence the rank distance is preserved. I hope that I didn't wrong here? } ).  Similarly, any $\mathbb{F}_q$-vector space isomorphism $\varphi$ preserves the code length $n$, the dimension $k$, and the rank distance. However, such maps do not necessarily preserve the algebraic structure of the codes. To ensure this, we must verify that
$$
\varphi(f(x) * g(x)) = \varphi(f(x)) * \varphi(g(x)).
$$
Let for example $ f_1=x^n-a_0$ and $f_2=x^n-b_0,$ and
    $$\varphi: \  \mathbb{F}_q[x;\sigma]/\langle f_2(x)\rangle \longrightarrow  \mathbb{F}_q[x;\sigma]/\langle f_1(x)\rangle; \ \ x \longrightarrow  \alpha x^{n-1}, $$ for some $ \alpha \in \Fq^*.$ We can verify easily that $ \varphi$ is  $ \Fq$-vector space homomorphism, and with a simple computation we obtain that 
    $$ \varphi(x^i)=   N_{i(n-1)}^{\sigma}(\alpha)  x^{i(n-1)}=   N_{i(n-1)}^{\sigma}(\alpha)  x^{in} x^{-i}=  N_{i(n-1)}^{\sigma}(\alpha)  N_{i}^{\sigma^n}(a_0)  x^{-i}  \mod (x^n-a_0),$$
    and so
    $$ \varphi(x^n)= N_{n(n-1)}^{\sigma}(\alpha) x^{n(n-1)}=  N_{n(n-1)}^{\sigma}(\alpha) N_{n-1}^{\sigma^n}(a_0)  \mod (x^n-a_0),$$
    on the other hand  as $ x^n-b_0= b_0 \mod (x^n-b_0),$ then 
    $$  \varphi(x^n)= \varphi(b_0 )= b_0 \mod( x^n-a_0)  ,$$ 
    Next we combine the two values of $ \varphi(x^n) $ and  the equivalence gives us this equality
    $$   b_0=  N_{n(n-1)}^{\sigma}(\alpha) N_{n-1}^{\sigma^n}(a_0)   $$
To preserve the length $n$ and dimension $k$, $ \varphi$ should be and $\Fq$- vector space isomorphism. Then for $   b_0=  N_{n(n-1)}^{\sigma}(\alpha) N_{n-1}^{\sigma^n}(a_0)$ we consider 
    $$\bar \varphi: \  \mathbb{F}_q[x;\sigma]  \longrightarrow  \mathbb{F}_q[x;\sigma]/\langle f_1(x)\rangle; \ \ x \longrightarrow x^{n-1}. $$
We can verify that 
$ \bar \varphi( x^n-b_0)= \varphi(x^n)-b_0=   N_{n(n-1)}^{\sigma}(\alpha) N_{n-1}^{\sigma^n}(a_0) -b_0=0.  $
and so we can prove $\varphi$  is an $\Fq$-isomorphism. Hence  it preserves both  rank and the Hamming metric.

\noindent Let check the $\Fq$-morphism property, i.e we need to prove that $$ \varphi(a_ix^i b_jx^j)= \varphi(a_ix^i) \varphi(b_j x^j).$$ Fist we have 
$$a_i x^i  b_j x^j = a_i  \sigma^i(b_j)  x^{i+j}.$$
Then 
$$ \varphi( a_i x^i  b_j x^j ) = a_i  \sigma^i(b_j)  N_{(i+j)(n-1)}^{\sigma}(\alpha)  N_{i+j}^{\sigma^n}(a_0)  x^{-(i+j)} , . $$
In the other hand, 
$$ \varphi( a_i x^i)  \varphi( b_j x^j ) = a_i  N_{i(n-1)}^{\sigma}(\alpha) N_i^{\sigma^n}(a_0)  \sigma^{-i}(b_j)  \sigma^{-i}(N_{j(n-1)}^{\sigma}(\alpha))   \sigma^{-i}(N_j^{\theta}(a_0))  x^{-(i+j)}. $$

If $b=1$ we obtain an nice relation with skew Reed-Solomon codes.

\begin{proposition}
     Let $g \in \Fq[x,\sigma]$ such that $ g(x) ~|_r~ f(x),$ $ C= \Fq[x,\sigma] g(x)$ a skew polycyclic code generated by $g(x)$ and $\widehat{C}= \mathbb{F}_{q_0^e}[x,\sigma] g(x) $ its extension code over $ \mathbb{F}_{q_0^e}.$ Let $\beta= \theta(\alpha)\alpha^{-1}$ with $\alpha$ be a normal element of $ \mathbb{F}_{q_0^e}$. 
     Suppose that there exist integers $ a,  \delta, r, k_0, \dots , k_r$ such that $\gcd( b,e)=1,\  k_j<k_{j+1}\ \text { for } 0 \leq j \leq r-1, \ k_r-k_0 \leq \delta+r-2, $ and  
    $$ (x-{\theta^{a+i+k_j}(\beta)}) \mid_r g(x)  ~:~ i=0,\ldots,{\delta-2}, j=0,\ldots, r-1,$$ 
    i.e., $$ \left\{ a+i+k_j   ~:~ i=0,\ldots,{\delta-2}, j=0,\ldots, r-1 \right\} \subseteq T_{\beta}(g). $$
    
     In particular $d_H(C) \geq \delta+r$ and $d_R(C) \geq \delta+r$.
\end{proposition}  




\subsection{ Proof of Corollary \ref{Caract_(1,1)} }
\begin{proof}\;
\begin{itemize}
    \item[(1) $\Rightarrow$ (2):] Follows from the equivalence of  (1) and (6) of Corollary \ref{Th.1FqSigma}.
 \item[(2) $\Rightarrow$ (3):] 
 If $(a_0, a_{\ell})\star\left(\alpha^n, \alpha^{n-\ell}\right) = (1, 1)$, then 
$$ a_0 \alpha^{n}=1 \quad \quad 
\text{and} \quad  
 \quad a_\ell \alpha^{n-\ell}=1.
$$
 It follows that
 $$  a_0  \alpha^{n}=1 \quad \text{ and } \quad  a_\ell= \alpha^{\ell-n} a_0  \alpha^{n}=\alpha^{\ell} a_0.
  $$

 
  i.e., $\beta:=\alpha^{-1}$ is an $n$-th root of $a_0$ and $  a_{\ell} \beta^{\ell}= a_0  $.
 

 \item[(3) $\Rightarrow$ (1):] Suppose that $\alpha $ is an $n$-th root of $a_0$ such that $ a_0 =  \alpha^{\ell} a_\ell.$

 $$ \alpha^n  = a_0\ \ \text{ and } \ \   
   a_{\ell}=  \alpha^{-\ell}a_0= \alpha^{n-\ell}. $$ 
 It follows that  
$$(a_0,a_{\ell} )\star \left(\alpha^{-n}, \alpha^{\ell-n} \right)= (1,1) .$$
 By the  statement (6) of Theorem \ref{Th.1FqSigma}, we conclude that  $ (a_0, a_{\ell}) \sim_{( n,\ell, \mathbb{F}_{q'})}(1,1) ,$ and the result holds.  
\end{itemize}
 \end{proof}
